%%% The main file. It contains definitions of basic parameters and includes all other parts.

%% Settings for single-side (simplex) printing
% Margins: left 40mm, right 25mm, top and bottom 25mm
% (but beware, LaTeX adds 1in implicitly)
% \documentclass[12pt,a4paper]{report}
% \setlength\textwidth{145mm}
% \setlength\textheight{247mm}
% \setlength\oddsidemargin{15mm}
% \setlength\evensidemargin{15mm}
% \setlength\topmargin{0mm}
% \setlength\headsep{0mm}
% \setlength\headheight{0mm}
% \openright makes the following text appear on a right-hand page
%\let\openright=\clearpage

%% Settings for two-sided (duplex) printing
\documentclass[12pt,a4paper,twoside,openright]{report}
\setlength\textwidth{145mm}
\setlength\textheight{247mm}
\setlength\oddsidemargin{14.2mm}
\setlength\evensidemargin{0mm}
\setlength\topmargin{0mm}
\setlength\headsep{0mm}
\setlength\headheight{0mm}
\let\openright=\cleardoublepage

%% Generate PDF/A-2u
\usepackage[a-2u]{pdfx}

%% Character encoding: usually latin2, cp1250 or utf8:
\usepackage[utf8]{inputenc}

%% Prefer Latin Modern fonts
\usepackage{lmodern}

%% Further useful packages (included in most LaTeX distributions)
\usepackage{amsmath}        % extensions for typesetting of math
\usepackage{amsfonts}       % math fonts
\usepackage{amsthm}         % theorems, definitions, etc.
\usepackage{bbding}         % various symbols (squares, asterisks, scissors, ...)
\usepackage{bm}             % boldface symbols (\bm)
\usepackage{graphicx}       % embedding of pictures
\usepackage{fancyvrb}       % improved verbatim environment
\usepackage[numbers]{natbib}         % citation
\usepackage[nottoc]{tocbibind} % makes sure that bibliography and the lists
			    % of figures/tables are included in the table
			    % of contents
\usepackage{dcolumn}        % improved alignment of table columns
\usepackage{booktabs}       % improved horizontal lines in tables
\usepackage{paralist}       % improved enumerate and itemize
\usepackage{xcolor}         % typesetting in color
\usepackage{booktabs}       % better horizontal lines for tables
\usepackage{siunitx}        % thousands separators for large numbers in text
\usepackage[labelfont=bf]{caption} % use bold captions lables (Figure X.X:)
\usepackage{subcaption}     % subfigures
\usepackage{tikz}           % for drawing (checkmark etc)
\usepackage{bytefield}      % packet structure
\usepackage{float}       % better positioning of figures
\usepackage{newfloat}       % better positioning of figures
\usepackage{filecontents}   % for auto creating of .xmpdata file
\usepackage{listings}       % code listings
\usepackage{xparse}         % for serious parsing magic
\usepackage[strings]{underscore}     % underscores which do not screw up hyphenation
\usepackage[toc,nopostdot,indexonlyfirst]{glossaries}     % lists of terms and abbreviations
\usepackage{dirtree}        % directory listings with comments
\usepackage{mdframed}       % fancy frames

%%% src: https://tex.stackexchange.com/questions/57353/how-to-get-caption-above-listing-with-minted
\captionsetup[listing]{position=top}

%%% Graphics path (for some reason needed on windows)
\graphicspath{{img/}}

%%% Make listings part of the table of contents
\renewcommand{\lstlistoflistings}{\begingroup
\tocfile{List of Listings}{lol}
\endgroup}

%%% Basic information on the thesis

% Thesis title in English (exactly as in the formal assignment)
\def\ThesisTitle{QUIC Protocol Implementation for .NET}

% Author of the thesis
\def\ThesisAuthor{Bc. Radek Zikmund}

% Year when the thesis is submitted
\def\YearSubmitted{2020}

% Name of the department or institute, where the work was officially assigned
% (according to the Organizational Structure of MFF UK in English,
% or a full name of a department outside MFF)
\def\Department{Department of Distributed and Dependable Systems}

% Is it a department (katedra), or an institute (ústav)?
\def\DeptType{Department}

% Thesis supervisor: name, surname and titles
\def\Supervisor{Mgr. Pavel Ježek, Ph.D.}

% Supervisor's department (again according to Organizational structure of MFF)
\def\SupervisorsDepartment{Department of Distributed and Dependable Systems}

% Study programme and specialization
\def\StudyProgramme{Computer Science}
\def\StudyBranch{Software and Data Engineering}

% An optional dedication: you can thank whomever you wish (your supervisor,
% consultant, a person who lent the software, etc.)
\def\Dedication{%
Dedication.
}

% Abstract (recommended length around 80-200 words; this is not a copy of your thesis assignment!)
\edef\Abstract{%
TODO: do the abstract last%
}

% 3 to 5 keywords (recommended), each enclosed in curly braces
\def\Keywords{%
{QUIC} {.NET} {network protocol}
}

%% The hyperref package for clickable links in PDF and also for storing
%% metadata to PDF (including the table of contents).
%% Most settings are pre-set by the pdfx package.
\hypersetup{unicode}
\hypersetup{breaklinks=true}

%% make \autoref use section instead of subsection and subsubsection
\let\subsectionautorefname\sectionautorefname
\let\subsubsectionautorefname\sectionautorefname

%% generate thesis.xmpdata file
%% see https://tex.stackexchange.com/questions/534035/macros-commands-inside-a-filecontents-environment-does-not-expand
%% for the \filecontentsspecials magic

\def\filecontentsspecials#1#2#3{
  \global\let\ltxspecials\dospecials
  \gdef\dospecials{\ltxspecials
    \catcode`#1=0
    \catcode`#2=1
    \catcode`#3=2
    \global\let\dospecials\ltxspecials
  }
}
\filecontentsspecials|[]
\begin{filecontents}[overwrite]{\jobname.xmpdata}
\Author{|ThesisAuthor}
\Title{|ThesisTitle}
\Keywords{QUIC\sep .NET\sep network protocol}
\Subject{|Abstract}
\Publisher{Charles University}
\end{filecontents}

%% also generate abstract.txt so that it can be included from the abstract.tex file
\filecontentsspecials|[]
\begin{filecontents}[overwrite]{abstract.txt}
|Abstract
\end{filecontents}

% Definitions of macros (see description inside)
%%% This file contains definitions of various useful macros and environments %%%
%%% Please add more macros here instead of cluttering other files with them. %%%

%%% Minor tweaks of style

% These macros employ a little dirty trick to convince LaTeX to typeset
% chapter headings sanely, without lots of empty space above them.
% Feel free to ignore.
\makeatletter
\def\@makechapterhead#1{
  {\parindent \z@ \raggedright \normalfont
   \Huge\bfseries \thechapter. #1
   \par\nobreak
   \vskip 20\p@
}}
\def\@makeschapterhead#1{
  {\parindent \z@ \raggedright \normalfont
   \Huge\bfseries #1
   \par\nobreak
   \vskip 20\p@
}}
\makeatother

% This macro defines a chapter, which is not numbered, but is included
% in the table of contents.
\def\chapwithtoc#1{
\chapter*{#1}
\addcontentsline{toc}{chapter}{#1}
}

% Draw black "slugs" whenever a line overflows, so that we can spot it easily.
\overfullrule=1mm

%%% macro for replacing
\ExplSyntaxOn
\NewDocumentCommand{\replace}{mmm}
 {
  \marian_replace:nnn {#1} {#2} {#3}
 }

\tl_new:N \l_marian_input_text_tl

\cs_new_protected:Npn \marian_replace:nnn #1 #2 #3
 {
  \tl_set:Nn \l_marian_input_text_tl { #1 }
  \tl_replace_all:Nnn \l_marian_input_text_tl { #2 } { #3 }
  \tl_use:N \l_marian_input_text_tl
 }
\ExplSyntaxOff

%%% Functional foreach construct
%%% (https://stackoverflow.com/questions/2402354/split-comma-separated-parameters-in-latex)

\makeatletter

% #1 - Function to call on each comma-separated item in #3
% #2 - Parameter to pass to function in #1 as first parameter
% #3 - Comma-separated list of items to pass as second parameter to function #1
\def\foreach#1#2#3{%
  \@test@foreach{#1}{#2}#3,\@end@token%
}

% Internal helper function - Eats one input
\def\@swallow#1{}

% Internal helper function - Checks the next character after #1 and #2 and
% continues loop iteration if \@end@token is not found
\def\@test@foreach#1#2{%
  \@ifnextchar\@end@token%
    {\@swallow}%
    {\@foreach{#1}{#2}}%
}

% Internal helper function - Calls #1{#2}{#3} and recurses
% The magic of splitting the third parameter occurs in the pattern matching of the \def
\def\@foreach#1#2#3,#4\@end@token{%
  #1{#2}{#3}%
  \@test@foreach{#1}{#2}#4\@end@token%
}

%%% foreach usage:
% Example-function used in foreach, which takes two params and builds hrefs
%\def\makehref#1#2{\href{#1/#2}{#2}}

% Using foreach by passing #1=function, #2=constant parameter, #3=comma-separated list
%\foreach{\makehref}{http://stackoverflow.com}{2409851,2408268}

% Will in effect do
%\href{http://stackoverflow.com/2409851}{2409851}\href{http://stackoverflow.com/2408268}{2408268}

\makeatother



%%% Macros for definitions, theorems, claims, examples, ... (requires amsthm package)

\theoremstyle{plain}
\newtheorem{thm}{Theorem}
\newtheorem{lemma}[thm]{Lemma}
\newtheorem{claim}[thm]{Claim}

\theoremstyle{plain}
\newtheorem{defn}{Definition}

\theoremstyle{remark}
\newtheorem*{cor}{Corollary}
\newtheorem*{rem}{Remark}
\newtheorem*{example}{Example}

%%% An environment for proofs

\newenvironment{myproof}{
  \par\medskip\noindent
  \textit{Proof}.
}{
\newline
\rightline{$\qedsymbol$}
}

%%% The field of all real and natural numbers
\newcommand{\R}{\mathbb{R}}
\newcommand{\N}{\mathbb{N}}

%%% Useful operators for statistics and probability
\DeclareMathOperator{\pr}{\textsf{P}}
\DeclareMathOperator{\E}{\textsf{E}\,}
\DeclareMathOperator{\var}{\textrm{var}}
\DeclareMathOperator{\sd}{\textrm{sd}}

%%% Transposition of a vector/matrix
\newcommand{\T}[1]{#1^\top}

%%% Various math goodies
\newcommand{\goto}{\rightarrow}
\newcommand{\gotop}{\stackrel{P}{\longrightarrow}}
\newcommand{\maon}[1]{o(n^{#1})}
\newcommand{\abs}[1]{\left|{#1}\right|}
\newcommand{\dint}{\int_0^\tau\!\!\int_0^\tau}
\newcommand{\isqr}[1]{\frac{1}{\sqrt{#1}}}

%%% Various table goodies
\newcommand{\pulrad}[1]{\raisebox{1.5ex}[0pt]{#1}}
\newcommand{\mc}[1]{\multicolumn{1}{c}{#1}}

% Outputs red TODOs in the document. Requires \usepackage{color}.
%
% Usage: \todo{Document the TODO command.}
%
% Comment out second line to disable.
\newcommand{\todo}[1]{}
\renewcommand{\todo}[1]{{\color{red} TODO: {#1}}}

% the symbol of .NET
\newcommand{\dotnet}{.NET\@}

% prettier C# typesetting
\newcommand{\csharp}{%
  {\settoheight{\dimen0}{C}C\kern-.05em \resizebox{!}{\dimen0}{\raisebox{\depth}{\#}}}}

% special items for itemize and enumerate environments which have a label
\newcommand\litem[2][\itshape]{\item{#1 #2}:}

% special item for description lists to make the description go on the next line
\newcommand\ditem[1]{\pagebreak[3] \item[{\parbox[t]{0.98\textwidth}{#1}}] \leavevmode \\ }
\newcommand\ditemWithComment[2]{\pagebreak[3] \item[#1] \hfill \textit{(#2)} \leavevmode \\}

% macro for typesetting library names
\newcommand\libname[1]{\texttt{\itshape#1}}

% macros for library names used throughout the thesis
\newcommand{\libmsquic}{\libname{MsQuic}}
\newcommand{\libcurl}{\libname{libcurl}}
\newcommand{\libopenssl}{\libname{OpenSSL}}
\newcommand{\libschannel}{\libname{Schannel}}
\newcommand{\xUnit}{\libname{xUnit}}

% helper for composing captions from short and long entries
\newcommand{\myCaption}[2]{% short, long
  \def\temp{#1}\ifx\temp\empty% Use different short caption if provided
    \caption{#2}%
  \else%
    \caption[#1]{#2}%
  \fi%
}
\newcommand{\myCaptionInline}[2]{% short, long
  \def\temp{#1}\ifx\temp\empty% Use different short caption if provided
    #2%
  \else%
    [#1]#2%
  \fi%
}
\newcommand{\myCaptionOf}[3]{% env, short, long
  \def\temp{#2}\ifx\temp\empty% Use different short caption if provided
    \captionof{#1}{#3}%
  \else%
    \captionof{#1}[#2]{#3}%
  \fi%
}

% environment setting common attributes of figures
\newcommand{\figureArgs}{} % dummy macro to be redefined, used to make caption go under the figure
\newenvironment{myFigure}[3][]{% label, caption
  \renewcommand{\figureArgs}{\myCaption{#1}{#3}\label{#2}}%
  \figure[H]{}
  \centering
}{ %
  \figureArgs{}
  \endfigure{}
}

\newenvironment{mySubfigure}[3]{% [width, label, caption]
  \renewcommand{\figureArgs}{\caption{#3}\label{#2}}%
  \subfigure[H]{#1}
  \centering
}{ %
  \figureArgs{}
  \endsubfigure{}
}

% environment setting common attributes of tables
\newenvironment{myTable}[5][\normalsize]{% [fontsize, label, caption, columnspec, header]
  \table[H]
  #1
  \caption{#3}\label{#2}% for tables, the caption goes to the top
  \centering
  \tabular{#4}
  \toprule
  #5 \\
  \midrule
}{ %
  \bottomrule
  \endtabular{}
  \endtable{}
}


\definecolor{colorunimportant}{gray}{0.5}

% use VS like colors/
\definecolor{codecolorclass}{rgb}{0.12, 0.53, 0.62}
\definecolor{codecolorenum}{rgb}{0.40, 0.63, 0.23}
\definecolor{codecolorinterface}{rgb}{0.40, 0.63, 0.23}
\definecolor{codecolorkeyword}{rgb}{0, 0, 0.5}
\definecolor{codecolormethod}{rgb}{0.45, 0.32, 0.12}
\definecolor{codecolorkeywordcontrolflow}{rgb}{0.55, 0.03, 0.79}
\definecolor{codecolornamespace}{gray}{0}
\definecolor{codecolorcomment}{rgb}{0, 0.5, 0}
\definecolor{codecolorstring}{rgb}{0.64, 0.08, 0.08}
\definecolor{codecolorxmltag}{rgb}{0.64, 0.08, 0.08}
\definecolor{codecolorxmlstring}{rgb}{0, 0, 1}
\definecolor{codecolorxmlattribute}{rgb}{1, 0, 0}

% allow adding known classes to the listings for better syntax highlight
% https://tex.stackexchange.com/questions/172048/how-can-i-define-a-custom-class-of-keywords-in-listings
\makeatletter
\lst@InstallKeywords k{types}{typestyle}\slshape{typestyle}{}ld
\lst@InstallKeywords k{enums}{enumstyle}\slshape{enumstyle}{}ld
\lst@InstallKeywords k{attributes}{attributestyle}\slshape{attributestyle}{}ld
\makeatother

\lstloadlanguages{csh}

% general settings
\lstset{
  captionpos=t,
  showspaces=false,
  showtabs=false,
  keepspaces=true,
  breaklines=true,
  showstringspaces=false,
  breakatwhitespace=true,
  frame=bt,
  framesep=1em,
  columns=fullflexible,
  numbers=left,
  numberstyle=\tiny,
  escapeinside={(*@}{@*)},
  commentstyle=\color{codecolorcomment},
  keywordstyle=\color{codecolorkeyword}\bfseries,
  keywordstyle=[2]\color{codecolorkeywordcontrolflow}\bfseries,
  stringstyle=\color{codecolorstring},
  typestyle=\color{codecolorclass},
  enumstyle=\color{codecolorenum},
  basicstyle=\ttfamily\footnotesize
}

\lstnewenvironment{myListingCsharp}[5][]{% [extra], label, caption, types, enums
  \noindent\minipage{\linewidth}\medskip
  \lstset{
    caption={#3},
    label={#2},
    language=[Sharp]C,
    morekeywords={
      abstract, event, new, struct,
      as, explicit, null, switch,
      base, extern, object, this,
      bool, false, operator, [2]throw,
      break, finally, out, true,
      byte, fixed, override, try,
      case, float, params, typeof,
      catch, [2]for, private, uint,
      char, [2]foreach, protected, ulong,
      checked, [2]goto, public, unchecked,
      class, [2]if, readonly, unsafe,
      const, implicit, ref, ushort,
      continue, in, [2]return, using,
      decimal, int, sbyte, virtual,
      default, interface, sealed, volatile,
      delegate, internal, short, void,
      do, is, sizeof, [2]while,
      double, lock, stackalloc,
      [2]else, long, static, unmanaged, Cdecl,
      enum, namespace, string, var},
    types={#4},
    enums={#5},
    %morecomment=[s]{//}{},
    morestring=*[s]{\$"}{"},
    moredelim=*[s]{\{}{\}},
    moredelim=[is][\color{codecolormethod}]{|}{|},
    #1
  }
}{%
  \endminipage
}

\lstnewenvironment{myListingC}[5][]{% [extra], label, caption, types, enums
  \noindent\minipage{\linewidth}\medskip
  \lstset{
    caption={#3},
    label={#2},
    language=C,
    types={#4},
    enums={#5},
    moredelim=[is][\color{codecolormethod}]{|}{|},
    #1
  }
}{%
  \endminipage
}

% custom XML language to get nicer formatting
\lstdefinelanguage{XML}
{
  basicstyle=\ttfamily\footnotesize,
  morestring=[b]",
  morecomment=[s]{<?}{?>},
  morecomment=[s]{<!--}{-->},
  commentstyle=\color{codecolorcomment},
  moredelim=[s][\color{black}]{/>}{<},
  moredelim=[s][\color{black}]{/>}{</},
  moredelim=[s][\color{black}]{>}{<},
  stringstyle=\color{codecolorxmlstring},
  identifierstyle=\color{codecolorxmltag},
}
\lstnewenvironment{myListingXml}[4][]{% [extra], label, caption, attributes
  \noindent\minipage{\linewidth}\medskip
  \lstset{
    caption={#3},
    captionpos=t,
    label={#2},
    language=XML,
    attributes={#4},
    attributestyle=\color{codecolorxmlattribute},
    #1
  }
}{%
  \endminipage
}

% general listing env

\newenvironment{myListing}[3][]{% [short caption], label, caption
  \begin{mdframed}[
    hidealllines=true,
    innerleftmargin=0,
    innerrightmargin=0
    ]
  \myCaptionOf{lstlisting}{#1}{#3}\label{#2}% for tables, the caption goes to the top
  \hrule\medskip
}{%
  \medskip\hrule
  \end{mdframed}
}

% settings for uniform verbatim environment
\DefineVerbatimEnvironment%
{myVerbatim}{Verbatim}
{frame=lines,fontsize=\footnotesize,framesep=1em,commandchars=\&\{\}}

\newenvironment{myListingVerbatim}[3][]{% [short caption], label, caption
  \mdframed[
    hidealllines=true,
    innerleftmargin=0,
    innerrightmargin=0
    ]%
  \myCaptionOf{lstlisting}{#1}{#3}\label{#2}% for tables, the caption goes to the top
  \myVerbatim
}{%
  \endmyVerbatim%
  \endmdframed%
}


% typesetting binary and other values
\newcommand{\binary}[1]{0x#1}

% macros for typesetting various code elements inside the text

\definecolor{bytefieldunused}{gray}{0.8}

\newcommand{\class}[1]{{\color{codecolorclass}\texttt{#1}}}
\newcommand{\genericClass}[2]{\texttt{\class{#1}\allowbreak{}<{#2}>}}
\newcommand{\interface}[1]{{\color{codecolorclass}\texttt{#1}}}
\newcommand{\keyword}[1]{{\color{codecolorkeyword}\texttt{\textbf{#1}}}}
\newcommand{\namespace}[1]{{\color{codecolornamespace}\texttt{#1}}}
\newcommand{\method}[1]{{\color{codecolormethod}{\texttt{#1}}}}
\newcommand{\enum}[1]{{\color{codecolorenum}{\texttt{#1}}}}
\newcommand{\xmltag}[1]{{\color{codecolorxmltag}{\texttt{#1}}}}

% macros for typesetting description lists with C# class members

\newcommand{\propgetset}{\{ \keyword{get}; \keyword{set}; \}}
\newcommand{\propget}{\{ \keyword{get}; \}}
\newcommand{\propset}{\{ \keyword{set}; \}}

\newcommand{\ditemsrcsize}{\footnotesize}


\newcommand{\ditemmethodseparator}{}

% method to apply in \foreach macro
\def\makeditemmethodarg#1#2{\ditemmethodseparator{}{#2}%
\renewcommand{\ditemmethodseparator}{, }%set proper separator for the next item
}

\newcommand{\ditemmethodmakearglist}[1]{% arglist
\renewcommand{\ditemmethodseparator}{}%clear the separator for the first item
\foreach{\makeditemmethodarg}{}{#1}%
}

\newcommand{\ditemmethod}[4][\class]% [type typeset] type, name, args (types only)
  {\ditem{\ditemsrcsize\texttt{#1{#2} \method{#3}(#4)}}}

\newcommand{\ditemmethodWithComment}[5][\class]% [type typeset] type, name, args (types only)
  {\ditemWithComment{\ditemsrcsize\texttt{#1{#2} \method{#3}(#4)}}{#5}}

\newcommand{\ditemctor}[2]% name, args
  {\ditem{\ditemsrcsize\texttt{\method{#1}(#2)}}}

\newcommand{\ditemproperty}[4][\class]% [type typeset] type, name, get/set
  {\ditem{\ditemsrcsize\texttt{#1{#2} \method{#3} #4}}}

% checkmark
\def\checkmark{\tikz\fill[scale=0.4](0,.35) -- (.25,0) -- (1,.7) -- (.25,.15) -- cycle;}

% colored bitboxes for the bytefield package
\newcommand{\colorbitbox}[3]{%
\rlap{\bitbox{#2}{\color{#1}\rule{\width}{\height}}}%
\bitbox{#2}{#3}}
\newcommand{\colorwordbox}[4]{%
\rlap{\wordbox[#2]{#3}{\color{#1}\rule{\width}{\height}}}%
\wordbox[#2]{#3}{#4}}

% macro for typesetting transport parameters
\newcommand{\quicTransportParameter}[1]{\texttt{#1}}
\newcommand{\MaxAckDelay}{\quicTransportParameter{max\_ack\_delay}}
\newcommand{\MaxIdleTimeout}{\quicTransportParameter{max\_idle\_timeout}}
\newcommand{\DisableActiveMigration}{\quicTransportParameter{disable\_active\_migration}}

% macro for typesetting QUIC frame types
\newcommand{\quicFrame}[1]{\texttt{#1}}
\newcommand{\PADDING}{\quicFrame{PADDING}}
\newcommand{\PING}{\quicFrame{PING}}
\newcommand{\ACK}{\quicFrame{ACK}}
\newcommand{\RESETSTREAM}{\quicFrame{RESET\_STREAM}}
\newcommand{\STOPSENDING}{\quicFrame{STOP\_SENDING}}
\newcommand{\CRYPTO}{\quicFrame{CRYPTO}}
\newcommand{\NEWTOKEN}{\quicFrame{NEW\_TOKEN}}
\newcommand{\STREAM}{\quicFrame{STREAM}}
\newcommand{\MAXDATA}{\quicFrame{MAX\_DATA}}
\newcommand{\MAXSTREAMDATA}{\quicFrame{MAX\_STREAM\_DATA}}
\newcommand{\MAXSTREAMS}{\quicFrame{MAX\_STREAMS}}
\newcommand{\DATABLOCKED}{\quicFrame{DATA\_BLOCKED}}
\newcommand{\STREAMDATABLOCKED}{\quicFrame{STREAM\_DATA\_BLOCKED}}
\newcommand{\STREAMSBLOCKED}{\quicFrame{STREAMS\_BLOCKED}}
\newcommand{\NEWCONNECTIONID}{\quicFrame{NEW\_CONNECTION\_ID}}
\newcommand{\RETIRECONNECTIONID}{\quicFrame{RETIRE\_CONNECTION\_ID}}
\newcommand{\PATHCHALLENGE}{\quicFrame{PATH\_CHALLENGE}}
\newcommand{\PATHRESPONSE}{\quicFrame{PATH\_RESPONSE}}
\newcommand{\CONNECTIONCLOSE}{\quicFrame{CONNECTION\_CLOSE}}
\newcommand{\HANDSHAKEDONE}{\quicFrame{HAND\-SHAKE\_DONE}}

% macros for commonly used class names
\newcommand{\ValueTaskOf}[1]{\genericClass{ValueTask}{#1}}
\newcommand{\MemoryOf}[1]{\genericClass{Memory}{#1}}
\newcommand{\ListOf}[1]{\genericClass{List}{#1}}
\newcommand{\ReadOnlySpanOf}[1]{\genericClass{ReadOnlySpan}{#1}}
\newcommand{\ReadOnlyMemoryOf}[1]{\genericClass{ReadOnlyMemory}{#1}}
\newcommand{\ReadOnlySequenceOf}[1]{\genericClass{ReadOnlySequence}{#1}}
\newcommand{\SpanOf}[1]{\genericClass{Span}{#1}}
\newcommand{\ChannelOf}[1]{\genericClass{Channel}{#1}}
\newcommand{\ArrayOf}[1]{\texttt{{#1}[]}}
\newcommand{\FuncOf}[1]{\genericClass{Func}{#1}}
\newcommand{\ArrayPoolOf}[1]{\genericClass{ArrayPool}{#1}}

\newcommand{\byte}{\keyword{byte}}
\newcommand{\Long}{\keyword{long}}
\newcommand{\Int}{\keyword{int}}
\newcommand{\bool}{\keyword{bool}}
\newcommand{\CancellationToken}{\class{CancellationToken}}

\newcommand{\Socket}{\class{Socket}}
\newcommand{\Stream}{\class{Stream}}
\newcommand{\QuicSocketContext}{\class{Quic\allowbreak{}Socket\allowbreak{}Context}}
\newcommand{\QuicServerSocketContext}{\class{Quic\allowbreak{}Server\allowbreak{}Socket\allowbreak{}Context}}
\newcommand{\QuicClientSocketContext}{\class{Quic\allowbreak{}Client\allowbreak{}Socket\allowbreak{}Context}}
\newcommand{\QuicConnectionContext}{\class{Quic\allowbreak{}Connection\allowbreak{}Context}}
\newcommand{\QuicReader}{\class{Quic\allowbreak{}Reader}}
\newcommand{\QuicWriter}{\class{Quic\allowbreak{}Writer}}
\newcommand{\CryptoSeal}{\class{Crypto\allowbreak{}Seal}}
\newcommand{\CryptoSealAlgorithm}{\class{Crypto\allowbreak{}Seal\allowbreak{}Algorithm}}
\newcommand{\ReceiveStream}{\class{Receive\allowbreak{}Stream}}
\newcommand{\SendStream}{\class{Send\allowbreak{}Stream}}
\newcommand{\RecoveryController}{\class{Recovery\allowbreak{}Controller}}
\newcommand{\RangeSet}{\class{RangeSet}}
\newcommand{\PacketNumberWindow}{\class{Packet\allowbreak{}Number\allowbreak{}Window}}
\newcommand{\PacketNumberSpace}{\class{Packet\allowbreak{}Number\allowbreak{}Space}}
\newcommand{\SentPacket}{\class{Sent\allowbreak{}Packet}}
\newcommand{\ICongestionController}{\interface{ICongestion\allowbreak{}Controller}}
\newcommand{\StreamCollection}{\class{Stream\allowbreak{}Collection}}
\newcommand{\ProcessPacketResult}{\enum{Process\allowbreak{}Packet\allowbreak{}Result}}
\newcommand{\TransportParameters}{\class{Transport\allowbreak{}Parameters}}
\newcommand{\ITls}{\interface{ITls}}
\newcommand{\QuicTlsProvider}{\class{QuicTlsProvider}}
\newcommand{\MockQuicTlsProvider}{\class{MockQuicTlsProvider}}
\newcommand{\OpenSslQuicTlsProvider}{\class{OpenSslQuicTlsProvider}}
\newcommand{\OpenSslTls}{\class{OpenSslTls}}
\newcommand{\MockTls}{\class{MockTls}}
\newcommand{\SSL}{\class{SSL}}
\newcommand{\SSLCTX}{\class{SSL\_CTX}}
\newcommand{\QuicTestBase}{\genericClass{Quic\allowbreak{}Test\allowbreak{}Base}{\class{T}}}
\newcommand{\QuicListenerTests}{\genericClass{Quic\allowbreak{}Listener\allowbreak{}Tests}{\class{T}}}
\newcommand{\QuicConnectionTests}{\genericClass{Quic\allowbreak{}Connection\allowbreak{}Tests}{\class{T}}}
\newcommand{\QuicStreamTests}{\genericClass{Quic\allowbreak{}Stream\allowbreak{}Tests}{\class{T}}}
\newcommand{\ManualTransmissionQuicTestBase}{\class{Manual\allowbreak{}Transmission\allowbreak{}Quic\allowbreak{}Test\allowbreak{}Base}}

\newcommand{\QuicConnection}{\class{Quic\allowbreak{}Connection}}
\newcommand{\QuicStream}{\class{Quic\allowbreak{}Stream}}
\newcommand{\QuicListener}{\class{Quic\allowbreak{}Listener}}

\newcommand{\QuicConnectionProvider}{\class{Quic\allowbreak{}Connection\allowbreak{}Provider}}
\newcommand{\QuicStreamProvider}{\class{Quic\allowbreak{}Stream\allowbreak{}Provider}}
\newcommand{\QuicListenerProvider}{\class{Quic\allowbreak{}Listener\allowbreak{}Provider}}
\newcommand{\QuicImplementationProvider}{\class{Quic\allowbreak{}Implementation\allowbreak{}Provider}}
\newcommand{\QuicImplementationProviders}{\class{Quic\allowbreak{}Implementation\allowbreak{}Providers}}

\newcommand{\MsQuicConnection}{\class{MsQuic\allowbreak{}Connection}}
\newcommand{\MsQuicStream}{\class{MsQuic\allowbreak{}Stream}}
\newcommand{\MsQuicListener}{\class{MsQuic\allowbreak{}Listener}}
\newcommand{\MsQuicImplementationProvider}{\class{MsQuic\allowbreak{}Implementation\allowbreak{}Provider}}

\newcommand{\ManagedQuicConnection}{\class{Managed\allowbreak{}Quic\allowbreak{}Connection}}
\newcommand{\ManagedQuicStream}{\class{Managed\allowbreak{}Quic\allowbreak{}Stream}}
\newcommand{\ManagedQuicListener}{\class{Managed\allowbreak{}Quic\allowbreak{}Listener}}
\newcommand{\ManagedQuicImplementationProvider}{\class{Managed\allowbreak{}Quic\allowbreak{}Implementation\allowbreak{}Provider}}

\newcommand{\QuicListenerOptions}{\class{Quic\allowbreak{}Listener\allowbreak{}Options}}
\newcommand{\ServerAuthenticationOptions}{\class{ServerAuthenticationOptions}}
\newcommand{\QuicClientConnectionOptions}{\class{QuicClientConnectionOptions}}

% macro for the System.Net.Quic.dll library
\newcommand{\SystemNetQuicDll}{\texttt{System.\allowbreak{}Net.\allowbreak{}Quic.dll}}

% macros for typesetting paths and filenames
\newcommand{\filename}[1]{\path{#1}}


\makeglossaries

\newglossaryentry{managed-code}
{
  name=managed code,
  description={Code written in one of the \dotnet{} languages running on the \dotnet{} virtual machine.}
}

\newglossaryentry{native-code}
{
  name=native code,
  description={Code written in ahead-of-time compiled language. This code runs directly on the target CPU.}
}

\newglossaryentry{network-path}
{
    name=network path,
    description={An imaginary path between two network addresses. Each end of a network address consists of a pair of local address and port number.}
}

\newglossaryentry{head-of-line-blocking}
{
    name=head-of-line blocking,
    description={Performance limiting phenomenon caused by a line of packets being held up by the first packet}
}

\newglossaryentry{quic-packet}
{
    name=QUIC packet,
    description={A complete processable unit of QUIC that can be
  encapsulated in a UDP datagram.  Multiple QUIC packets can be
  encapsulated in a single UDP datagram.}
}

\newglossaryentry{out-of-order-packet}
{
  name=out-of-order pakcet,
  description={A packet that does not arrive directly after the packet that was
  sent before it.  A packet can arrive out of order if it is delayed, if earlier packets are
  lost or delayed, or if the sender intentionally skips a packet number.}
}

\newglossaryentry{cid}
{
  name=Connection ID,
  description={An opaque identifier that is used to identify a QUIC
  connection at an endpoint.  Each endpoint sets a value for its
  peer to include in packets sent towards the endpoint.}
}

\newglossaryentry{ack-eliciting-packet}
{
  name=ack-eliciting packet,
  description={A QUIC packet which contains at least one frame which is not \PADDING{},  \ACK{}, or \CONNECTIONCLOSE{}}
}

\newacronym{0rtt}{0-RTT}{Zero Round Trip Time Resumption}
\newacronym{scid}{SCID}{Source Connection ID}
\newacronym{dcid}{DCID}{Destination Connection ID}


% Title page and various mandatory informational pages
\begin{document}
%%% Title page of the thesis and other mandatory pages

%%% Title page of the thesis

\pagestyle{empty}
\hypersetup{pageanchor=false}
\begin{center}

\centerline{\mbox{\includegraphics[width=166mm]{img/logo-en.pdf}}}

\vspace{-8mm}
\vfill

{\bf\Large MASTER THESIS}

\vfill

{\LARGE\ThesisAuthor}

\vspace{15mm}

{\LARGE\bfseries{}QUIC Protocol Implementation\\for .NET}

\vfill

\Department

\vfill

{
\centerline{\vbox{\halign{\hbox to 0.45\hsize{\hfil #}&\hskip 0.5em\parbox[t]{0.45\hsize}{\raggedright #}\cr
Supervisor of the master thesis:&\Supervisor \cr
\noalign{\vspace{2mm}}
Study programme:&\StudyProgramme \cr
\noalign{\vspace{2mm}}
Study branch:&\StudyBranch \cr
}}}}

\vfill

% Zde doplňte rok
Prague \YearSubmitted

\end{center}

\newpage

%%% Here should be a bound sheet included -- a signed copy of the "master
%%% thesis assignment". This assignment is NOT a part of the electronic
%%% version of the thesis. DO NOT SCAN.

%%% A page with a solemn declaration to the master thesis

\openright
\hypersetup{pageanchor=true}
\pagestyle{plain}
\pagenumbering{roman}
\vglue 0pt plus 1fill

\noindent
I declare that I carried out this master thesis independently, and only with the cited
sources, literature, and other professional sources. It has not been used to obtain another
or the same degree.

\medskip\noindent
I understand that my work relates to the rights and obligations under the Act No.~121/2000 Sb.,
the Copyright Act, as amended, in particular the fact that the Charles
University has the right to conclude a license agreement on the use of this
work as a school work pursuant to Section 60 subsection 1 of the Copyright~Act.

\vspace{10mm}

\hbox{\hbox to 0.5\hsize{%
In \hbox to 6em{\dotfill} date \hbox to 6em{\dotfill}
\hss}\hbox to 0.5\hsize{\dotfill\quad}}
\smallskip
\hbox{\hbox to 0.5\hsize{}\hbox to 0.5\hsize{\hfil Author's signature\hfil}}

\vspace{20mm}
\newpage

%%% Dedication

\openright

\noindent
\Dedication

\newpage

%%% Mandatory information page of the thesis (EN)

\openright

\vbox to 0.5\vsize{
\setlength\parindent{0mm}
\setlength\parskip{5mm}

Title:
\ThesisTitle

Author:
\ThesisAuthor

\DeptType:
\Department

Supervisor:
\Supervisor, \SupervisorsDepartment

Abstract:
\Abstract

Keywords:
\Keywords

\vss}

\newpage

\openright
\pagestyle{plain}
\pagenumbering{arabic}
\setcounter{page}{1}


%%% A page with automatically generated table of contents of the master thesis

\tableofcontents

%%% Each chapter is kept in a separate file
\chapter{Introduction}

The internet, as we know it today, heavily relies on the use of the HTTP protocol. Is it used not
only by web browsers for interacting with web applications, but it may also be used at the server
side for communication between individual nodes in a cluster. Such use of HTTP is commonly seen in,
e.g., microservices architecture. Last but not least, HTTP is also used as a transport medium for
RESTful web APIs and technologies such as gRPC and GraphQL.

The latest version of the HTTP protocol is HTTP/2, published in 2015~\cite{rfc7540}. HTTP/2 improved
on its predecessor HTTP/1.1~\cite{rfc7230} by introducing request multiplexing over a single TCP
connection, header compression, and server-push features. These features led to reduced loading
times of web pages and generally improved the efficiency of the web~\cite{deSaxce2015}.

\section{Remaining Performance Issues in HTTP/2}

HTTP/2 does not solve all the performance issues HTTP/1.1 had. HTTP/2 introduced new features and
changes only to the topmost layer of the HTTP stack, and these changes could not address the
performance limiting phenomenons caused by the layers underneath. For the context of this thesis,
the following two performance problems are the most relevant.

\subsection*{Head-of-Line Blocking}

Request multiplexing was introduced in HTTP/2 to reduce the number of web servers' resources
required to serve requests made by web browsers. Because modern web pages are composed of many parts
(HTML, images, Javascript files, CSS style sheets), web browsers have to make multiple HTTP requests
to load the page. In HTTP/1.1, multiple parallel HTTP connections were made to download individual
parts of the webpage, one connection per HTTP request. In HTTP/2, a single connection can be used to
make multiple HTTP requests in parallel. Although this change improved HTTP's performance, its
current design is limited by a phenomenon known as \textit{\gls{head-of-line-blocking}}.

The individual HTTP frames which make up the HTTP requests are interleaved and transferred over TCP
in a single stream of data. When a TCP packet carrying part of this stream is lost, delivery of all
following packets is delayed until the lost packet is retransmitted. This will cause a delay in all
HTTP requests currently in progress, even if the lost packet contained data belonging only to a
single HTTP request. More information on head-of-line blocking can be found, e.g., on a dedicated
Wikipedia page~\cite{wiki:head-of-line-blocking}.

\subsection*{HTTPS Connection Establishment Latency}

HTTPS~\cite{rfc2818} is an extension to HTTP, which makes the connection encrypted by inserting the
TLS protocol layer between HTTP and TCP. Therefore, establishing an HTTPS connection requires first
establishing a TCP connection --- performing the three-way handshake --- and then performing another
separate handshake for the TLS layer. As can be seen on the illustration in
\autoref{fig:https-packets}, a minimum of three round trips is needed before the HTTP request can be
even sent.

\begin{myFigure} {fig:https-packets} {Packets sent during HTTPS connection establishment}

  \resizebox{0.8\linewidth}{!}{\input{img/01-https-handshake.pdf_tex}}

\end{myFigure}

More and more websites enforce the use of HTTPS in order to protect the privacy of their users. In
August 2020, 96 out of the top 100 viewed websites actively redirect to HTTPS, and more than 75\% of
all network traffic from the Chrome web browser used HTTPS~\cite{googleTransparency}. Even server to
server communication in the cloud is trending towards HTTPS to protect against partially compromised
networks. With HTTPS becoming the norm, almost all connections suffer from the increased latency
caused by the additional TLS handshake.

\section{HTTP/3 and QUIC}

The next version of the HTTP protocol --- HTTP/3~\cite{draft-ietf-quic-http} --- addresses the issues
mentioned above by replacing the TCP and TLS layers with a brand new UDP-based protocol named
QUIC\footnote{Originally intended as the acronym for Quick UDP Internet Connection. However, it has
been changed to be the protocol's actual name during the standardization process}. The QUIC protocol
features allow moving multiplexing capability from the application layer into the transport layer of
the HTTP protocol stack. The responsibilities and relationships between protocols on HTTP/2 and
HTTP/3 stacks are illustrated in \autoref{fig:http2-vs-http3-stack}.

\begin{myFigure} {fig:http2-vs-http3-stack} {Comparison between HTTP/2 and HTTP/3 protocol stacks}

  \input{img/01-http2-vs-http3-stack.pdf_tex}

\end{myFigure}

Although QUIC development is tied with that of HTTP/3, it is designed as a general-purpose transport
layer protocol that can be used for other application-layer protocols. The following points
summarize the main improvements of QUIC over TCP+TLS:

\begin{itemize}

    \litem{Stream multiplexing} QUIC provides an abstraction of multiple streams of data multiplexed
on a single connection. Moreover, because QUIC itself also implements loss detection, packet loss
can be managed in a way that reduces the scope of the \gls{head-of-line-blocking} problem described
above to the single affected stream.

    \litem{Faster connection establishment} QUIC also performs TLS handshake, but it does so in
parallel with the base protocol's handshake. The combined handshake requires fewer round-trips and
is, therefore, faster than the combination of TCP and TLS\@.

    QUIC also supports opt-in \gls{0rtt} from TLS 1.3. The 0-RTT mode of operation allows a client
to cache some session information, allowing it to send application layer data with the first packet
in future connections to the same server. 0-RTT effectively reduces the latency by another round
trip but makes connections vulnerable to repeat attacks. A more detailed description of 0-RTT can be
found online, e.g., on Cloudflare's blog~\cite{cloudflare-0rtt}.

    \litem{Always encrypted} TLS handshake is a mandatory part of connection establishment, and
encryption, therefore, cannot be turned off. This makes QUIC secure-by-default.

    \litem{Separating connection identity from peer's IP address} QUIC protocol does not use peers'
IP addresses to identify connections but instead uses Connection IDs, which are 8 to 20-byte
sequences negotiated during connection establishment.

    This makes QUIC very attractive for mobile devices, which can change IP addresses due to
switching between Wi-Fi and cellular data network. In TCP and --- by extension --- HTTP/2, the existing
connection must be terminated, and a new connection established from the new IP address. QUIC, on
the other hand, can migrate the connections in a way that is transparent to the application layer.

    This feature also enables QUIC extensions like Multipath
QUIC~\cite{draft-deconinck-quic-multipath-04}, which allows the simultaneous use of multiple network
interfaces for a single connection to achieve greater throughput.

\end{itemize}

As of August 2020, the specifications of HTTP/3 and QUIC are still at the draft stage, but the
standardization process is believed to be very close to complete. There are already multiple
implementations of QUIC being developed based on the draft versions of the standard. These
implementations are backed by big companies such as Google, Cloudflare, Facebook, and Microsoft.

Experiments with these implementations allowed early performance comparison between HTTP/3 and
HTTP/2, yielding promising results. In 2015, the Chromium team's experimental implementation showed
a 3\% improvement in mean page load time and 30\% fewer rebuffer events when watching YouTube
videos~\cite{Wilk2015}. Cloudflare launched preliminary support for HTTP/3 in April 2020 and has
measured a 12.4\% decrease in the \textit{time to the first byte} metric~\cite{Tellakula2020}, which
is consistent with the QUICs promise of reduced latency. In measurements done by Orange Labs, QUIC
protocol significantly outperforms TCP in unstable networks such as wireless mobile
networks~\cite{Cook2017}.

\section{Support for QUIC in \dotnet{}}

Microsoft's \dotnet{} development team has long-term plans to provide full support for QUIC and
HTTP/3 in \dotnet{}. Support for HTTP/3 should be completely transparent to users because the
implementation of \class{HttpClient} automatically chooses the best available HTTP
version~\cite{HttpClientDocs}. However, since QUIC can be used to build other protocols than HTTP/3,
its implementation will be exposed via public classes residing most likely inside the
\namespace{System.Net.Quic} namespace.

The QUIC support has been initially intended for the \dotnet{}~5 release planned for November 2020.
However, it turned out that the standardization process would not be completed in time for QUIC to
be implemented for the release. \dotnet{}~5, therefore, ships with only a preview (non-production
ready) HTTP/3 and QUIC support. A production-ready HTTP/3 and QUIC implementation was postponed
until \dotnet{}~6 release.

\subsection*{Existing QUIC Implementation in \dotnet{}}

The work-in-progress implementation of QUIC is still present in the master development branch of the
\dotnet{} runtime repository~\cite{dotnetGithub} but has been made accessible only internally. This
implementation is a wrapper around the \libmsquic{} library~\cite{msquicGithub}, which is a C
implementation of QUIC developed by Microsoft. The \libmsquic{} library is designed for
high-performance scenarios and has been recently made open-source.

The decision to use \libmsquic{} as the QUIC protocol implementation is not final. There are
compelling arguments for implementing the QUIC protocol in managed \dotnet{} code --- and, more
specifically, in \csharp{} --- for the production release.

The existing QUIC implementation in \dotnet{} uses a layer of indirection which allows multiple
implementations to exist side by side. In theory, this makes it possible to switch between QUIC
implementations at runtime, allowing performance comparisons between candidate QUIC implementations.

\subsection*{Motivation for Implementing QUIC in Managed Code}

Code written in ahead-of-time compiled languages such as C or C++ (referred to as
\textit{\gls{native-code}}) is likely to be faster than code written in \dotnet{} languages
(referred to as \textit{\gls{managed-code}}), which rely on the just-in-time compilation. However,
there are other aspects than raw performance to be considered when deciding to use native libraries
such as \libmsquic{}:

\begin{itemize}

    \litem{Cross-platform compatibility/availability} The \dotnet{} platform officially supports
multiple versions of Windows, macOS, and several Linux distributions. If the native library does not
support all these platforms, then the implementation must use an alternative library on other
platforms, introducing more complexity into the codebase and possible incompatibilities between
platforms.

    \litem{Support for different library versions} Currently, no native libraries that managed
\dotnet{} libraries depend on are part of the \dotnet{} distribution. Instead, \dotnet{} runtime
expects that the library is already installed on the target machine and can be dynamically loaded.
There is no way to enforce a specific version of the library, which means that the \dotnet{} code
must work correctly with multiple versions of the native library.

    \litem{Maintainability} Maintenance of the interop code requires the developer to read and
understand the language in which the native library is written. Debugging the code around the
language boundaries can be difficult if the necessary tooling for mixed-language debugging is not
available.

    \litem{Library's API} One of the aspects which greatly influences the resulting performance is
the similarity between the public API of the \dotnet{} library and the API of the native library.
Some aspects are the following:

    \begin{itemize}
      \item Event-based (using callbacks) vs.\ method based
      \item Blocking vs.\ non-blocking
      \item Synchronous vs.\ asynchronous
      \item Which side allocates memory buffers
    \end{itemize}

    \Gls{managed-code} needs to translate the differences between APIs, which may result in
noticeable performance overhead and even negate the performance gained from the native library's
use.

    \litem{Overhead of interop calls} When execution transitions between managed and unmanaged code,
the runtime must ensure the correct behavior of exceptions and garbage collector. The cost of the
so-called managed-to-native transition is relatively small. However, it may be noticeable when the
transition occurs very frequently.

    \litem{Future development} Exposing new features from newer versions of the native library can
be problematic because the code needs to work with older versions of the library.

\end{itemize}

In the past, the \dotnet{} development team has encountered multiple problems with the
\libcurl~\cite{curlGithub} library, which was used to implement HTTP request handling on macOS and
Linux. Different Linux distributions contained different versions of the \libcurl{} library and,
therefore, supported different features and had different bugs. The \dotnet{} development team had
to expend many resources to make sure the \gls{managed-code} written by other \dotnet{} developers
behaved consistently on all platforms.

Starting with \dotnet{} Core~2.1, the default implementation of HTTP request handling does not rely
on native libraries like \libcurl{}. Instead, the functionality has been rewritten in pure managed
code on top of sockets. This implementation offered better performance and consistent behavior
across all \dotnet{} platforms~\cite{SocketsHttpHandlerDocs}.

\section{Goals of this Thesis}\label{sec:01-goals}

The reasons mentioned in the previous section motivated Microsoft developers to consider
implementing QUIC in purely managed \csharp{} code. Before the final decision on which
implementation of QUIC will be used, it is necessary to investigate the performance characteristics
of managed QUIC implementation. This thesis seeks to create a partial implementation of QUIC whose
performance can be analyzed and --- if the managed implementation approach is found viable --- can form
the basis of the production QUIC implementation in future \dotnet{} versions.

Because this thesis's work can potentially become part of future \dotnet{}~6 release, the QUIC
implementation development will occur inside a development branch of the \dotnet{} runtime
repository. The result should be a compilable branch of the runtime, which exposes the QUIC API for
other \dotnet{} developers to use.

We have mentioned that the design of the QUIC API was interrupted early. Although the current
version does not expose all the QUIC protocol features, it is sufficient for evaluating our
implementation. This thesis will, therefore, avoid making any modifications to the API\@.

Implementing the full QUIC specification is outside the scope of a master thesis. Fortunately, to
evaluate the managed QUIC implementation's performance, many parts of the specification can be
omitted from our prototype. Therefore, this thesis will focus on the functionality needed to fully
support simple scenarios and leave out advanced features such as connection migration.

Because we do not expect readers to be thoroughly familiar with the QUIC protocol specification and
all its features, we will present an overview of the protocol in \autoref{chap:02-quic} and defer
the selection of the protocol features we will implement in this thesis to the beginning of the
analysis in \autoref{chap:03-analysis}. However, the implementation design should be such that the
rest of the specification can be implemented in the future.

\todo{these two paragraphs have changed}

It would make sense to evaluate our QUIC implementation's functionality by providing a partial
implementation of HTTP/3. However, even supporting the most straightforward GET requests would be
too complicated. Instead, we will implement a simple application in which server will echo all data
back to the client.

Because the indirect nature of the existing QUIC API could be used to alternate between
\libmsquic{}-based QUIC and our QUIC implementation at runtime. This would allow us to create a
benchmarking application to compare the performance of the two QUIC implementations. A small step
from there would allow us to compare the performance of the two QUIC implementations with that of
TCP+TLS-based \class{SslStream}. This benchmarking will be an optional goal of this thesis.

\subsection*{Summary of the Goals}

\todo{small changes in goals were done}

The following list summarizes the goals of this thesis presented in the previous subsection.

\begin{enumerate}

  \item Select a sufficient subset of QUIC specification needed to support the most basic data
  transfer and implement it inside \dotnet{} runtime codebase.

  \item Allow switching between the new managed implementation and the existing \libmsquic{}-based one.

  \item Evaluate the managed QUIC implementation by using it to implement a simple client-server
echo application.

  \item \textit{(optional)} Try to compare the performance of the new implementation with the
previous \libmsquic{}-based one and with TCP+TLS-based \class{SslStream}.

\end{enumerate}

\chapter{QUIC Protocol}\label{chap:02-quic}

This chapter is intended as a summary of the QUIC protocol specification and aims to provide
sufficient background for the design of our implementation. This text is based on version 29 of the
draft specification documents from June 2020, more specifically on the documents describing the core
transport protocol~\cite{draft-ietf-quic-transport}, TLS integration~\cite{draft-ietf-quic-tls}, and
congestion control mechanism~\cite{draft-ietf-quic-recovery}. Readers familiar with these documents
may skip this chapter.

We will start this chapter by first providing a high-level overview of QUIC and then providing a
more detailed description of the protocol's individual parts.

\section{Introduction}

QUIC protocol provides reliable and secure transport of multiple streams of data over a single
connection. QUIC provides following services:

\begin{itemize}

  \item Stream multiplexing

  \item Stream and connection-level flow control

  \item Low-latency connection establishment

  \item Connection migration and resilience to NAT rebinding

\end{itemize}

QUIC is implemented on top of UDP, which provides only unreliable transfer of datagrams. Therefore,
in addition to stream multiplexing, QUIC also implements loss recovery, congestion control,
transport security, and other features known from TCP or TLS protocols.

\subsection{Basic Concepts}

Throughout this chapter, and later when analysing the QUIC protocol specification, this text will
use several terms in a very specific meaning. This meaning is the same as defined in the
main specification document~\cite[Section~1.2]{draft-ietf-quic-transport}.

\todo{again, following list is almost verbatim from RFC}

\begin{description}

  \ditem{QUIC packet}  A complete processable unit of QUIC that can be
    encapsulated in a UDP datagram.  Multiple QUIC packets can be
    encapsulated in a single UDP datagram.

  \ditem{Out-of-order packet}  A packet that does not arrive directly after the packet that was
    sent before it.  A packet can arrive out of order if it is delayed, if earlier packets are
    lost or delayed, or if the sender intentionally skips a packet number.

   \ditem{Endpoint}  An entity that can participate in a QUIC connection by
      generating, receiving, and processing QUIC packets.  There are
      only two types of endpoint in QUIC: client and server.

   \ditem{Client}  The endpoint initiating a QUIC connection.

   \ditem{Server}  The endpoint accepting incoming QUIC connections.

  \ditem{Address}  When used without qualification, the tuple of IP version,
    IP address, UDP protocol, and UDP port number that represents one
    end of a network path.

  \ditem{Connection ID}  An opaque identifier that is used to identify a QUIC
    connection at an endpoint.  Each endpoint sets a value for its
    peer to include in packets sent towards the endpoint.

  \ditem{Stream}  A unidirectional or bidirectional channel of ordered bytes
    within a QUIC connection.  A QUIC connection can carry multiple
    simultaneous streams.

  \ditem{Application}  An entity that uses QUIC to send and receive data.

\end{description}

\subsection{Notational Conventions}

Numeric constants in hexadecimal always use network byte order (big-endian) and are prefixed by 0x.
In case of constants of non-conventional integer sizes, such as 2-bit integers used in packet
fields, the notation will omit unnecessary leading zeros. For example, \binary{FE} represents
decimal number 254 and \binary{2} represents decimal number 2.

This chapter contains graphical diagrams of several elements such as QUIC packets. In those
schematics, individual fields include length information as follows:

\begin{itemize}

  \litem{X (A)} Indicates that X is A bits long.

  \litem{X (i)} Indicates that X is encoded using the QUIC's variable-length encoding. This encoding is described in \autoref{sec:02-variable-length-encoding}.

  \litem{X (A..B)} Indicates that X can be any length from A to B, either of these limits can be omitted to indicate minimum size of zero, or no upper size limit, respectively.

  \litem{X (A) = C} Indicates that X has a fixed value C.

  \litem{X (A) ...} Indicates that X has is repeated zero or more times and each instance is A bits long.

\end{itemize}

\subsection{Wire Encoding}\label{sec:02-wire-encoding}

The process of encoding QUIC packets to be sent via the network is optimized for size. All values
sent over the network are encoded in big endian order. There are two nontrivial encodings used:
variable-length integer encoding and packet number encoding.

\subsubsection{Variable-Length Integer Encoding}\label{sec:02-variable-length-encoding}

Almost all numeric values stored in QUIC packets are encoded using a \textit{variable-length integer
  encoding}. This encoding uses the first byte's two most significant bits to encode whether the
value is encoded as 1, 2, 4, or 8-byte integer. This encoding supports only positive numbers. The
ranges available for individual encoding lengths are listed in \autoref{tab:02-quic-varint-length}.

\begin{myTable} {tab:02-quic-varint-length} {Variable-length integer encoding lengths}
  {ccr}
  {Most significant bits & Encoding length (B) & Maximum value}
  \binary{0}             & 1                   & \num{63}         \\
  \binary{1}             & 2                   & \num{16383}      \\
  \binary{2}             & 4                   & \num{1073741824} \\
  \binary{3}             & 8                   & $2^{62}-1$       \\
\end{myTable}

In order to optimize the size of QUIC packets, the implementations are encouraged to always choose
the shortest encoding necessary to represent the given number.

\subsubsection{Packet Number Encoding}\label{sec:02-packet-number-encoding}

\todo{check this for legibility}

Packets in QUIC are sequentially numbered and, therefore, packet numbers received by an endpoint
form a mostly ascending sequence with occasional reordering or gaps due to network unreliability.
This causes the packet number of a newly received packet to be close to that of the previously
received packet. QUIC leverages this fact by sending only the lower bytes of the packet number. The
number of bytes sent is variable between one to four bytes and is always chosen so that the receiver
can reconstruct the original packet number using the previously received packet numbers and the
bytes from the arriving packet.

When determining how many bytes of the packet number need to be included in the packet, the sender
uses the value of the highest acknowledged packet number. The encoding length must use enough bytes
to be able to represent at least twice the difference between the current packet number and the
highest acknowledged number.

As an example, consider following situation: as sender prepares to send packet \binary{4417}, its
highest acknowledged packet number is \binary{43a0}. The difference between the two numbers is
\binary{77} which is still small enough to require only the least significant byte (\binary{17}) of
the packet number to be sent in the packet. Suppose that by the time the packet arrives at the
receiver, its highest received packet number is \binary{43f5}. Based on the encoded packet number
length, the deconstructed packet number must be from the range \binary{4375}--\binary{4475}. The only
number from this range which matches the least significant byte sent in the packet is \binary{4417}
which is the correct packet number.

\section{QUIC Packets}

In QUIC connection, endpoints exchange \textit{QUIC packets} enclosed in UDP datagrams. A single UDP
datagram can contain multiple packets, although in most cases they contain only one. Packing
multiple QUIC packets into a single UDP datagram is called \textit{coalescing}. A single QUIC packet
cannot span multiple UDP datagrams. QUIC uses a total of six different packet types:

\todo{The order of the packet here now matches the subsubsections later, but I think different order here might be better}

\begin{enumerate}

  \litem{Version Negotiation} sent by the server when the client tries to establish a connection
  with QUIC version that is not supported by the server;

  \litem{Retry} optionally used by servers for Client Address Validation \todo{forward reference
  address validation} when establishing new connections;

  \litem{Initial} used to initiate a new connection and exchange the initial information;

  \litem{Handshake} used during the connection handshake;

  \litem{0-RTT} when TLS 1.3 0-RTT mode of operation is enabled, 0-RTT packets carry \textit{early
  data} --- application data that is sent before the TLS handshake is complete in order to reduce
  latency; and

  \litem{1-RTT} main packet type used throughout the lifetime of QUIC connection. 1-RTT and 0-RTT
  packets are the only packet types which carry application data.

\end{enumerate}

Version Negotiation and Retry packets are sent as one-time responses in special scenarios and are
not individually numbered. QUIC is designed in a way that lets the server send these packets without
maintaining any state for the connection. These packets are not used in an established connection.

Packets of the other types --- Initial, Handshake, 0-RTT and 1-RTT --- are individually numbered. The
packet numbers do not form a single sequence, as is the case with TCP packets. Instead, QUIC
organizes these packets into three separate \textit{packet number spaces}, and uses a separate
sequence of numbers for each packet number space. These packet number spaces are:

\begin{enumerate}

  \litem{Initial} Used for initiating new connections and exchanging initial cryptographic
  information. Contains only Initial packets.

  \litem{Handshake} Used during the connection handshake process. Contains only Handshake packets.

  \litem{Application} Used throughout the lifetime of the connection to transfer application data.
  This packet space contains both 0-RTT and 1-RTT packets.

\end{enumerate}

In addition to being numbered separately, packets from each packet number space are processed
completely independently of the other packet number spaces. For example, Initial packets can be
acknowledged only by another Initial packet. In their payload, the Initial, Handshake, 0-RTT and
1-RTT packets carry \textit{QUIC Frames} which are low level protocol messages carrying e.g.
acknowledgements and parts of streams sent by the application.

Each packet type has a header and a payload. QUIC tries to optimize the packet encoding in order to
maximize the amount of application data that can be sent in a single UDP datagram. Therefore, two
types of packet headers exist. 1-RTT packets use a \textit{short header} and all other packet types
use the \textit{long header}. The short header contains only a subset of fields present in the long
header and omits data which are relevant only during connection establishment, like protocol version
identifier.

The type of the header is specified by the most significant bit of the first byte in the packet ---
the \textit{Header Form} bit. The second most significant bit --- the \textit{Fixed Bit} --- is always
set to 1 in valid QUIC packets. Following sections describe the structure of the rest of the packet
for individual packet types.

\newcommand{\longFieldHeight}{1}

\subsection{Long Packet Header}

Long headers contain information that is necessary for the connection establishment. It includes
both the Source and Destination Connection IDs. It also contains \textit{Version} and \textit{Packet
  Type} fields. A few leftover bits of the header are reserved for type-specific information. The
structure of the long header is illustrated in \autoref{fig:02-long-header}.

\begin{myFigure}{fig:02-long-header}{Long packet header structure}

  \begin{bytefield}[bitwidth=2.5em]{8}
    \bitheader{0-7} \\
    \bitbox{1}{H=1} & \bitbox{1}{F=1} & \bitbox{2}{T (2)} & \bitbox{4}{Type Specific Bits (4)} \\
    \wordbox[tlrb]{\longFieldHeight}{Version (32)} \\
    \wordbox[tlrb]{1}{Destination Connection ID Length (8)} \\
    \wordbox[tlrb]{\longFieldHeight}{Destination Connection ID (..) } \\
    \wordbox[tlrb]{1}{Source Connection ID Length (8)} \\
    \wordbox[tlrb]{\longFieldHeight}{Source Connection ID (..) } \\
  \end{bytefield}

  H = Header Form, F = Fixed Bit, T = Packet Type

\end{myFigure}

The semantics of the individual fields of the long header is as follows:

\begin{description}

    \ditem{Header Form Bit} used to distinguish between packets with so called \textit{long header}
    and \textit{short header} format. In case of the long header packets. the form bit is set to \binary{1}.

    \ditem{Fixed Bit} A bit that is always set to \binary{1} in valid QUIC packets.

    \ditem{Packet Type} Discriminator of the packet type. the possible values are listed in \autoref{tab:02-packet-type}.

\begin{myTable}{tab:02-packet-type}{Values of the Packet Type field in long packet header}
  {cc}
  {Packet Type & Value}
  Initial      & \binary{0} \\
  0-RTT        & \binary{1} \\
  Handshake    & \binary{2} \\
  Retry        & \binary{3} \\
\end{myTable}

    \ditem{Version} Indicates which version of QUIC is in use. Location of this field will be
    the same across all QUIC versions However, the structure of the rest of the packet may be different
    in future versions of QUIC.

    \ditem{Destination Connection ID Length} Length of the Destination Connection ID field.

    \ditem{Destination Connection ID}  The Connection ID issued by the recipient of the packet.

    \ditem{Source Connection ID Length}  Length of the Source Connection ID field.

    \ditem{Source Connection ID} The Connection ID issued by the sender of the packet.

    \ditem{Reserved Bits} Bits reserved for use in the future QUIC versions. In the initial QUIC
    version, these bits must be set to 0.

\end{description}

\subsection{Version Negotiation Packet}\label{sec:02-version-negotiation-packet}

The Version Negotiation packet is sent by the server when it receives a long header packet
requesting a version of QUIC unsupported by the server. As an exception to other long header
packets, Version Negotiation packet is not discriminated by a specific value in the Packet Type
field in the header, but by a special value \binary{00000000} in the Version field.

After the header, the packet contains a list of supported versions, each one listed as 32-bit
integer in big-endian. Also, since future QUIC versions may allow larger Connection IDs than 20
bytes, a valid Version Negotiation packet can contain up to 255-byte Connection IDs. The structure
of the Version Negotiation packet is illustrated in \autoref{fig:02-version-negotiation-packet}.

\begin{myFigure}{fig:02-version-negotiation-packet}{Version Negotiation packet structure}

  \begin{bytefield}[bitwidth=2.5em]{8}
    \bitheader{0-7} \\
    \begin{rightwordgroup}{long \\ header}
      \bitbox{1}{H=1} & \colorbitbox{bytefieldunused}{7}{Unused} \\
      \wordbox[tlrb]{\longFieldHeight}{Version (32) = \binary{00000000}} \\
      \wordbox[tlrb]{1}{Destination Connection ID Length (8)} \\
      \wordbox[tlrb]{\longFieldHeight}{Destination Connection ID (0..2040) } \\
      \wordbox[tlrb]{1}{Source Connection ID Length (8)} \\
      \wordbox[tlrb]{\longFieldHeight}{Source Connection ID (0..2040) }
    \end{rightwordgroup} \\
    \wordbox[tlrb]{\longFieldHeight}{Supported Version 1 (32)} \\
    \wordbox[tlrb]{\longFieldHeight}{Supported Version 2 (32)} \\
    \wordbox[tlrb]{1}{...} \\
  \end{bytefield}

  H = Header Form

\end{myFigure}

The identifier of the initial QUIC version is \binary{00000001}. There are also special identifiers
for draft versions of the protocol valid only until the QUIC specification is finalized.

\subsection{Retry Packet}\label{sec:02-retry-packet}

Retry packets are distinguished by value \binary{3} in the \textit{Packet Type} field of the long
header. Retry packets are sent by a server as part of the optional \textit{Address Validation}
mechanism used to protect against traffic amplification attack. \todo{forward ref to address
  validation?} The structure of the packet is illustrated in \autoref{fig:02-retry-packet}.

\begin{myFigure}{fig:02-retry-packet}{Retry packet structure}

  \begin{bytefield}[bitwidth=2.5em]{8}
    \bitheader{0-7} \\
    \begin{rightwordgroup}{long \\ header}
      \bitbox{1}{H=1} & \bitbox{1}{F=1} & \bitbox{2}{T=\binary{3}} & \colorbitbox{bytefieldunused}{4}{Unused} \\
      \wordbox[tlrb]{\longFieldHeight}{Version (32)} \\
      \wordbox[tlrb]{1}{Destination Connection ID Length (8)} \\
      \wordbox[tlrb]{\longFieldHeight}{Destination Connection ID (0..160) } \\
      \wordbox[tlrb]{1}{Source Connection ID Length (8)} \\
      \wordbox[tlrb]{\longFieldHeight}{Source Connection ID (0..160) }
    \end{rightwordgroup} \\
    \wordbox[tlrb]{\longFieldHeight}{Retry Token (..)} \\
    \wordbox[tlrb]{\longFieldHeight}{Retry Integrity Tag (128)} \\
  \end{bytefield}

  H = Header Form, F = Fixed Bit, T = Packet Type

\end{myFigure}

The semantics of the fields specific to the Retry packet are:

\begin{description}

    \ditem{Retry Token} Contains an opaque token generated by the server. This token must be
    echoed back in an Initial packet during next connection attempt. The contents of the token must be
    difficult to guess by the attacker and verifyable in a stateless manner (i.e. without saving it in
    memory for future comparisons).

    \ditem{Retry Integrity Tag} Tag used to check the integrity of the packet. For the details on
    how the ingegrity tag is calculated, refer to the TLS integration
    specification~\cite[Section~5.8]{draft-ietf-quic-tls}.

\end{description}

\subsection{Initial, Handshake and 0-RTT Packets}

Initial, Handshake and 0-RTT packets are almost identical in structure. All three types use the
\textit{Type-Specific Bits} from the long header to store \textit{Reserved Bits} and the
\textit{Packet Number Length}. After the long packet header, these packets contain the
\textit{Length} field containing the packet length, \textit{Packet Number} field and the actual
payload consisting of QUIC frames. The only exception to this structure is the Initial packet, which
contains two additional fields just after the long header: \textit{Token Length} and \textit{Token},
which are used to carry the Retry Token from the Retry packet in case Address Validation is
requested by the server. The structure of these three packets is illustrated in
\autoref{fig:02-initial-handshake-0rtt-packets}.

\begin{myFigure}{fig:02-initial-handshake-0rtt-packets}{Structure of the Initial, Handshake, and 0-RTT packets}

  \begin{bytefield}[bitwidth=2.5em]{8}
    \bitheader{0-7} \\
    \begin{rightwordgroup}{long \\ header}
      \bitbox{1}{H=1} & \bitbox{1}{F=1} & \bitbox{2}{T (2)} & \bitbox{2}{R=\binary{0}} & \bitbox{2}{L (2)} \\
      \wordbox[tlrb]{\longFieldHeight}{Version (32)} \\
      \wordbox[tlrb]{1}{Destination Connection ID Length (8)} \\
      \wordbox[tlrb]{\longFieldHeight}{Destination Connection ID (0..160) } \\
      \wordbox[tlrb]{1}{Source Connection ID Length (8)} \\
      \wordbox[tlrb]{\longFieldHeight}{Source Connection ID (0..160) }
    \end{rightwordgroup} \\
    \begin{leftwordgroup}{Initial \\ only}
      \colorbitbox{bytefieldunused}{8}{Token Length (i)} \\
      \colorwordbox{bytefieldunused}{tlrb}{\longFieldHeight}{Token (..)}
    \end{leftwordgroup} \\
    \wordbox[tlrb]{1}{Length (i)} \\
    \wordbox[tlrb]{1}{Packet Number (8..32)} \\
    \wordbox[tlrb]{\longFieldHeight}{Packet Payload (..)} \\
    \wordbox[tlrb]{\longFieldHeight}{Integrity Tag (16)} \\
  \end{bytefield}

  H = Header Form, F = Fixed Bit, T = Packet Type,\\
  R = Reserved Bits, L = Packet Number Length

\end{myFigure}

The semantics of the new fields in these packets are:

\begin{description}

    \ditemWithComment{Token Length}{Initial only} Length of the Token field.

    \ditemWithComment{Token}{Initial only} Contains an opaque token if the server provided
    one in the Retry packet as part of Address Validation.

    \ditem{Reserved Bits} Bits reserved for use in the future QUIC versions. In the initial QUIC
    version, these bits must be set to 0.

    \ditem{Packet Number Length} The length of the encoding used for the packet number.

    \ditem{Length} The length of the remainder of the packet. This includes the Packet Number,
    the payload of the packet, and --- for encrypted packets --- the AEAD integrity tag.

    \ditem{Packet Number} The sequence number of this packet in the respective packet number space.
    This field uses a special encoding which is described in \autoref{sec:02-packet-number-encoding}.

    \ditem{Packet Payload} Serialized sequence of QUIC frames.

    \ditem{Integrity Tag} Opaque checksum produced by the AEAD cipher during packet encryption.
    Packet encryption is described in detail in \autoref{sec:02-packet-protection}.

\end{description}

\subsection{1-RTT Packet}

1-RTT packets are the only packet that uses the short header to make more space for application data
in the UDP datagram. Besides the Header Form, Fixed, Reserved Bits and Packet Number Length fields
which have the same meaning as in the long header, short header contains \textit{Spin Bit} and
\textit{Key Phase Bit} fields.

1-RTT packets can be sent only after the connection has been successfully established. This implies
that connection IDs used for the connection by both endpoints are already known, and there is no
need to repeat the Source Connection ID or specify the length of the Destination Connection ID.
1-RTT packets also lack the Length field. It is assumed that 1-RTT packets fill the rest of the UDP
datagram. Therefore, after the short header, only the Packet Number field and the QUIC frame payload
follows. The structure of the 1-RTT frame is illustrated in \autoref{fig:02-1rtt-packet}.

\begin{myFigure}{fig:02-1rtt-packet}{1-RTT packet structure}

  \begin{bytefield}[bitwidth=2.5em]{8}
    \bitheader{0-7} \\
    \begin{rightwordgroup}{short \\ header}
      \bitbox{1}{H=0} & \bitbox{1}{F=1} & \bitbox{1}{S} & \bitbox{2}{R=\binary{0}} & \bitbox{1}{K} & \bitbox{2}{L (2)} \\
      \wordbox[tlrb]{\longFieldHeight}{Destination Connection ID (0..160) }
    \end{rightwordgroup} \\
    \wordbox[tlrb]{1}{Packet Number (8..32)} \\
    \wordbox[tlrb]{\longFieldHeight}{Packet Payload (..)} \\
    \wordbox[tlrb]{\longFieldHeight}{Integrity Tag (16)} \\
  \end{bytefield}

  H = Header Form, F = Fixed Bit, S = Spin Bit, \\
  R = Reserved Bits, K = Key Phase Bit, L = Packet Number Length

\end{myFigure}

The semantics of the fields specific to 1-RTT packets are:

\begin{description}

    \ditem{Spin Bit} A bit used for an optional QUIC feature which allows on-path nodes to measure
    connection latency by observing changes in this bit. \todo{describe the SPIN bit feature?}

    \ditem{Key Phase Bit} A bit used to communicate that the keys used for the packet encryption
    need to be updated. Updating the keys is necessary after using them to encrypt certain number of
    packets to provide stronger protection. The mechanism of Key Update is described in
    \autoref{sec:02-key-update}.

\end{description}

\section{QUIC Frames}\label{sec:02-quic-frames}

QUIC frames are low-level QUIC protocol messages carried in the payload of Initial, Handshake, 0-RTT
and 1-RTT packets. Examples of these frames include, e.g., \ACK{} frames carrying acknowledgments
for received packets, \STREAM{} frames carrying the application data, and \CRYPTO{} frames carrying
data for the TLS handshake.

During the lifetime of the connection, all QUIC packets have to be acknowledged by sending ACK in
another packet in the same packet number space. However, not all packets have to be acknowledged
immediately, e.g., acknowledging packets containing only ACK would cause an endless flood of ACK
packets. Instead, the ACK frame is sent later together with a more urgent data. Frame types that
require immediate acknowledgments are called \textit{ack-eliciting frames}, and the packets with at
least one such frame are called \textit{ack-eliciting packets}.

Because packets in differnet packet number spaces offer different level of confidentiality, not all
frames can be sent in any packet type. For example, application data in \STREAM{} frames cannot be
sent in Initial and Handshake to avoid compromising security. \autoref{tab:02-frame-types} lists all
frame types, whether they are ack-eliciting and in which packets they can be sent.

\begin{myTable}[\small] {tab:02-frame-types} {QUIC frame types}
  {l@{\hskip -0.1in}ccccc}
  {                     &               & \multicolumn{4}{c}{Allowed in packet type} \\ \cmidrule(lr){3-6}
    Frame type          & Ack-eliciting & Initial & Handshake & 0-RTT & 1-RTT}
  \PADDING{}            &               & \checkmark{}     & \checkmark{}       & \checkmark{}   & \checkmark{}          \\
  \PING{}               & \checkmark{}           & \checkmark{}     & \checkmark{}       & \checkmark{}   & \checkmark{}          \\
  \ACK{}                &               & \checkmark{}     & \checkmark{}       &       & \checkmark{}          \\
  \RESETSTREAM{}        & \checkmark{}           &         &           & \checkmark{}   & \checkmark{}          \\
  \STOPSENDING{}        & \checkmark{}           &         &           & \checkmark{}   & \checkmark{}          \\
  \CRYPTO{}             & \checkmark{}           & \checkmark{}     & \checkmark{}       &       & \checkmark{}          \\
  \NEWTOKEN{}           & \checkmark{}           &         &           &       & \checkmark{}          \\
  \STREAM{}             & \checkmark{}           &         &           & \checkmark{}   & \checkmark{}          \\
  \MAXDATA{}            & \checkmark{}           &         &           & \checkmark{}   & \checkmark{}          \\
  \MAXSTREAMDATA{}      & \checkmark{}           &         &           & \checkmark{}   & \checkmark{}          \\
  \MAXSTREAMS{}         & \checkmark{}           &         &           & \checkmark{}   & \checkmark{}          \\
  \DATABLOCKED{}        & \checkmark{}           &         &           & \checkmark{}   & \checkmark{}          \\
  \STREAMDATABLOCKED{}  & \checkmark{}           &         &           & \checkmark{}   & \checkmark{}          \\
  \STREAMSBLOCKED{}     & \checkmark{}           &         &           & \checkmark{}   & \checkmark{}          \\
  \NEWCONNECTIONID{}    & \checkmark{}           &         &           & \checkmark{}   & \checkmark{}          \\
  \RETIRECONNECTIONID{} & \checkmark{}           &         &           & \checkmark{}   & \checkmark{}          \\
  \PATHCHALLENGE{}      & \checkmark{}           &         &           & \checkmark{}   & \checkmark{}          \\
  \PATHRESPONSE{}       & \checkmark{}           &         &           & \checkmark{}   & \checkmark{}          \\
  \CONNECTIONCLOSE{}    &               & \checkmark{}     & \checkmark{}       & \checkmark{}   & \checkmark{}          \\
  \HANDSHAKEDONE{}      & \checkmark{}           &         &           &       & \checkmark{}          \\
\end{myTable}

\todo{Short description of each packet type and what it is used for?}


\section{QUIC Connection}

\todo{expand this part} This section describes the details of connection lifetime.

\subsection{Connection ID}

Traditional network protocols use the combination of remote endpoint IP address and port to identify
the connection. QUIC, on the other hand, uses a dedicated Connection ID identifier. This in essence
enables migrating the connection to different network paths and interfaces, e.g. from cellular data
to a Wi-Fi network, because the connection identity does not depend on the peer's IP address.

Connection IDs are opaque byte sequences between 8 to 20 bytes in length. Each endpoint in a QUIC
connection independently selects Connection IDs it will use to identify the QUIC connection. These
Connection IDs are then used to populate the Source Connection ID and Destination Connection ID of
the QUIC packets.

The first pair of Connection IDs is exchanged during connection establishment. Additional Connection
IDs can be issued independently by each endpoint during the lifetime of the connection. These
additional Connection IDs are primarily used when migrating the connection to a new network path.
QUIC requires that the same Connection ID be used only on one network path. When migrating a
connection, the endpoints must start using different Connection ID for sending packets to prevent
correlation of the network traffic by external observers.

Connection IDs can be also retired by an endpoint. By retiring a Connection ID, the endpoint
communicates that it will no longer use the Connection ID, and the other endpoint should drop any
incoming packets that use it. Retiring a Connection ID serves as a request to the peer to issue a
new Connection ID as a replacement.

Alternatively, endpoints can use a zero-length Connection ID. In that case, the connection identity
is tied to the IP address and port of the remote endpoint. Using zero-length Connection ID saves
space in the sent datagrams, but imposes several limitations on the connection. For example, if an
endpoint uses a zero-length Connection ID, it cannot issue additional connection IDs and, therefore,
it cannot migrate a connection to a new local address.

\subsection{Matching Packets to Connections}

When a packet arrives to an endpoint, it needs to be associated with an existing connection, or ---
for servers --- potentially initiate a new connection. If packet contains a non-zero-length Connection
ID in the DCID field, the Connection ID is used to find an existing connection. If the packet uses a zero-length conneciton id, then the local and remote addresses determine the target connection.

\autoref{fig:02-connection-multiplexing} illustrates how server endpoint processes incoming packets
from client connections. Server maintains a table mapping between DCIDs corresponding connections,
and dispatches the incoming packets accordingly. If a packet cannot be associated with an existing
connection, then it may be a new connection attempt, otherwise the packet is discarded.

\begin{myFigure}{fig:02-connection-multiplexing}{Multiple QUIC connections on the same machine port}

\input{img/02-socket-multiplexing.pdf_tex}

\todo{Check if the image is more clear now}

\end{myFigure}

In case the packet cannot be associated with an existing connections, client endpoints simply ingore
the packet. Server, on the other hand, check the packet type and version of the protocol it
requires. For valid Initial packets with supported verions, server proceeds with the handshake as
described in the previous section. For invalid Initial packets, the server responds with an Initial
packet containing a \CONNECTIONCLOSE{} frame signaling the refusal of the connection.

In case the client packet requests an unsupported QUIC version, server replies with a Version
Negotiation packet (described in \autoref{sec:02-version-negotiation-packet}). After receiving such
a packet, the client can try again using one of the supported versions. Figure
\autoref{fig:02-version-negotiation-flow} illustrates such an exchange. In the figure, the client
tries to establish connection using an unsupported version denoted by X. Server then replies with
Version Negotiation packet listing versions 1,2, and 3. Client the tries to establish connection
again using the version 1.

\begin{myFigure}{fig:02-version-negotiation-flow}{Example Version Negotiation packet exchange}

\resizebox{\linewidth}{!}{\input{img/02-version-negotiation-flow.pdf_tex}}

\end{myFigure}

Lastly, the server may choose to send a Stateless Reset (see later in
\autoref{sec:02-stateless-reset}) for any other packet that cannot be matched to an existing
connection.

\subsection{QUIC Transport Parameters}\label{sec:02-transport-parameters}

QUIC has several options for parameterizing the connection. These are called \textit{Transport
  Parameters} and they essentially represent constraints for the other endpoint. Each endpoint sets
the transport parameters for the other endpoint. QUIC leverages the extensibility of the TLS
protocol to exchange transport parameters during the connection handshake.

\todo{should I mention them by name? e.g. I mention max\_idle\_timeout later}
Example transport parameters are:

\begin{itemize}

  \item Initial flow control limits

  \item Whether connection migration is allowed

  \item Maximum delay before sending an acknowledgment for ack-eliciting packets

  \item Maximum idle timeout before the connection is silently closed

  \item Maximum size of a UDP packet the endpoint is willing to receive

\end{itemize}

Many of the parameters have a default value, which is used when the given transport parameter is not
sent. Other transport parameters are mandatory. The exhaustive list of the transport parameters can be found in the core transport specification~\cite[Section~7.4]{draft-ietf-quic-transport}.

\subsection{Connection Establishment}

In order to initiate a new connection, clients sends an Initial packet to the server, which
initiates the handshake process. After the handshake completes, both peers have derived protection
keys necessary to send and receive 1-RTT packets with application data.

An example handshake flow is illustrated in \autoref{fig:02-example-handshake-flow}. The figure
shows client and server endpoints exchanging UDP datagrams with QUIC packets. For illustration, we
list the contents of the individual QUIC frames. The figure also lists contents of the \CRYPTO{}
frames sent. However, these are only for illustrative purposes because QUIC does not interpret the
\CRYPTO{} frames' contents.

\begin{myFigure} {fig:02-example-handshake-flow} {Example QUIC handshake flow}

\resizebox{\linewidth}{!}{\input{img/02-handshake-flow.pdf_tex}}

\end{myFigure}

In its first datagram, the client sends an Initial packet with initial information consisting of a
single \CRYPTO{}(CH) frame. This frame contains the \textit{Client Hello} TLS message.

The server replies with a datagram containing three coalesced QUIC packets. The first is an Initial
packet that acknowledges the client's Initial packet using the \ACK{}(0) frame, and a \CRYPTO{}(SH)
frame with a \textit{Server Hello} message. The contents of Client Hello and Server Hello messages
are used by TLS to derive Handshake protection keys. The server advances the TLS handshake further
by sending another \CRYPTO{} frame in the Handshake packet. The server also has enough information
to derive also the 1-RTT keys, so it can also start sending data on Stream 1 using the \STREAM{}(1,
``...'') frame in a 1-RTT packet.

Because the server's Initial packet contained an ack-eliciting \CRYPTO{}(SH) frame, the client needs
to acknowledge it by sending an Initial frame with an \ACK{}(0) frame. The client can posesses both
Client and Server Hello mesages and can derive the Handshake keys, enabling him to process the
server's Handshake packet. The Handshake packet needs to be separately acknowledged by another
\ACK{}(0) frame, and a reply from the TLS layer must be sent using another \CRYPTO{} frame. From the
information in the server's Handshake packet's \CRYPTO{} frame, the client derives 1-RTT protection
keys and processes the server's 1-RTT packet. In addition to sending an \ACK{}(0) frame for the
server's 1-RTT packet, the client can now start sending application data on stream 0 using a
\STREAM{}(0, ``...'') frame.

After server detects that Handshake has successfully completed, it sends a \HANDSHAKEDONE{} frame to
communicate the fact to the client.

\todo{retry? or do it as part of Security section?}

\subsection{Connection Termination}

QUIC connection can be terminated in three ways:

\begin{itemize}

  \item Idle timeout

  \item Immediate close

  \item Stateless reset

\end{itemize}

\subsubsection{Idle Timeout}\label{sec:idle-timeout}

If idle timeout is enabled, the endpoint silently closes the connection if it does not receive a
packet from the peer for a specified time period. Each peer may advertise a timeout period using the
\MaxIdleTimeout{} transport parameter, but the effective value is the minimum of the two
values.

In order to prevent timeouts, endpoints can send a \PING{} or another ack-eliciting frame to test
the liveness of the connection. However, sending \PING{} frames should be initiated by the
application protocol, not QUIC implementation, to prevent unnecessary network traffic.

\subsubsection{Immediate Close}

An immediate close can be initiated both by QUIC implementation and by the application protocol.
Either endpoint can initiate an immediate close by sending a \CONNECTIONCLOSE{} frame. By sending a
\CONNECTIONCLOSE{} frame, the peer enters a \textit{closing state}, in which it includes the
\CONNECTIONCLOSE{} frame in all packets it sends in reply to incoming packets. The closing state is
received also if the endpoint receives a \CONNECTIONCLOSE{} frame from the peer. In that case the
endpoint also echoes the \CONNECTIONCLOSE{} frame back to the other endpoint.

The closing state lasts until the endpoint is sure the other endpoint is also in the closing state ---
e.g., until it also receives \CONNECTIONCLOSE{} --- or until a \textit{closing timeout} expires. The
closing timeout period is calculated from the current estimate of the round-trip time of the
connection.

The \CONNECTIONCLOSE{} frame carries an error code and, optionally, a human-readable error phrase.
When initiated by QUIC, the error codes semantics defined by the QUIC specification. However, when
initiated by the application protocol, the semantics of all possible error code values sent in the
frame are defined by the application protocol itself. This implies that the implementations must
require an error code when closing the connection and do not provide a default error code value.

\subsubsection{Stateless Reset}\label{sec:02-stateless-reset}

A stateless reset is an option of last resort for an endpoint that does not have access to the state
of a connection, possibly resulting from a crash or outage. An endpoint may send a stateless reset
in response to receiving a packet that it cannot associate with an active connection.

In such cases, the endpoint sends a specially crafted packet that ends with a Stateless Reset Token
associated with his Connection ID\@. The Stateless Reset Token requirements are quite complex, and
we encourage readers to read the full specification if they are interested in details~\cite[Section~10.4]{draft-ietf-quic-transport}.

\subsection{Connection Migration}

A novel feature of QUIC is the ability to migrate connections to a different network path. This is
enabled by using a dedicated Connection ID, instead of using the endpoint's address to identify the
connection. In the initial QUIC versio, only client endpoints can migrate the connection to a
different address.

Prior to migrating a connection, the endpoint can optionally check the reachability of the other
endpoint using the process called \textit{path validation}. Path validation consists of exchanging
\textit{probing packets}, and is described in \autoref{sec:02-path-validation}.

If the path is validated, the endpoint can migrate the connection by starting to send packets from
the mew local address. After server receives the first non-probing packet from the new client's
address, it starts sending all future packets to that new address.

In order to prevent network traffic being correlated by the outside observers, QUIC requires each
Connection ID to be used for only one combination of local and remote endpoint addresses. Therefore,
when connection is migrated, both endpoints must switch to using different Connection IDs.

The connection migration process is illustrated in \autoref{fig:02-connection-migration-flow}.
Client and servers use Connection IDs \textbf{C1} and \textbf{S1}, respectively. Client first probes
the reachability of the server with a probing packet containing a \PATHCHALLENGE{} frame. Because
this packet is sent from a different local address, it uses a different Connection ID (\textbf{S2})
that was issued previously by the server. Likewise, server uses different Connection ID \textbf{C2}
to reply with a packet containing a \PATHRESPONSE{} frame, confirming the reachability from the
clients new local address. After that, client migrates the connection by sending all packets via the
new local address. After receiving next non-probing packet from the new address, server switches to
the new client's address as well. Connection IDs \textbf{C1} and \textbf{S1} are no longer used in
the rest of the connection and can be retired later.

\begin{myFigure}{fig:02-connection-migration-flow}{Flow of packets during connection migration}

\resizebox{\linewidth}{!}{\input{img/02-connection-migration-flow.pdf_tex}}

\end{myFigure}

\section{QUIC Streams}

QUIC can transport multiple streams of data in a single connection. Each stream is identified by its
\textit{Stream ID} and is processed independently of the other streams. Each QUIC packet can carry
data for one or more QUIC streams. \autoref{fig:02-stream-multiplexing} illustrates how QUIC may
pack two streams into frames such that those streams are transported in parallel.

\begin{myFigure}{fig:02-stream-multiplexing}{Stream multiplexing in QUIC}

  \input{img/02-stream-multiplexing.pdf_tex}

\end{myFigure}

\subsection{Streams Types}

Streams transported by QUIC can be either unidirectional or bidirectional. Unidirectional streams
carry data from the initiator to its peer, and bidirectional streams carry data in both directions.
Both client and server can open new streams. QUIC recognizes four types of streams, and the type of
the stream is encoded in the two least significant bits of the Stream ID. The stream types and their
associated encoding is summarized in \autoref{tab:02-stream-id-type-map}.

\begin{myTable} {tab:02-stream-id-type-map} {Mapping of QUIC Stream types to Stream ID bits}
  {cc}
  {Stream type                     & Least significant bits}
  Client-Initiated, Bidirectional  & \binary{0} \\
  Server-Initiated, Bidirectional  & \binary{1} \\
  Client-Initiated, Unidirectional & \binary{2} \\
  Server-Initiated, Unidirectional & \binary{3} \\
\end{myTable}

Bidirectional streams can be viewed as the combination of two unidirectional streams. After opening
the stream, each direction of the stream behaves as separate inbound and outbound unidirectional
streams. This implies that the sending and receiving part of the stream can be closed independently
of each other.

\subsection{Stream Lifetime}

Opening a stream does no require any special action. Streams are opened simply by sending the first
\STREAM{} frame carrying data for that stream. However, streams of a particular type can be opened
only in ascending order of their Stream IDs. For example, stream with ID 2 must be opened before
opening stream 6. Sending data for higher numbered streams will automatically open all lower
numbered streams of the same stream type.

Streams can be closed either gracefully or abortively. Graceful stream close is signaled by a
\textit{Fin} bit in the \STREAM{} frame, signaling that data carried by this packet are the last
part of the stream. The stream is gracefully closed once all stream data is confirmed received by
the other endpoint.

Abortive stream close is achieved using the \RESETSTREAM{} frame and, therefore, this action is also
reffered to as \textit{resetting the stream}. Streams can be reset only by the sender. Receiver can
request abortive stream close by sending a \STOPSENDING{} frame if it no longer wishes to receive
data on that stream. Both \RESETSTREAM{} and \STOPSENDING{} frames carry an application-level error
code, that is then read by the application using QUIC.

After the stream is closed, it's Stream ID cannot be reused. Instead, the next available Stream ID
must be used. QUIC uses the variable-length integer encoding (see
\autoref{sec:02-variable-length-encoding}) and, therefore, there is no shortage of available stream
IDs, which range from 0 to $2^{62}-1$.

\subsection{Required Operations on Streams}

The implementation should provide the following operations on sending part of the stream:

\begin{itemize}

  \item write data;

  \item end the stream by specifying that all data has been written; and

  \item terminate the stream with an application-level error code.

\end{itemize}

On receiving part of the stream, application protocols must to be able to:

\begin{itemize}

  \item read data; and

  \item abort reading with an application-level error code.

\end{itemize}

\section{Flow Control}

QUIC aims to be a general-purpose transport protocol to be used over a potentially untrusted
network, and as such, it needs to protect endpoints from malicious peers. To prevent malicious
senders from exhausting all available memory on the receiver by sending large amounts of data, or
fast senders from overwhelming slow receivers, QUIC employs a credit-based flow control scheme.

All QUIC streams are flow controlled both individually and together as an aggregate. Each endpoint
also controls the number of streams the other peer is allowed to open. All flow control limits are
communicated to the peer using three types of frames:

\begin{itemize}

\litem{\MAXSTREAMDATA{}} maximum offset of data sent on a stream with specified ID\@.

\litem{\MAXDATA{}} maximum sum of all offsets of data sent on all streams.

\litem{\MAXSTREAMS{}} maximum number of opened streams of a particular stream type.

\end{itemize}

Endpoints can only increase the flow control limits. Their peers must ignore any attempts to
decrease the flow control limits to ensure consistency when two consecutive QUIC packets with flow
control updates are reordered during transit. In case the peer violates any of the control flow
limits mentioned above, the QUIC implementation must immediately terminate the connection.

\section{Loss Detection and Recovery}

Because UDP is an unreliable transport protocol, QUIC must implement measures to recover from packet
loss. The packet loss detection is implemented similarly to TPC --- each endpoint sends
acknowledgments for each received packet. However, an essential difference from TCP is that QUIC
endpoints do not retransmit entire lost packets using the same packet number. Instead, each QUIC
frame in the original packet is updated and sent in some future packet, or dropped altogether if the
information contained in the frame is no longer relevant.

The acknowledgement is communicated using the \ACK{} frame which contains ranges of received packet
numbers that the endpoint acknowledges. A packet number can sent multiple times (in \ACK{} frames in
multiple packets) until the endpoint can determine that the other endpoint received the
acknowledgement.

For packets containing an \textit{ack-eliciting} frame (see \autoref{sec:02-quic-frames}),
acknowledgements must be sent before the delay specified by the \MaxAckDelay{} transport
parameter. For other packets, such as packets containing only \ACK{} or \PADDING{} frames, the
acknowledgement can be delayed until an ack-eliciting packet is received.

\autoref{fig:02-packet-loss-example} illustrates the loss detection and retransmission process in
action. When server receives the acknowledgment for packet 3, but not for packet 2 sent earlier, it
infers that packet 2 never reached the receiver and retransmits the payload \textbf{B} in packet 4.
The server does not have to retransmit the \ACK{} for packet 1 because it was already sent in packet
3 which was already acknowledged by the client. Therefore, packet 4 includes ACK only for packet 2.

\begin{myFigure}{fig:02-packet-loss-example}{Loss detection and retransmission example}

\input{img/02-retransmission-example.pdf_tex}

\end{myFigure}

The exact criteria for a packet to be deemed lost by a QUIC endpoint are following:

\begin{enumerate}

  \item The packet was not acknowledged.

  \item A packet which was sent later has been acknowledged.

  \item Either the packet has been sent long enough in the past, or its packet number is
sufficiently smaller than the highest acknowledged packet number.

\end{enumerate}

The third condition is dependent on the values of particular constants. The specification recommends
that the packet is considered lost if the gap between its packet number and the highest acknowledged
packet number is at least 3, or if it was sent for longer than $9/8$ times the estimate of the
current round-trip time.

The first condition mentioned above requires receipt of a packet from the peer to declare any packet
as lost. However, the loss of the last packet in a sequence could go undetected because there is no
following packet that can be acknowledged. To avoid possible deadlocks in such scenarios, QUIC
endpoint sends up to two ack-eliciting \textit{probe packets} if it does not receive a packet from a
peer in a period called \textit{probe timeout} (PTO for short). The PTO duration doubles each time
probe packets are sent until either a reply is received or the connection is terminated due to idle
timeout (see \autoref{sec:idle-timeout})

Similarly to TCP, QUIC also uses congestion control to manage the \textit{congestion window} --- the
amount of data that can be in-flight. The selection of the congestion control algorithm is left on
the implementation. As an example, The QUIC specification document for loss detection and recovery~\cite[Section~7]{draft-ietf-quic-recovery} describes a congestion control algorithm similar to TCP NewReno~\cite{rfc6582} algorithm.

\section{Security}

This section describes the mechanisms used to ensure security of the protocol. Besides encrypting
all packets sent throughout the lifetime of the connection, QUIC uses additional mechanisms to
ensure that the servers using the protocol are resistant to denial-of-service and other
cyber-attacks.

\subsection{TLS Integration}

Instead of designing a new handshake protocol, QUIC offloads the encryption negotiation to TLS
protocol (more precisely, TLS version 1.3). The low-level messages used in TLS, such as
\textit{Server Hello} and \textit{Client Hello} are transported by QUIC inside \CRYPTO{} frames and
passed to a TLS implementation on the other side. This way, QUIC is able to offer the same
confidentiality level as conventional TLS connections.

The TLS protocol is extensible, among the standard extensions which are also used by QUIC are
Application Level Protocol Negotiation~\cite{rfc7301} (ALPN for short) and Server Name
Indication~\cite{rfc6066} (SNI for short).

Application Level Protocol Negotiation is used when multiple application protocols are supported on
the same TCP or UDP port. ALPN allows the application layer to negotiate --- as part of the TLS
handshake --- which application protocol will be used in the established connection.

Server Name Indication is used by clients to specify the hostname of the server to which they are
connecting. When multiple websites are hosted on the same IP address and port, SNI allows the server
to customize the security configuration for each hosted website. During connection establishment,
proper security configuration, such as the SSL certificate to be used, can be selected based on the
value of the hostname.

A custom TLS extension is used by QUIC to exchange transport parameters during the handshake. The
transport parameters are described in \autoref{sec:02-transport-parameters}.

\subsection{Packet Protection}\label{sec:02-packet-protection}

All QUIC packets of type Initial, Handshake, 1-RTT and 0-RTT are encrypted to ensure integrity and
confidentiality of the transmitted data. Negotiation of the cryptographic ciphers and the encryption
keys is handled by the TLS handshake. This section focuses on how the negotiated encryption is
applied to QUIC packets.

\subsubsection{Authenticated Encryption with Associated Data}

QUIC uses type of encryption called Authenticated encryption with associated data~\cite{rfc5116}
(AEAD for short). This type of encryption ensures both confidentiality and authenticity of the
encrypted data. In addition to encrypted data --- called \textit{ciphertext} --- AEAD encryption
outputs also an authentication tag which is used to check the integrity of the payload during
decryption. The encryption can be authenticated by supplying additional authentication data (AAD for
short) which are not encrypted, but influence the authentication tag and, therefore, must be
supplied also during decryption. As an additional protection, AEAD also accepts a \textit{nonce}
parameter, which is an additional input which is supposed to be unique for each encrypted packet.

The programming interface for AEAD provides following operations:

\begin{itemize}

  \item Encryption:

  \begin{itemize}

    \item input: plaintext, key, nonce, AAD (optional)

    \item output: ciphertext, authentication tag

  \end{itemize}

  \item Decryption:

  \begin{itemize}

    \item input: ciphertext, key, nonce, authentication tag, AAD (if provided during encryption)

    \item output: plaintext or error if the authentication tag does not match the supplied
      ciphertext and AAD

  \end{itemize}

\end{itemize}

\subsubsection{Deriving QUIC Protection Keys}\label{sec:02-encryption-key-derivation}

QUIC derives multiple distinct keys from the secrets negotiated by TLS handshake. The derivation
uses the HKDF-Expand-Label function~\cite{rfc5869} to derive keys of desired length. The keys are
derived as follows:

\begin{equation*}
  \begin{split}
  key & = \operatorname{HKDF-Expand-Label}(secret, \texttt{"quic key"}, \texttt{""}, 32) \\
  iv  & = \operatorname{HKDF-Expand-Label}(secret, \texttt{"quic iv"}, \texttt{""}, 12)  \\
  hp  & = \operatorname{HKDF-Expand-Label}(secret, \texttt{"quic hp"}, \texttt{""}, 32)  \\
  \end{split}
\end{equation*}

The following subsections describe how the $key$, $iv$, and $hp$ keys are used in the actual process
of packet encryption.

\subsubsection{Packet Protection Procedure}

When actually encrypting the packets, QUIC first encrypts the packet payload using the AEAD cipher
negotiated by the TLS implementation. QUIC specification allows use of all AEAD ciphers allowed in
TLS 1.3. These ciphers are:

\begin{itemize}

  \item TLS\_AES\_128\_GCM\_SHA256

  \item TLS\_AES\_256\_GCM\_SHA384

  \item TLS\_CHACHA20\_POLY1305\_SHA256

  \item TLS\_AES\_128\_CCM\_SHA256

  \item TLS\_AES\_128\_CCM\_8\_SHA256

\end{itemize}


The parameters for AEAD for packet payload protection are:

\begin{itemize}

    \litem{key} the $key$ derived in previous section

    \litem{nonce} the $iv$ derived in previous section, last 8 bytes XORed with the packet number

    \litem{AAD} the contents of the packet header

    \litem{plaintext} the payload of the packet.

\end{itemize}

The produced authentication tag is appended to the protected payload and is included in the payload
size.

QUIC also protects the header of the packet. The header protection mechanism is more complex than
that of the payload. The process requires calculating the \textit{header protection mask}, which is
then applied using XOR to parts of the packet header. The details of the header protection mask
calculation depend on the negotiated cipher, but it always uses the $hp$ key and a 16-byte sample of
the encrypted payload.

The parts of the packets protected by the payload encryption and header
protection mechanisms are illustrated in \autoref{fig:02-protected-fields}.


\begin{myFigure}{fig:02-protected-fields}{Fields protected by payload encryption and header protection}

  \newcommand{\legendsquare}[1]{%
    \textcolor{#1}{\rule{0.7em}{0.7em}}%
  }

  \definecolor{hp}{rgb}{0.7, 1, 0.7}
  \definecolor{pp}{rgb}{0.7, 0.7, 1}

  \legendsquare{hp}~Header protection \hspace{1cm} \legendsquare{pp}~Payload encryption

  \vspace{5mm}

  \begin{mySubFigure}{\textwidth}{fig:02-protected-fields-long}{Long header packets}

    \begin{bytefield}[bitwidth=2.5em]{8}
      \bitheader{0-7} \\
      \begin{rightwordgroup}{long \\ header}
        \bitbox{1}{H=1} & \bitbox{1}{F=1} & \bitbox{2}{T (2)} & \colorbitbox{hp}{2}{R=\binary{0}} & \colorbitbox{hp}{2}{L (2)} \\
        \wordbox[tlrb]{\longFieldHeight}{Version (32)} \\
        \wordbox[tlrb]{1}{Destination Connection ID Length (8)} \\
        \wordbox[tlrb]{\longFieldHeight}{Destination Connection ID (0..160) } \\
        \wordbox[tlrb]{1}{Source Connection ID Length (8)} \\
        \wordbox[tlrb]{\longFieldHeight}{Source Connection ID (0..160) }
      \end{rightwordgroup} \\
      \begin{leftwordgroup}{Initial \\ only}
        \colorbitbox{bytefieldunused}{8}{Token Length (i)} \\
        \colorwordbox{bytefieldunused}{tlrb}{\longFieldHeight}{Token (..)}
      \end{leftwordgroup} \\
      \wordbox[tlrb]{1}{Length (i)} \\
      \colorwordbox{hp}{ltrb}{\longFieldHeight}{Packet Number (8..32)} \\
      \colorwordbox{pp}{tlrb}{\longFieldHeight}{Packet Payload (..)} \\
      \wordbox[tlrb]{\longFieldHeight}{Integrity Tag (16)} \\
    \end{bytefield}

  \end{mySubFigure}

  \vspace{5mm}

  \begin{mySubFigure}{\textwidth}{fig:02-protected-fields-short}{Short header packets}

    % this one has to be off-center to make the image aligned with the top one
    \hspace{1.8cm}\begin{bytefield}[bitwidth=2.5em]{8}
      \bitheader{0-7} \\
      \begin{rightwordgroup}{short \\ header}
        \bitbox{1}{H=0} & \bitbox{1}{F=1} & \bitbox{1}{S} & \colorbitbox{hp}{2}{R=\binary{0}} & \colorbitbox{hp}{1}{K} & \colorbitbox{hp}{2}{L (2)} \\
        \wordbox[tlrb]{\longFieldHeight}{Destination Connection ID (0..160) }
      \end{rightwordgroup} \\
      \colorwordbox{hp}{ltrb}{\longFieldHeight}{Packet Number (8..32)} \\
      \colorwordbox{pp}{tlrb}{\longFieldHeight}{Packet Payload (..)} \\
      \wordbox[tlrb]{\longFieldHeight}{Integrity Tag (16)} \\
    \end{bytefield}

  \end{mySubFigure}

\end{myFigure}

\subsubsection{Updating 1-RTT Protection Keys}\label{sec:02-key-update}

Some AEAD functions have limits for how many packets can be encrypted using the same keys. Analysis
in the specification document shows that updating the protection keys after sending at most $2^{23}$
packets provides the same confidentiality level as a standard TLS
connection~\cite[Appendix~B]{draft-ietf-quic-tls}.

The key update is signalled by flipping a \textit{Key Phase} bit in the 1-RTT packet header. And
after observing the change in the \textit{Key Phase}, an endpoint derives new secret and $key$ and
$iv$ keys using the HKDF-Expand-Label function:

\begin{equation*}
  \begin{split}
  secret_{n+1} & = \operatorname{HKDF-Expand-Label}(secret_{n}, \texttt{"quic ku"}, \texttt{""}, 32) \\
  key_{n+1} & = \operatorname{HKDF-Expand-Label}(secret_{n+1}, \texttt{"quic key"}, \texttt{""}, 32) \\
  iv_{n+1}  & = \operatorname{HKDF-Expand-Label}(secret_{n+1}, \texttt{"quic iv"}, \texttt{""}, 12)  \\
  \end{split}
\end{equation*}

Because the \textit{Key Phase} bit is protected by header protection, the $hp$ key must remain
unchanged to ensure that the other endpoint can correctly remove the header protection.

\subsection{Client Address Validation}\label{sec:02-address-validation}

After receiving the first Initial packet from a new client, server can request address validation by
sending a Retry packet (see \autoref{sec:02-retry-packet}). The Retry packet carries a token which
must be echoed back to the server in all following Initial packets. As long as an attacker cannot
generate a valid token for its address and the client is able to return that token, this exchange
proves to the server that the client has received the token.

\autoref{fig:02-client-address-validation-flow} illustrates the use of Retry packet to valididate
client address during connection establishment. By default, clients do not fill the Token field of
the Initial packet. The server rejects the initial connection attempt and issues a Retry Token
(denoted ``ABCD'' in the figure). Client then tries again with another Initial with the provided
token and the server proceeds with the usual handshake.

\begin{myFigure}{fig:02-client-address-validatin-flow}{Client address validation using a Retry packet.}

\resizebox{\linewidth}{!}{\input{img/02-client-address-validation-flow.pdf_tex}}

\end{myFigure}

\subsection{Path Validation}\label{sec:02-path-validation}

QUIC is layered on top of UDP which is a connection-less protocol. This means that changes in
endpoint address can also happen without active migration on the endpoint's part, e.g., because of
NAT rebinding along the network path.

The endpoint address may also be spoofed by the other endpoint in an attempt to perform
\textit{traffic amplification attack} against the spoofed address. For this reason, the amount of
data sent to the new endpoint address must be limited until path validation determines that the
address belongs to the endpoint. This path validation is performed by sending a \textit{probing
  packet} containing a \PATHCHALLENGE{} frame with an unpredictable token. The other endpoint must
echo the token back in a \PATHRESPONSE{} frame. After that the new address is considered validated
and the sending rate restrictions are lifted.

\chapter{Analysis}\label{chap:03-analysis}

This chapter analyzes the protocol and selects the necessary subset needed to evaluate a \dotnet{}
implementation of the QUIC protocol. Afterward, we design the architecture and outline the
implementation of its major parts. Lastly, we investigate the means for unit testing and debugging
of the implementation.

\section{Implemented Feature Subset Selection}\label{sec:03-feature-selection}

The decision which features of the QUIC protocol we implement is guided by the goals we set in
\autoref{sec:01-goals}. For feature selection, these goals can be rephrased into the following:

\begin{enumerate}

  \item Support basic data transport, enabling some experimentation with QUIC as the transport for
application layer protocols. \textit{(goals 1 and 2)}

  \item Enable performance measurements that are representative of the potential full QUIC
implementation. \textit{(goals 1 and 3)}

\end{enumerate}

The first goal requires full implementation of the multiplexed stream abstraction as defined by the
QUIC specification. It also requires implementing loss detection and recovery to ensure that no data
gets lost during the transport.

In order to get representative performance measurements, we should implement all performance
affecting features. The most important are packet protection and flow control because they influence
performance throughout the lifetime of the QUIC connection. To summarize, the thesis should
implement at least the following features:

\begin{itemize}

    \item Connection lifetime support (establishment, termination)

    \item Stream multiplexing

    \item Packet protection

    \item Loss detection and recovery

    \item Flow control

\end{itemize}

On the other hand, we can disregard many QUIC features that react to one-time events or do not
otherwise influence the implementation's performance. These features include:

\begin{itemize}

    \item Connection migration, and therefore multiple Connection IDs support

    \item Complex (token-based) address validation

    \item Network path MTU detection

    \item Version negotiation

    \item 1-RTT key updates

    \item Advanced security measures (see Section 21 in transport
specification~\cite{draft-ietf-quic-transport})

\end{itemize}

The rest of the features form a grey area that can be implemented fully, partially, or even not at
all if convenient.

\section{Design Considerations}

Before we start the actual analysis, we will briefly outline the design principles used for the
actual design of the implementation and their rationale.

\subsection{Performance}

One of the critical factors in the decision between managed \dotnet{} implementation of QUIC or
using an external library like \libmsquic{} is performance. Therefore, the implementation design
decisions should focused on greater performance, possibly sacrificing maintainability if the
trade-off was justified.

As a general rule, the implementation will:

\begin{itemize}

    \litem{Avoid excessive heap allocations} Although heap allocation is cheap in \dotnet{}, the
actual price is paid during garbage collection. The more objects are allocated, the more frequent
garbage collections are. Each collection introduces a small stall into the program which could
disrupt internal timing of the QUIC implementation. Therefore, heap allocations on hot paths of the
code executions should be minimized.

    \litem{Prefer return codes over exceptions} Throwing an exception is an expensive operation, and
their frequent use would have negative impact on the performance.

\end{itemize}

\subsection{Testability}

The second design aspect we focusd on is the testability of the implementation. Ideally, the design
would minimize the need for live debugging of the implementation. This is especially important
because stopping the implementation on a breakpoint inevitably disrupts the connection, possibly
leading to termination because of idle timeout (see \autoref{sec:02-idle-timeout}).

The design intention was to allow writing deterministic automated tests that can inspect the packets
sent by the endpoint and verify that they are consistent with the behavior defined by the QUIC
specification.

\section{Target \dotnet{} API}

The public API, which was designed by the \dotnet{} team for use by other developers, uses similar
concepts like TCP and UDP \dotnet{} networking APIs. There is a \QuicListener{} class intended to be
used by servers for listening for incoming connections, similarly to using the \class{TcpListener}
class. The actual connection is represented by the \QuicConnection{} class, which can be compared to
\class{TcpClient} class. The individual QUIC streams are exposed using the \QuicStream{} class,
which implements the \Stream{} abstraction.

The API is expected to be used asynchronously using the \keyword{async}/\keyword{await} model. Most
methods return a \ValueTask{} which should be \keyword{await}ed to efficiently wait until the
operation completes. The full list of methods with detailed description can be found later in
\autoref{sec:06-api}.

\section{High-Level Architecture}\label{sec:03-high-level-architecture}

The \QuicConnection{} implementation will require some sort of background processing
thread\footnote{In order to reduce verbosity, this text will be using the term \textit{thread} to
mean both a dedicated \class{Thread} instance or a \class{Task} running on a thread-pool thread.} to
send acknowledgments for incoming packets and to resend data after being determined lost due to a
timeout. Because the correctness of a multithreaded code is hard to test, our implementation of
\QuicConnection{} will not interface with the underlying \Socket{} instance directly. Instead, the
implementation will provide an internal interface for exchanging the datagrams to be sent/received.
The actual sending of these datagrams will be handled by a separate class named
\QuicSocketContext{}. By separating the socket I/O from the connection logic implementation, we can
write unit tests that inspect the sent datagrams and assert that their content conforms to the
protocol specification.

On the server's side, the \QuicSocketContext{} class will also handle any stateless packet
processing, such as sending Retry and Version Negotiation packet. It will be also responsible for
matching packets to the appropriate \QuicConnection{} instances and queueing new connections to
\QuicListener{} to be read by the application. \autoref{fig:03-architecture} illustrates the
relationship between \QuicSocketContext{}, \QuicConnection{} and \QuicListener{} classes.

\begin{myFigure}{fig:03-architecture}{High-level background processing architecture.}

\resizebox{0.76\textwidth}{!}{\input{img/03-architecture.pdf_tex}}

\end{myFigure}

\subsection{Servicing the Socket}

The \QuicSocketContext{} class will implement the necessary background processing management. Its
responsibilies are:

\begin{itemize}

  \item routing incoming datagrams from \Socket{} to the proper \QuicConnection{} instance;

  \item calling timeout handlers after a timeout set by the connection expires; and

  \item sending outgoing datagrams provided by the \QuicConnection{}.

\end{itemize}

For performance reasons, it may be better to process timeouts and send the datagrams using one
thread and process received datagrams using another thread. However, the cost of synchronization of
the shared state may outweigh the performance gain. Our prototype implementation will use a single
thread to service both types of events but will allow for the possibility of future experimentation
with separate threads for sending and receiving.

\subsubsection{Processing Multiple Connections in Parallel}

For client connections, the background processing provided by the \QuicSocketContext{} class needs
to service only a single connection. However, on the server's side, there can be multiple
connections receiving datagrams from the same socket and, therefore, being served by a single
\QuicSocketContext{} instance. Servicing multiple \QuicConnection{} instances by only one thread
could limit the server's throughput.

In order to allow processing multiple QUIC connections in parallel, our implementation will use a
dedicated background \Task{} for each \QuicConnection{} instance. Our implementation separates the
per-connection logic into separate \QuicConnectionContext{} class. The \QuicSocketContext{} class
will be directly responsible only for stateless packet processing like sending a Version Negotiation
packet for incoming packets with unsupported versions. Packets belonging to existing connections
will be queued for processing by the appropriate \QuicConnectionContext{} instance using the
\ChannelOf{} class which provides efficient implemntation of a producer-consumer queue. A possible
runtime structure of the \QuicSocketContext{} servicing two separate connections is illustrated in
\autoref{fig:03-socket-context-architecture}. The two \QuicConnectionContext{} instances are
serviced by separate dedicated threads. The \QuicSocketContext{} itself does not need a dedicated
thread. Instead, its logic is performed in whatever thread that completes the pending asynchronous
socket receive call.

\begin{myFigure}{fig:03-socket-context-architecture}{Architecture of the server-side background
processing.}

\resizebox{0.88\textwidth}{!}{\input{img/03-server-socket-context.pdf_tex}}

\end{myFigure}

\subsubsection{Future Support for Connection Migration}

Lastly, we need to review the architecture for possible support for the connection migration
feature. Thanks to the \QuicConnectionContext{} not depending directly on a particular \Socket{}
instance, it is conceivable that a single \QuicConnectionContext{} instance could be part of two
\QuicSocketContext{}s --- one for the old endpoint address, and the other for the new one. Use of
the \ChannelOf{} class already ensures that the incoming datagrams could be queued concurrently by
both \QuicSocketContext{} instances. The only modification needed on the architecture level would be
allowing the \QuicConnectionContext{} to select the right \Socket{} from which the outgoing
datagrams should be sent.

\subsection{Public API Threading model}

The target API uses the \class{ValueTask} type designed for efficient asynchronous methods. The user
code will start an asynchronous operation that can be completed by a background thread servicing the
connection. This way, the user code cannot block the connection's background thread, which could
otherwise cause timeouts to be missed.

The \QuicConnection{} can be potentially used in a multithreaded environment, and the implementation
should allow concurrent usage of \QuicConnection{} and \QuicStream{} classes when it makes sense.
For example, it makes sense for an application code to process each \QuicStream{} in a different
thread. However, it does not make sense to concurrently write into one \QuicStream{} from two
threads without any synchronization. Our implementation will, therefore, provide the following
thread-safety guarantees for the API\@:

\begin{itemize}

  \item Individual streams can be used concurrently from different threads. However, each direction
of the stream (reading and writing) can be accessed only by one thread at a time and must be,
therefore, synchronized.

  \item Accepting/opening new streams on a connection can be done concurrently from multiple
threads.

  \item All other operations on \QuicConnection{} must be synchronized, including, e.g., starting
and aborting the connection.

\end{itemize}

\section{Packet Serialization/Deserialization}

Special care needs to be taken when implementing the QUIC packets' serialization to the wire format.
Inefficient data representation can negatively impact overall performance because processing
individual QUIC packets and frames will likely be a hot path in the implementation. The packet and
frame representation should be, therefore, carefully designed to avoid allocations.

The incoming packets can contain arbitrary encoding errors that should be handled efficiently and
gracefully, i.e., without throwing exceptions. Possible errors include values being outside of the
range of allowed values and incomplete or damaged packets.

\subsection{QUIC Packet and Frame Representation}\label{sec:03-data-representation}

Some packets contain many fields. Therefore, passing them around as individual variables would make
the implementation harder to maintain. The QUIC frames form coherent messages that should be
represented by individual \dotnet{} types. This gives us also an opportunity to keep the
serialization and deserialization code next to each other, making it easier to ensure that, e.g.,
the order of the serialized fields matches in both methods.

Representing QUIC frames as individual classes would introduce many heap allocations for every
received QUIC packet. This thesis, therefore, will model QUIC frames and the QUIC packet headers as
value types. Also, the payload of some frames may consist of large blocks of memory. Examples
include \STREAM{} and \CRYPTO{} frames, which can fill the entire payload of the QUIC packet.
Duplicating this block of memory into a separate \ArrayOf{\Byte{}} would be another unnecessary
memory allocation. \dotnet{} Core~2.1 introduced the \SpanOf{\class{T}} type, which can be used to
efficiently reference an arbitrary memory block, including memory allocated on the stack using the
\keyword{stackalloc} keyword.

The \SpanOf{\class{T}} type is a \keyword{ref struct} --- a special kind of value type which can be
stored only in local variables or inside another \keyword {ref struct}s. Limitations on the usage of
\keyword{ref struct}s also forbid their use as generic type arguments, e.g., for \FuncOf{} and other
generic delegates. This limitation is acceptable for \QuicConnection{} implementation because the
QUIC frames can be processed one after another right after being deserialized from the QUIC packet.

\subsection{QuicReader and QuicWriter}

Both serialization and deserialization require maintaining the current position in the buffer being
read from or written into. In order to simplify the serialization or deserialization code, our
implementation introduces \QuicReader{} and \QuicWriter{} classes as a primary means of reading and
writing QUIC primitives to memory. We have chosen \QuicReader{} and \QuicWriter{} to be reference
types allocated on the heap, because we expect to cache and reuse their instances multiple times.

The \QuicReader{} and \QuicWriter{} classes also take care of converting the endianity of the data
between big-endian used by QUIC and the endianity used by the machine running the application. This
is done by forwarding the calls to respective methods on the \class{BinaryPrimitives} class.

\section{Stream Implementation}\label{sec:03-stream-implementation}

QUIC is a transport protocol. Therefore, its entire purpose is transferring streams of data. Since
it is likely to be the hot path of the implementation, the internal handling of stream data must be
efficient and avoid unnecessary copying of stream data blocks.

QUIC recognizes four types of streams. These streams can have a sending part, receiving part, or
both. The fact which endpoint initiated the stream controls only which flow control limits apply to
that stream. Otherwise, all streams are handled equally. As mentioned in
\autoref{sec:02-stream-types}, bidirectional streams can be implemented as two unidirectional
streams. Therefore, our implementation will divide the QUIC stream logic into \SendStream{} and
\ReceiveStream{} classes, which will handle the sending and receiving behavior of the stream.

\subsection{Receiving Part of the Stream}\label{sec:03-receive-stream}

The implementation of \ReceiveStream{} must buffer the received data before delivering them to the
application in case the data are received out-of-order. It also needs to track the amount of
buffered memory and how much data was delivered to the application in order to correctly update flow
control limits for the peer.

\subsubsection{Stream Data Buffering}

Ideally, the stream implementation would be structured so that the stream data were copied straight
from the decrypted packet to the memory provided by the application. This would be possible if the
API used an event-based model with callbacks for incoming data. However, the current API design
utilizes a method-based model of the \Stream{} abstraction. If the application does not call the
\method{Read} or \method{ReadAsync} method, there is no application buffer to deliver into.

This implies that in order to avoid intermediate copies of the stream data, the buffer holding the
QUIC packet in which the stream data arrived cannot be reused to receive other QUIC packets until
the contained stream data are delivered to the application. Tracking which buffers can be reused for
receiving the following QUIC packets can be complicated because QUIC packets can carry multiple
\STREAM{} frames.

We believe that the implementational complexity of avoiding intermediate copies outweighs the
possible performance gain. For this reason, our pilot implementation will use a two-copy approach:
the first copy from the packet into an intermediate buffer, the second copy from the intermediate
buffer into the destination memory provided by the application.

\subsubsection{Packet Reordering and Data Deduplication}

Unreliability of the UDP protocol can cause QUIC packets to be reordered, lost or received multiple
times. Because of that, sections of the stream may be received multiple times, and the contents of
\STREAM{} frames can arbitrarily overlap.

Even though QUIC Flow control provides an upper limit on the data that needs to be buffered at any
given moment, allocating one large buffer may be a waste of memory, as the entire buffer might not
be needed at any given point in time. Therefore, our implementation uses a list of smaller buffers
and allocates only the necessary number of buffers needed to buffer currently received data. In
order to reduce the pressure on the \gls{gc}, buffers are reused using the \ArrayPoolOf{\Byte{}}
class to avoid their frequent allocation.

\subsubsection{Reading Data by the Application}

Because user code runs on a different thread from the internal \QuicConnection{} logic, access to
the buffered data must be synchronized. Our implementation uses the \ChannelOf{} class, which
provides an efficient implementation of the producer-consumer queue with support for asynchronous
operations using the \ValueTaskOf{} type.

Each time data are received for the stream, the implementation checks if there is a contiguous block
of memory that could be delivered to the application. For this reason, the implementation needs to
keep track of the parts of the stream which have already been received. If there is a new part of
the stream that can be delivered, a view into the relevant region of the buffer is queued into the
\ChannelOf{} using the \MemoryOf{\Byte{}} type.

\autoref{fig:03-receive-stream} illustrates how all parts of the \ReceiveStream{} work together. The
data from the incoming \STREAM{} frame are copied to the proper offset in intermediate buffers.
These buffers are rented from an \ArrayPoolOf{\Byte{}} as needed. A \MemoryOf{\Byte} instance
providing a view into the newly deliverable parts of the stream is then queued into a \ChannelOf{}
instance for delivery. The application thread retrieves the memory regions as needed and copies the
data into the buffer provided by the application. Once all data from an intermediate buffer are
delivered, the intermediate buffer is returned to the \ArrayPoolOf{\Byte{}} to be reused for
buffering future data.

\begin{myFigure}{fig:03-receive-stream}{Implementation of the receiving part of the stream}

  \resizebox{\linewidth}{!}{\input{img/03-receive-stream.pdf_tex}}

\end{myFigure}

\subsubsection{Flow Control Considerations}

QUIC Flow control limits for streams specify the maximum offset of data that an endpoint can send.
An endpoint needs to update these limits by sending a \MAXSTREAMDATA{} frame after the data are
delivered to the application. However, sending updated limits in every packet could waste space that
could be used by the application data. On the other hand, updating the limits too late could lead to
the other endpoint being blocked and signal the fact by sending a \STREAMDATABLOCKED{} frame. In
such a state, the endpoint cannot send more application data until it receives an update via
\MAXSTREAMDATA{} frame.

Our implementation will send a \MAXSTREAMDATA{} frame after the client uses more than half of the
limit provided since the last update. This should prevent updates being too often and at the same
time early enough that the sender does not run out of flow control credit.

\subsection{Sending Part of the Stream}\label{sec:03-send-stream}

The sending part of a QUIC stream must keep track of the state of all outbound data on the stream.
Any part of the stream can be lost during transmission, requiring its retransmission, and any
previously sent part of the stream can be acknowledged.

\subsubsection{Stream Data Buffering}

Similarly to buffering received data, we will analyze how many times the outgoing data need to be
copied before they are sent. The semantics of the \method{Write} and \method{WriteAsync} methods on
the \Stream{} class are such that when the method completes, the provided memory can be reused for
other purposes.

In this case, a single-copy implementation approach would require that the \method{Write} and
\method{WriteAsync} methods complete after the peer acknowledges that the data were received. In the
best-case scenario where no packet loss occurs, this approach implies a full round trip delay for
calling either of the methods. This approach would make the \QuicStream{} implementation
inconsistent with the other network-enabled \Stream{} implementations. Moreover, substantial time
would be spent waiting for acknowledgment instead of sending more data. In order to fully saturate
the available bandwidth, users of the \QuicStream{} would have to either provide very large data
buffers or use the \method{WriteAsync} method and overlap its execution by maintaining multiple
outstanding \ValueTask{}s.

In order to fit with the other \Stream{} implementations, our implementation will create an
intermediate copy of the data, and the \method{Write} and \method{WriteAsync} will complete once the
data is internally buffered. This allows applications to queue enough data to maximally utilize the
connection's bandwidth. In order to be memory-efficient, our implementation will, as in
\ReceiveStream{} implementation, use an \ArrayPoolOf{\Byte{}} to reduce the pressure on the garbage
collector.

\subsubsection{Acknowledgement, Loss, and Retransmission}

Each byte that in the stream can be, conceptually, in three different states:

\begin{itemize}

  \litem{Pending} The byte needs to be sent to the peer in some future \STREAM{} frame.

  \litem{In-flight} The byte has been sent, but it is uncertain if the containing packet was
received.

  \litem{Acknowledged} The packet containing the byte has been acknowledged by the other endpoint.
It is no longer necessary to buffer this byte.

\end{itemize}

The transitions between these states are straightforward. The data transition from pending to
in-flight by sending them in a \STREAM{} frame. When the packet is deemed lost by the loss detection
algorithm, the data transition back to the pending state, and if the packet has been acknowledged,
the data transition to the acknowledged state. The state of the individual parts of the stream can
be tracked by maintaining three sets of \textit{ranges} --- starts and ends of the blocks of data in
the same state.

\subsubsection{Writing Data by the Application}

Similarly to the receiving part of the stream, our implementation will use \ChannelOf{} type to
provide synchronization with the application thread. However, since \ChannelOf{} does not allow
random access to the provided data, a separate list of buffers is maintained for data which were
sent but not yet acknowledged.

\autoref{fig:03-send-stream} illustrates the process of writing data into the \SendStream{}. The
application data are copied to an internal buffer. This buffer comes from a shared
\ArrayPoolOf{\Byte{}}. When the buffer is full or the stream is flushed, the buffer is enqueued into
the synchronization \ChannelOf{}. The background thread then dequeues buffers from the \ChannelOf{}
and appends them into a list of retransmission buffers. Each time a new \STREAM{} frame for this
stream is written into a QUIC packet, the earliest not-yet-sent data are selected and copied into
the outgoing QUIC packet buffer. Once all data in a particular retransmission buffer are
acknowledged, the buffer can be removed and returned to the \ArrayPoolOf{\Byte}.

\begin{myFigure}{fig:03-send-stream}{Implementation of the sending part of the stream}

  \resizebox{\linewidth}{!}{\input{img/03-send-stream.pdf_tex}}

\end{myFigure}


\subsection{Abort/Dispose Model for Streams}

The \QuicStream{} class implements the \interface{IDisposable} interface which should provide
automatic closing of the stream when \QuicStream{} is used in a \keyword{using} statement or
\keyword{using} block. This implies resetting the writable part of the stream using \RESETSTREAM{}
frame and requesting reset of the readable part of the stream using \STOPSENDING{} frame.

However, both these frames require an application-level error code. QUIC specification does not
define any transport-level error codes for these stream operations and there is no way for the
application to specify which error code should be used. This has been established as a flaw of the
current API and is expected to change in future iterations design
iterations~\cite{dotnetGithubQuicShutdown}.

Leaving \interface{IDisposable} unimplemented would be misleading to users in the prototype
implementation. Therefore, our implementation will deviate from the specification by using error
code 0 regardless of its semantics in the application protocol, unless the stream has been
explicitly closed by the application using the \method{AbortRead} or \method{AbortWrite} methods.

Another area where the QUIC API is not well defined is the behavior of pending \method{AcceptStream}
calls on \QuicConnection{} when the connection is closed. In case the connection has been closed
gracefully from the application protocol's perspective, it would be better if this method did not
throw an exception but returned \keyword{null} reference. However, the semantics of the error codes
are defined by the application protocol and, therefore, the QUIC implementation cannot distinguish
between graceful or abortive connection close. Until the behavior of these methods in such scenarios
is better defined, our implementation will throw \class{QuicConnectionAborted\allowbreak{}Exception}
for all pending async calls.

\section{TLS Implementation}

TLS handshake forms an integral part of the QUIC connection establishment. Because correct TLS
implementation is crucial for ensuring the security of the resulting implementation, this thesis
should avoid implementing TLS algorithms by itself. Instead, it should reuse some existing and
well-tested implementation.

The novel way QUIC integrates with TLS requires a specific API from the TLS implementation. Below is
a non-exhaustive list of operations the TLS library's API must provide:

\begin{itemize}

  \item Querying the current state of the handshake

  \item Retrieving both the application-level secrets and secrets used to protect the handshake process

  \item Retrieving raw, unencrypted TLS messages to be sent to the other endpoint

  \item Obtaining the negotiated cipher

  \item Specifying protocols used for ALPN

  \item Specifying a custom TLS extension in order to exchange QUIC transport parameters

\end{itemize}

The \dotnet{} runtime libraries use different native libraries to provide TLS functionality on
different operating systems. On Windows, \dotnet{} uses \libname{Secure Channel}~\cite{Schannel}
(\libschannel{} for short) which is part of the Windows operating system. On Linux and macOS
systems, the \libopenssl{} library~\cite{OpenSSLWeb} is used.

\begin{description}

    \ditem{\libname{Secure Channel}} The \libschannel{} versions present in the latest Windows 10
builds support only TLS 1.2. However, future updates will also implement TLS 1.3. A preview of the
\libschannel{} with TLS 1.3 support can be obtained by installing an Insider build of Windows 10.
Because \libmsquic{} uses the \libschannel{} library when compiled for Windows, we assume that
\libschannel{} exposes the API necessary for a QUIC implementation.

    \ditem{\libopenssl{}} None of the mainstream versions of \libopenssl{} library expose necessary
API for integration into QUIC, and there are no plans to include such API in the next \libopenssl{}
3.0.0 release~\cite{OpensslBlogNoQuic}. However, developers at Akamai maintain a fork of
\libopenssl{} which adds the QUIC-enabling API~\cite{AkamaiOpensslGithub}. This modified version of
\libopenssl{} is used by \libmsquic{} (in Linux builds) and some other QUIC libraries like
\libname{quiche} from Cloudflare~\cite{quicheGithub}. Akamai's changes may be merged into
\libopenssl{} for the following 3.1.0 or later releases.

\end{description}

In conclusion, the APIs required for our QUIC implementation are currently only accessible in only
the preview versions of Windows 10. Relying solely on \libschannel{} for TLS 1.3 support would
severely impact cross-platform availability of our prototype implementation. We have, therefore,
decided to integrate with the modified \libopenssl{} library which supports all platforms supported
by \dotnet{}. This solution, however, has some drawbacks:

\begin{itemize}

  \item The modified \libopenssl{} binary must be present on the machine running our QUIC
implementation. This implies that the library must be compiled from source beforehand for the target
machine because pre-built binaries for the modified library are not available.

  \item Only a limited integration with X.509 certificates is possible because the
\class{X509Certificate} class implementation will be using a different binary --- \libname{CryptoAPI}
on Windows, unmodified \libopenssl{} on Linux and macOS\@.

\end{itemize}

These drawbacks are acceptable for the prototype implementation and will be eliminated once the
support for QUIC is released in mainstream versions of the \libopenssl{} and \libschannel{}
libraries.

\section{Packet Protection}\label{sec:03-packet-protection}

As described in \autoref{sec:02-packet-protection}, the packet encryption process consists of two
phases --- payload protection and header protection. The combined process requires the following
inputs:

\begin{itemize}

  \item Protection keys derived using the process explained in
\autoref{sec:02-encryption-key-derivation}

  \item Negotiated cipher

  \item Packet number (for encryption), or expected packet number (for decryption)

  \item The QUIC packet to encrypt or decrypt

\end{itemize}

The cipher is negotiated once and cannot be changed during the lifetime of the connection.
Protection keys can be changed only for the 1-RTT packets using the process of
\textit{\gls{key-update}} (see \autoref{sec:02-key-update}), which can be expected to be relatively
infrequent. The rest of the inputs change with every packet. Therefore, our implementation
encapsulates the packet protection implementation in a \class{CryptoSeal} class, which does not
depend on the rest of the QUIC implementation.

The process of receiving packets requires intermediate validation of the header fields. The
individual steps --- header protection and payload protection --- must be, therefore, exposed
separately. Also, the encryption and decryption should happen in-place to avoid unnecessary
allocations and copying of the packets.

Important consideration needs to be made for performing the actual encryption/decryption once all
inputs to \gls{aead} have been gathered. \dotnet{} does not contain an implementation of the CHACHA
family of ciphers. Fortunately, the \libopenssl{} library can be configured to not allow this cipher
to be negotiated and, therefore, our prototype can work without the CHACHA cipher support.

The other ciphers are based on the AES family of ciphers, which are supported by \dotnet{}. However,
individual classes implementing these ciphers do not share a common base class or interface.
Therefore, our implementation wraps the concrete AES implementations in classes derived from an
abstract \class{CryptoSealAlgorithm} class which defines a common interface required by
the \CryptoSeal{} class. \autoref{fig:03-crypto-seal} illustrates how the \CryptoSeal{} and
\class{CryptoSeal\allowbreak{}Algorithm} classes are connected.

\begin{myFigure}{fig:03-crypto-seal}{Relationship between CryptoSeal and CryptoSealAlgorithm
classes}

  \resizebox{\linewidth}{!}{\input{img/03-crypto-seal.pdf_tex}}

\end{myFigure}

The \gls{key-update} operation can be implemented by replacing the existing instance of the
\CryptoSeal{} class with a new one with the updated protection keys.

\section{Loss Detection and Loss Recovery}

In order to detect lost packets and retransmit any lost data, the \QuicConnection{} implementation
must keep track of which data have been sent in which packets and the timestamp when the packet was
sent for timeout detection. Our implementation will encapsulate this information and the loss
detection algorithm in a dedicated \RecoveryController{} class, which does not depend on the
\QuicConnection{} class. This way, its implementation can be unit tested separately from the rest of
the connection logic.

Another responsibility of the \RecoveryController{} class is maintaining the congestion window.
Although the specification defines only a single algorithm for congestion control based on TCP
NewReno~\autocite[Section~7]{draft-ietf-quic-recovery}, there are already experiments with other
algorithms like HyStart++ and CUBIC~\cite{cloudflareCubic}. The ability to support multiple such
algorithms could provide an opportunity for future experimentations. For that reason, the
\RecoveryController{} will use the \gls{strategy-pattern}~\cite{wiki:strategy-pattern} to allow
choosing the congestion control algorithm at runtime.

\section{Automated Testing}

It is necessary to implement automated tests that assert that the implementation conforms to the
QUIC protocol specification. Additionally, the public API makes use of existing concepts, namely the
\class{Stream} abstraction for QUIC streams, and thus it is necessary to ensure that \QuicStream{}
conforms to the expected \class{Stream} behavior.

The \dotnet{} runtime repository into which we wish to integrate our QUIC implementation uses the
\xUnit{} testing framework~\cite{xunit}. The \xUnit{} framework is one of the most mature testing
frameworks for \dotnet{}. Therefore, we will use it as well for writing tests for our
implementation.

\subsection{Testing the QUIC Protocol}

The QUIC protocol specification~\cite{draft-ietf-quic-transport} mostly specifies the protocol
behaviors in terms of what endpoint may or may not send in specific scenarios. This implies that the
correctness of the implementation as a whole can be tested by inspecting the contents of the
generated UDP datagrams.

Because our implementation separates QUIC connection logic from the socket IO, the unit tests can be
written against the internal \QuicConnection{} API\@. Exchange of QUIC packets between two
\QuicConnection{} instances can be simulated by the unit testing code. An advantage of using this
internal API is that the unit tests can be written as single-threaded --- without any background
processing thread which could introduce nondeterminism to the tests. However, in order to make the
tests completely deterministic, the testing code also must be able to control the exact timing of
events. In order to achieve that, our \QuicConnection{} implementation will not obtain the current
timestamp directly, e.g., by calling the \class{Stopwatch}\texttt{.\method{GetTimestamp}()} method.
Instead, the current timestamp will be provided as an argument to the internal API, either by the
\QuicConnectionContext{} class or the unit testing code.

\subsubsection{Testing Harness}\label{sec:03-testing-harness}

A large number of the unit tests will consist of inspecting contents of the QUIC packets exchanged
between a pair of \QuicConnection{} instances or checking the internal state of a connection after a
particular packet exchange. However, correctly implemented \QuicConnection{} will never violate the
protocol and, therefore, additional logic is required to test error conditions.

In order to test the behavior in erroneous conditions, the testing code needs to be able to either
compose an invalid packet to be sent or intercept a valid packet and modify it to elicit an error
response. However, doing this is difficult for two reasons:

\begin{itemize}

  \item All packets are encrypted and protected against modification. The testing code must first
retrieve the correct \CryptoSeal{} from the sender \QuicConnection{} instance and unprotect the
packet. After any modification, the encryption must be reapplied.

  \item Modifying a particular QUIC frame may change the encoded frame's size because of the
variable-length encoding used (see \autoref{sec:02-variable-length-encoding}). If the size of the
QUIC frame changes, then all following frames must be shifted, and the Length field in the packet
header must be adjusted.

\end{itemize}

Doing this manually would make unit testing function very verbose and make the code hard to
understand. To keep the tests concise and easy to understand, we will implement a \textit{testing
harness} which will provide the functionality mentioned above via a set of helper methods. These
methods will provide a declarative way to specify what the QUIC packet or datagram must contain and
register callbacks for packet modification to elicit error responses.

\autoref{lst:03-desired-unit-test} shows an example how the desired testing harness would be used in
conjunction with \xUnit{} testing framework to test if the first UDP datagram sent by the client has
the correct size\footnote{As part of the prevention against traffic amplification attacks, all UDP
datagrams containing an Initial packet sent by client must be larger than 1200
bytes~\autocite[Section~8.1]{draft-ietf-quic-transport}.}. The \method{GetDatagramToSend} method
will get the next UDP datagram to be sent and present it in a structured manner. Later, the
\method{ShouldHaveFrame<\class{TFrame}>} method will internally check if a frame represented by
\class{TFrame} type is present in the packet and will invoke the provided callback for further
assertions.

\begin{myListingCsharpNoPageBreak}{lst:03-desired-unit-test}{Example unit test inspecting QUIC packets contents.}{Fact, Assert, Datagram, InitialPacket, CryptoFrame, QuicConstants}{}
    [Fact]
    public void |ClientInitialDatagramHasInitialPacketWithCryptoFrame|()
    {
        var datagram = |GetDatagramToSend|(Client);
        Assert.|Equal|(
            QuicConstants.MinimumClientInitialDatagramSize,
            datagram.Size);

        var initial = Assert.|IsType|<InitialPacket>(
            Assert.|Single|(datagram.Packets));
        Assert.|Equal|(0, initial.PacketNumber);

        initial.|ShouldHaveFrame|<CryptoFrame>(crypto =>
        {
            Assert.|Equal|(0, crypto.Offset);
            Assert.|NotEmpty|(crypto.CryptoData);
        });
    }
\end{myListingCsharpNoPageBreak}

However, in \autoref{sec:03-data-representation}, we mentioned that the types representing QUIC
frames must be \keyword{ref struct}s in order to contain \SpanOf{\class{T}} instances, and that
\keyword{ref struct}s cannot be used as type arguments for generic classes or methods. This could be
overcome by using \MemoryOf{\class{T}} instead of \SpanOf{\class{T}}, but the use of structs for
frame types still poses a problem for modification of the QUIC frames by a callback.

Structs are normally passed by value; therefore, the callback would either have to accept the frame
by reference using the \keyword{ref} keyword or return the modified frame as the return value.
Neither of these solutions is perfect. Parameters types with \keyword{ref} modifiers cannot be
inferred for lambdas and need to be stated explicitly, and having to return the modified frame is
unintuitive for non-modifying callbacks. Furthermore, \keyword{struct}s representing QUIC frames
should be preferably made \keyword{readonly} to allow more compiler optimizations, so modifying the
frame would require creating a new instance of the frame, overwriting the old value.

For the above reasons, we decided to duplicate the types for frame representation using mutable
\keyword{class}es. This duplication will allow expressing the unit tests in a succinct, natural
declarative manner. This is illustrated in the test in \autoref{lst:03-unit-test-intercept} which
uses an \method{Intercept1RttFrame<\class{TFrame}>} method to intercept a frame of type
\class{TFrame} and modify it. In this case, we shift the range of acknowledged packets by one and,
by doing so, simulate acknowledging a packet that the server has not sent yet, which is a violation
of the protocol which should result in connection termination via a \CONNECTIONCLOSE{} frame.

\begin{myListingCsharpNoPageBreak}{lst:03-unit-test-intercept}{Example unit test simulating error conditions.}{Fact, Assert, InitialPacket, QuicError, ConnectionCloseFrame, AckFrame}{TransportErrorCode, FrameType}
    [Fact]
    public void |ConnectionCloseWhenAckingFuturePacket|()
    {
        // ... setup ommited for brevity

        |Intercept1RttFrame|<AckFrame>(Client, Server, ack =>
        {
            // ack one packet more than originally intended
            ack.LargestAcknowledged++;
        });

        |Get1RttToSend|(Server).|ShouldHaveFrame|<ConnectionCloseFrame>(f =>
        {
            Assert.|Equal|(TransportErrorCode.ProtocolViolation, f.ErrorCode);
            Assert.|Equal|(FrameType.Ack, f.ErrorFrameType);
            Assert.|Equal|(QuicError.InvalidAckRange, f.ReasonPhrase);
        });
    }
\end{myListingCsharpNoPageBreak}

\subsection{Testing the Public API}

The \dotnet{} runtime repository already contains a suite of tests used to test the
\libmsquic{}-based QUIC implementation. Another large separate suite of tests exists for ensuring
that all \class{Stream} implementation behave consistently. All these tests can be used to ensure
that our implementation of \QuicListener{}, \QuicConnection{}, and \QuicStream{} classes conforms to
the public API specification.

\section{Diagnostics}

QUIC is a very complex protocol, and bug investigation can be complicated. By stopping the
application on a breakpoint, the developer inadvertently changes the application's behavior.
Extended pauses for inspecting internal connection state make the implementation miss important
timeouts, possibly leading to connection being closed due to the idle timeout. This makes
interactive debugging of issues which span multiple roundtrips almost impossible.

There are two possible ways of gaining better insight into the connection's behavior --- externally
inspecting the packets sent over the network and producing verbose logs (called traces) by the
implementation.

\subsection{Inspecting the Sent Packets}

Inspecting the connection behavior by observing the packets sent over the network is a non-invasive
way of diagnosing issues in any networking protocol. An example of a tool that can be used for this
task is Wireshark~\cite{web:wireshark}, which also supports QUIC\@.

However, since QUIC is always encrypted, inspecting QUIC packets via Wireshark is not as
straightforward as for other network protocols like plain TCP\@. Implementations must leak the
encryption secrets used in the connection, e.g., into a log file, and these keys must be provided to
Wireshark to enable decrypting the packets. While this approach is relatively simple to implement,
it provides little concrete information about the internal state of the connection. The developer
must infer the internal connection state from his knowledge about the implementation and the
contents of the captured QUIC packets.

\subsection{Producing Traces from the Connection}

A better insight into the connection's behavior can be gained by emitting verbose logs that can be
later analyzed. However, such traces can consist of thousand lines of output which may be difficult
to read and reason about. Fortunately, there are tools which can visualize individual events in the
connection in a graphic format that is easier to understand. One of these tool is
\textit{qvis}~\cite{web:qvis} which is part of the \textit{quiclog} suite~\cite{githubquiclog}. The
\textit{qvis} tool consumes logs in a JSON format, which can be produced directly or generated from
implementation-specific log format or even from generated from network traffic captured using
Wireshark, provided that the encryption secrets used in the connection are known.

For our implementation, we have decided to keep the tracing implementation simple by directly
emitting traces in the JSON format which can be directly consumed by the \textit{qvis} tool without
any further conversions. However, serializing packet information into JSON format is likely to incur
noticeable overhead and, therfore, a better logging format should be considered for future
development.

\section{Integration into .NET Runtime Codebase}

The work of this thesis aims to be eventually mergeable into the \dotnet{} runtime codebase. As of
writing this thesis, there already exists a project for \textit{System.Net.Quic.dll}. Therefore, the
source code may be placed there directly without additional changes.

There is, however, the issue with the \libopenssl{} dependency. The \libopenssl{} library is written
in C and must be compiled for the target machine. Integrating the \libopenssl{} compilation into the
\dotnet{} runtime build process would impose additional requirements for the \dotnet{} runtime
compilation. The additional dependencies would complicate building the \dotnet{} runtime as part of
the continuous integration process. For this reason, it is necessary to provide also a mock TLS
implementation that would be used in automated tests as part of continuous integration.

The mock TLS implementation can also be used locally to avoid building the \libopenssl{} dependency.
The only limitation of the mock TLS implementation is that it does not allow interop with other QUIC
implementations.

\chapter{Evaluation}\label{chap:04-evaluation}

This chapter evaluates the performance of the managed QUIC implementation developed in this thesis
and compares it to the existing \libmsquic{}-based implementation. We will also compare both managed
and \libmsquic{}-based QUIC implementations to \class{SslStream} which uses a combination of TCP and
TLS\@. In our evaluation, we will focus on two performance characteristics relevant for network
communication:

\begin{itemize}

        \litem{Throughput} The rate at which the application data are sent. High total throughput is
essential for servers when handling a large number of parallel connections.

        \litem{Round Trip Time Latency} The delay between sending a request and receiving a
response. This metric is important mainly for clients, where lower latencies mean faster responses
from the server. Rather than measuring the average latency, it is more common to measure, e.g., the
99th percentile of latency, which indicates minimum expected latencies in the slowest 1\% replies.
This metric is vital in the context of webpages and communication with browsers, especially for
websites for which web browsers generate a large number (even hundreds) of HTTP requests for a
single page access. In such cases, the probability of at least one of the requests being slower than
the 99th percentile raises significantly and, therefore, a high value of the 99th percentile would
negatively affect the total page load time for a significant number of users~\cite{Treat2015}.

\end{itemize}

\section{Evaluation Environment}

We have evaluated our implementation in two different environments. Our primary evaluation
environment consists of two blade servers on the same LAN network. The relevant software and
hardware parameters of the servers are:

\begin{itemize}

        \litem{CPU} Intel\textsuperscript{\textregistered{}} Xeon\textsuperscript{\textregistered{}}
E3-1270 v6 @ \SI{3.80}{\giga\hertz} (\SI{4.20}{\giga\hertz} turbo boost, 4 cores, 2 threads per
core)

        \litem{RAM} \SI{32}{\giga\byte} (\SI{2400}{\mega\hertz})

        \litem{Network Interface Card} Intel\textsuperscript{\textregistered{}} Ethernet Controller
I350 (\SI[per-mode=symbol]{1}{\giga\bit\per\second})

        \litem{OS} Fedora 32

\end{itemize}

Unfortunately, the two servers are connected only by \SI[per-mode=symbol]{1}{\giga\bit\per\second}
Ethernet, which implies a theoretical upper limit for TCP's throughput at
\SI[per-mode=symbol]{117.5}{\mega\byte\per\second} and similar limit for UDP\@. In some
measurements, we reached values that approached this limit. Therefore, when appropriate, we also
evaluated our implementation using a loopback network interface on a desktop workstation PC with the
following specifications:

\begin{itemize}

        \litem{CPU} Intel\textsuperscript{\textregistered{}} Core\textsuperscript{\textregistered{}}
i7-8700 @ \SI{3.20}{\giga\hertz} (\SI{4.60}{\giga\hertz} turbo boost, 6 cores, 2 threads per core)

        \litem{RAM} \SI{32}{\giga\byte} (\SI{2400}{\mega\hertz})

        \litem{OS} Windows 10 version 20H2\footnote{Update to Windows 10 20H2 includes significant
TCP and UDP performance improvements~\cite{theregister20h2}. Measurements taken on older versions of
Windows would yield significantly different results.}

\end{itemize}

In both environments, the testing application was run in a a 64-bit process on preview
\dotnet{}~6.0.0 host (CoreCLR 6.0.20.55508, CoreFX 6.0.20.255508).

\subsection{Drawbacks of Measuring Loopback Performance}\label{sec:04-localhost-mtu}

Measuring over the loopback interface removes any upper limit on bandwidth because the loopback
interface is implemented purely in software inside the operating system. It also minimizes the
transport latency and guarantees zero packet loss. The software-based implementation of the loopback
interface allows optimizations that are not possible on real networks. As an example, on Windows
operating system, loopback connections can send \SI{64}{\kibi\byte} IP packets --- which is the
maximum size allowed by the IP protocol --- without any link-layer fragmentation. On the other hand,
Ethernet connections can carry only up to \SI{1500}{\byte} IP datagrams without fragmentation. The
greater size of datagrams on the loopback interface reduces overhead when TCP is used for
inter-process communication on the same machine.

Both managed and \libmsquic{}-based QUIC implementations send IP datagrams which are at most
\SI{1500}{\byte} even on loopback, because of the hardcoded limits in the codebase. In order to make
the loopback comparison more fair, we leveraged the ability to limit the maximum outbound IP packet
size to \SI{1500}{\byte} on Windows sockets API\@. From \csharp{}, this can be achieved using
following statement:

\lstsetcsharpsettings{}
\begin{lstlisting}[numbers=none,enums={SocketOptionLevel,SocketOptionName}]
socket.|SetSocketOption|(SocketOptionLevel.IP,
  (SocketOptionName)76 /* IP_USER_MTU */, 1500 /* size */);
\end{lstlisting}

However, the \texttt{IP_USER_MTU} option is only available on Windows 10 version 20H2 or newer. On
other platforms, the \method{SetSocketOption} call above would throw an exception. In order to keep
cross-platform compatibility, the attached source code does not contain this statement. However, it
was used to obtain the results presented in this chapter.

\subsection{Running with MsQuic Support}

\libmsquic{}-based QUIC implementation does not work in the default installation of \dotnet{}
because \dotnet{} does not distribute the \libmsquic{} binary. Programs that want to use the
\libmsquic{}-based QUIC implementation should reference the
\texttt{System.Net\allowbreak{}.Experimental.MsQuic} NuGet
package~\cite{SystemNetQuicExperimentalMsquic} which provides the necessary \libmsquic{} binary.
However, the Windows \libmsquic{} binary distributed in the NuGet package requires a Windows Insider
Preview build of Windows because it depends on newer \libschannel{} version for TLS implementation,
which is not available in the mainstream Windows version.

Unfortunately, we cannot install a preview build of Windows on the workstation we will use for
running the tests. Fortunately, in order to support Linux, \libmsquic{} can also use the same
modified \libopenssl{} version for the TLS implementation as our managed QUIC implementation. Even
though building \libmsquic{} with \libopenssl{} for Windows is not officially supported, we were
able to produce such a build with small changes in the \libmsquic{} build configuration. Therefore,
the \libmsquic{} library used in experiments in this chapter has been compiled locally using source
code from commit \texttt{dc2a6cf0dd12e27} from the official \libmsquic{}
repository~\cite{msquicGithub} and configured to use the same \libopenssl{} library as the managed
QUIC implementation. It should be noted that our custom Windows build of \libmsquic{} may show
slightly different performance characteristics than the official build for Windows, which uses the
\libschannel{} library.

On Linux, we could use the above mentioned
\texttt{System.Net\allowbreak{}.Experimental.\allowbreak{}MsQuic} NuGet package, but that would
complicate the build process because the package would have to be referenced only on Linux builds.
Therefore, we have decided to compile the Linux build of \libmsquic{} locally as well. Like the
Windows build, the Linux build of \libmsquic{} used in our tests was compiled from commit
\texttt{dc2a6cf0dd12e27} but did not require changes to build configuration, as it uses the modified
\libopenssl{} library by default.

All measurements of the \libmsquic{}-based QUIC support for \dotnet{} include any overhead
introduced by the interop layer bridging the native library API to \dotnet{} QUIC API\@. The
measured results, therefore, do not represent raw \libmsquic{} library performance, but the
performance observed by \dotnet{} applications which use the \libmsquic{}-based implementation
provider which can be validly compared to the performance of our managed QUIC implementation.

\subsection{Benchmarking Application}\label{sec:04-benchmark-app}

The actual throughput and latency measurements are done using a dedicated benchmarking \dotnet{}
application whose source code can be found in the
\filename{src/supplementary/benchmark/ThroughputTests} directory in the attachments of this thesis.
The application implements a trivial echo server and clients which exchange messages of the
specified size. The first parameter to the application selects one of the following modes:

\begin{description}

  \ditem{\texttt{server}} Starts only the server-side part of the application.

  \ditem{\texttt{client}} Starts only the client-side part of the application. In this mode, the
application spawns multiple clients and reports the measurement results on the standard output.

  \ditem{\texttt{inproc}} Starts both the server and client in the same process. All network
communication is done over the loopback network interface.

\end{description}

The benchmarking application accepts multiple parameters. A list of the most important follows.

\begin{description}

    \ditem{\texttt{-e, --endpoint}} The endpoint on which to listen (server) or to which to connect
(client). Not applicable for \texttt{inproc} mode.

    \ditem{\texttt{--certificate-file}} Path to the X.509 certificate file.

    \ditem{\texttt{--key-file}} Path to the X.509 certificate private key file.

    \ditem{\texttt{-t, --tcp}} Use TCP instead of QUIC\@.

    \ditem{\texttt{-c, --connections}} The number of connetions to create.

    \ditem{\texttt{-s, --streams}} The number of streams to create in each connection. Applicable
only when QUIC is used.

    \ditem{\texttt{-m, --message-size}} The size of sent messages in bytes.

    \ditem{\texttt{-w, --warmup-time}} Time before starting to take measurements.

    \ditem{\texttt{-d, --test-duration}} Time after which the measurement should stop.

    \ditem{\texttt{-i, --reporting-interval}} Delay between reports of intermediate measurements.

    \ditem{\texttt{-n --no-wait}} Whether clients wait for a reply before sending another message.

\end{description}

The full list of parameters can be found in the \filename{Program.cs} file, or by running the
program with the \texttt{--help} parameter. By default, the implementation uses managed QUIC
implementation. Switching to \libmsquic{}-based implementation can be achieved by defining the
\texttt{DOTNETQUIC_PROVIDER} environment variable to \texttt{msquic}.

The client mode of the application switches between two behaviors. By default, clients do not wait
for a reply from the server before sending another message. This lets us measure the total
throughput of the system. When using the \texttt{-n} switch, clients wait until a response from the
server is received and measure the latency as the delay between sending the message and receiving
the reply.

\section{Measurement Results}\label{sec:04-perf-results}

In the following subsections, we present the results of the individual performance experiments. The
throughput and latencies were measured as follows:

\begin{itemize}

  \litem{Throughput} Throughput of the entire server. This is the total amount of application data
echoed back to clients across all connections.

  \litem{Latency} The delay between client writing the message to the \Stream{} and reading back the
full reply. In our tests, we will measure the 99th percentile of latency. The latency measurements
will be taken with the use of \texttt{-n} flag in the benchmark application.

\end{itemize}

We will perform three kinds of performance comparisons:

\begin{itemize}

        \litem{Multiple Parallel Streams Performance} These tests compare how the managed and
\libmsquic{}-based QUIC implementations scale with increasing server load. These tests will utilize
the stream multiplexing of QUIC to send messages using multiple parallel QUIC streams.

        \litem{Single Stream Performance} These tests use only a single stream in a connection.
These tests should allow us to compare the performance between QUIC implementations and
TCP+TLS-based \SslStream{}.

        \litem{Performance in Simulated Cellular Network} In these tests, we will try to approximate
the characteristics of a real-world cellular network by increasing lag and packet loss in the LAN
environment.

\end{itemize}

In all experiments in this section, client and server parts of the benchmarking application were run
in separate processes. Each test case had a \SI{5}{\second} warm-up time and collected data for
another \SI{15}{\second}. These intervals were long enough to produce stable measurements across
multiple test runs.

In the benchmarking application, clients and server exchange messages of the size specified by the
\texttt{-m} parameter. In our tests, we will use mainly two message sizes: \SI{256}{\byte} and
\SI{4096}{\byte}. The \SI{256}{\byte} message are small enough to fit into a single QUIC packet,
while the larger \SI{4096}{\byte} messages are guaranteed to be split across multiple QUIC packets.
We expect that increasing the message size should increase the latency because more packets need to
be sent for both the message and the reply. Also, increasing the message should increase throughput
because there are fewer calls to the \method{Write} and \method{WriteAsync} methods for the same
amount of data.

Running the benchmarking application for tests in this section has been automated using a
PowerShell~\cite{powershell} script. The script can be found at
\filename{src/supplementary/benchmark/ThroughputTests/run.ps1} in this thesis' attachments. Comments
at the top of the file explain the usage of the script.

\subsection{Multiple Parallel Streams Performance}\label{sec:04-multi-stream-perf}

The first set of tests compared the managed and \libmsquic{}-based QUIC implementations. These tests
were run with increasingly larger messages and with a greater number of parallel connections and
streams to see how the two implementations scale.

\autoref{fig:04-lab-multi-stream-throughput} shows the measured throughput.
\autoref{fig:04-lab-multi-stream-throughput-a} shows the base-line performance when \SI{256}{\byte}
messages are sent using a single stream per connection. The other figures show measurements after
increasing message size to \SI{4096}{\byte} (\autoref{fig:04-lab-multi-stream-throughput-b}),
increasing number of streams to 32 (\autoref{fig:04-lab-multi-stream-throughput-c}), or both
(\autoref{fig:04-lab-multi-stream-throughput-d}). An immediate observation can be made that the
managed implementation produced very similar pattern in all four test runs, with the highest
throughput at 4 parallel connections and then decreasing. On the other hand, the \libmsquic{}-based
implementation maintains almost the same throughput regardless of the number of connections, but
increases throughput both when increasing the message size and number of streams in a connection.

\begin{myFigure}{fig:04-lab-multi-stream-throughput}{Multiple stream QUIC throughput measurements (LAN)}
\begin{mySubfigure}{0.49\linewidth}{fig:04-lab-multi-stream-throughput-a}{1 stream, \SI{256}{\byte} messages}
\footnotesize
\input{plots/04-lab-multi-stream-throughput-a.tex}
\end{mySubfigure}
\begin{mySubfigure}{0.49\linewidth}{fig:04-lab-multi-stream-throughput-b}{1 stream, \SI{4096}{\byte} messages}
\footnotesize
\input{plots/04-lab-multi-stream-throughput-b.tex}
\end{mySubfigure}

\begin{mySubfigure}{0.49\linewidth}{fig:04-lab-multi-stream-throughput-c}{32 stream, \SI{256}{\byte} messages}
\footnotesize
\input{plots/04-lab-multi-stream-throughput-c.tex}
\end{mySubfigure}
\begin{mySubfigure}{0.49\linewidth}{fig:04-lab-multi-stream-throughput-d}{32 stream, \SI{4096}{\byte} messages}
\footnotesize
\input{plots/04-lab-multi-stream-throughput-d.tex}
\end{mySubfigure}
\end{myFigure}

In case with \SI{4096}{\byte} messages, the \libmsquic{}-based implementation consistently saturates
the \SI{1}{\giga\bit} network connection between the two blade servers. To give more perspective on
the relative performance between the two implementations, \autoref{fig:04-multi-stream-throughput}
shows results of the two test runs on the Windows workstation over the loopback interface. Depending
on the number of connections, the throughput of the \libmsquic{}-based implementation was up to four
times greater than our implementation. Results of runs with \SI{256}{\byte} messages were left out
for brevity because they did not differ from the ones in the LAN environment.

\begin{myFigure}{fig:04-multi-stream-throughput}{Multiple stream QUIC throughput measurements (Loopback)}
\begin{mySubfigure}{0.49\linewidth}{fig:04-multi-stream-throughput-b}{1 stream, \SI{4096}{\byte} messages}
\footnotesize
\input{plots/04-multi-stream-throughput-b.tex}
\end{mySubfigure}
\begin{mySubfigure}{0.49\linewidth}{fig:04-multi-stream-throughput-d}{32 stream, \SI{4096}{\byte} messages}
\footnotesize
\input{plots/04-multi-stream-throughput-d.tex}
\end{mySubfigure}
\end{myFigure}

The peak of the throughput of our managed QUIC implementation at four connections can be explained
by the fact that we are spawning a long-running background task for each connection. Four
connections lead to a number of parallel \class{Task}s that were enough to completely utilize the
entire CPU\@.

As for the decline of the throughput of managed QUIC implementation when 16 or more connections were
used, profiling showed that a substantial amount of time was spent in the garbage collection. At 256
connections, up to 50\% of the total CPU time was spent in GC, divided into relatively large pauses
of \SI{20}{\milli\second} or more. The pauses introduced by GC led to QUIC packets being considered
lost, which led to the collapse of the congestion window in the QUIC connection, further reducing
the rate at which data were sent. This points to the fact that even though we carefully avoided
needless allocations in our implementation, there are still enough sources of allocation left to
degrade a server's performance under heavy load.

An example of a large source of allocations in the managed QUIC implementation is pooling the
buffers for sending or receiving QUIC packets. Even though the implementation uses
\ArrayPoolOf{\Byte}\texttt{.Shared} to reuse allocated buffers, this does not work well with a large
number of connections. This is because \ArrayPoolOf{\Byte}\texttt{.Shared} retains only a small
number of \ArrayOf{\Byte} instances for pooling and lets GC collect the rest\footnote{In
\dotnet{}~5, the implementation of \ArrayPoolOf{\Byte}\texttt{.Shared} maintains a separate pool for
each CPU core. Each per-CPU core pool organizes pooled arrays into 17 buckets with array sizes
ranging from \SI{16}{\byte} to \SI{1}{\mebi\byte}. Each bucket retains up to 8 array instances of
similar size.}. With a large number of connections, the buffers are often rented and returned in
large bursts, leading to a lot of large arrays being discarded upon return to the
\ArrayPoolOf{\Byte}\texttt{.Shared} and subsequently allocated anew.

Another source of allocations are \Socket{}\texttt{.\method{SendTo}} and
\Socket{}\texttt{.\method{ReceiveFrom}} methods which allocate of a new \class{EndPoint} instance
for each call. In order to remove this particular source of allocations, an allocation-free API for
sending or receiving UDP datagrams needs to be designed and implemented on the \Socket{} class.

\libmsquic{}-based implementation shows similar throughput regardless of the number of parallel
connections. Instead, it dramatically increases with message size and only slightly with the number
of streams in the connections. When we tried sending messages larger than \SI{4096}{\byte}, the
total throughput did not increase significantly anymore. Upon closer inspection of the \libmsquic{}
network traffic using Wireshark, it turns out that in tests with \SI{256}{\byte} messages, many of
the sent packets were very small, possibly containing only \ACK{} frames. Thus, the total available
network bandwidth was mostly unutilized. When \SI{4096}{\byte} messages are used, more space in QUIC
packets is utilized, which increases the throughput. Our investigation, however, did not uncover why
the total throughput of \libmsquic{} does not increase with the number of connections in the tests
with \SI{256}{\byte} messages.

\autoref{fig:04-lab-multi-stream-latency} shows the latencies measured in the LAN environment for
the same four test cases as above. In all tests, the \libmsquic{}-based implementation exhibits a
lower 99th percentile of latency, especially in test cases with a high number of connections. The
same measurement on Windows loopback did not yield significantly different results.

\begin{myFigure}{fig:04-lab-multi-stream-latency}{Multiple stream QUIC latency measurements (LAN)}
\begin{mySubfigure}{0.49\linewidth}{fig:04-lab-multi-stream-latency-a}{1 stream, \SI{256}{\byte} messages}
\footnotesize
\input{plots/04-lab-multi-stream-latency-a.tex}
\end{mySubfigure}
\begin{mySubfigure}{0.49\linewidth}{fig:04-lab-multi-stream-latency-b}{1 stream, \SI{4096}{\byte} messages}
\footnotesize
\input{plots/04-lab-multi-stream-latency-b.tex}
\end{mySubfigure}

\begin{mySubfigure}{0.49\linewidth}{fig:04-lab-multi-stream-latency-c}{32 stream, \SI{256}{\byte} messages}
\footnotesize
\input{plots/04-lab-multi-stream-latency-c.tex}
\end{mySubfigure}
\begin{mySubfigure}{0.49\linewidth}{fig:04-lab-multi-stream-latency-d}{32 stream, \SI{4096}{\byte} messages}
\footnotesize
\input{plots/04-lab-multi-stream-latency-d.tex}
\end{mySubfigure}
\end{myFigure}

When measuring latency, the benchmarking application waits for the server to respond before sending
another message. However, with many parallel connections and streams, the total network traffic
becomes similar to that produced when measuring throughput. This was the case, especially when we
see significant differences between the latencies of the two implementations. Therefore, the change
in latency can be attributed to the frequent and long pauses introduced by the GC that we have
described above. On the other hand, in the test runs where the managed implementation is close to
that of \libmsquic{}, the GC activity was under 5\% and did not cause long frequent pauses.

By increasing the message size, we increased the probability that the application message will be
affected by the packet loss introduced by the GC pauses. This happens because the message is spread
over multiple QUIC packets, and the loss of any of those packets will delay the entire message until
the missing part is retransmitted. This explains why increasing the message size drastically affects
the latency with a large number of parallel connections.

Increasing the number of streams also increases network traffic, making our implementation
susceptible to the increased GC pauses. However, after a certain threshold, increasing the number of
streams does not increase the throughput of the implementation anymore. Instead, it leads to an
increase in latency because it still increases the number of outstanding parallel requests.

Lastly, we would like to remind the reader that we measured the 99th percentile of the latency,
which considers only the slowest 1\% of requests. To give more perspective on the latency
distribution on the two implementations, \autoref{fig:04-lab-multi-stream-latency-median} shows the
median measurements of the latency. For brevity, we include only the measurements for the cases with
\SI{4096}{\byte} messages. The median measurements of latencies of our implementation are similar or
lower than that of the \libmsquic{}-based one.

\begin{myFigure}{fig:04-lab-multi-stream-latency-median}{Multiple stream QUIC latency median measurements (LAN)}
\begin{mySubfigure}{0.49\linewidth}{fig:04-lab-multi-stream-latency-median-a}{1 stream, \SI{4096}{\byte} messages}
\footnotesize
\input{plots/04-lab-multi-stream-latency-median-a.tex}
\end{mySubfigure}
\begin{mySubfigure}{0.49\linewidth}{fig:04-lab-multi-stream-latency-median-b}{32 stream, \SI{4096}{\byte} messages}
\footnotesize
\input{plots/04-lab-multi-stream-latency-median-b.tex}
\end{mySubfigure}
\end{myFigure}

We believe that this difference in the latency distributions between the two implementations is due
to their different architectures. \libmsquic{} spawns several worker threads (depending on the
number of CPU cores) which process events for individual connections from an event queue. The event
queue ensures that the connections are serviced evenly and at regular intervals. When the load on
the server increases, the extra latency is spread evenly across all active connections. This is
further supported by the fact that the median values for the latency of \libmsquic{} are only
slightly lower than its 99th percentile from \autoref{fig:04-lab-multi-stream-latency}.

Our implementation, on the other hand, relies solely on the \dotnet{} \class{Task} scheduler, which
schedules all background tasks independently and without any priority information. This allows for
greater variance in the latency of our QUIC implementation. In the future, our implementation's
architecture should be improved to reduce the scheduling variance and, therefore, the 99th
percentile of the latency.

In summary, our QUIC implementation outperforms the \libmsquic{}-based one in scenarios where small
messages are exchanged between client and server using a small number of streams. However, our
implementation does not scale as well as the \libmsquic-based one. A large number of parallel
connections and streams increases pressure on the GC, which introduces frequent and long pauses in
the application. Lastly, our implementation's latency distribution is much less uniform than that of
the \libmsquic{}-based one because of the way background work is scheduled on the \dotnet{}
thread-pool.

\subsection{Single Stream Performance}

In the second test, we compared the single-stream performance of managed QUIC implementations to the
performance of the combination of TCP and TLS available via the \TcpClient{} and \SslStream{}
\dotnet{} classes. We will also include the measurements for \libmsquic{} for completeness.

\autoref{fig:04-lab-single-stream-throughput} shows the results of the throughput measurements taken
between the two blade servers connected by \SI{1}{\giga\bit} Ethernet. In both \SI{256}{\byte}
messages (\autoref{fig:04-lab-single-stream-throughput-a}) and \SI{4096}{\byte} messages
(\autoref{fig:04-lab-single-stream-throughput-b}) scenario, TCP implementation outperforms both QUIC
implementations. Also, out of the three implementations, our QUIC implementation is the only one
which is not bottlenecked by the \SI{1}{\giga\bit} LAN\@.

\begin{myFigure}{fig:04-lab-single-stream-throughput}{Single stream throughput measurements (LAN)}
\begin{mySubfigure}{0.49\linewidth}{fig:04-lab-single-stream-throughput-a}{\SI{256}{\byte} messages}
\footnotesize
\input{plots/04-lab-single-stream-throughput-a.tex}
\end{mySubfigure}
\begin{mySubfigure}{0.49\linewidth}{fig:04-lab-single-stream-throughput-b}{\SI{4096}{\byte} messages}
\footnotesize
\input{plots/04-lab-single-stream-throughput-b.tex}
\end{mySubfigure}
\end{myFigure}

Because the throughput of some implementations in the previous figure was limited by the
\SI{1}{\giga\bit} network, \autoref{fig:04-single-stream-throughput} shows the results measured over
the loopback interface on the Windows workstation. We can see clearly that TCP still provides
greater throughput than the \libmsquic{}-based implementation\footnote{Without limiting the loopback
IP packet size by setting the \texttt{IP_USER_MTU} option to 1500 (see
\autoref{sec:04-localhost-mtu}), the TCP throughput with \SI{4096}{\byte} messages would reach up to
\SI[per-mode=symbol]{900}{\mebi\byte\per\second}.}

\begin{myFigure}{fig:04-single-stream-throughput}{Single stream throughput measurements (Loopback)}
\begin{mySubfigure}{0.49\linewidth}{fig:04-single-stream-throughput-a}{\SI{256}{\byte} messages}
\footnotesize
\input{plots/04-single-stream-throughput-a.tex}
\end{mySubfigure}
\begin{mySubfigure}{0.49\linewidth}{fig:04-single-stream-throughput-b}{\SI{4096}{\byte} messages}
\footnotesize
\input{plots/04-single-stream-throughput-b.tex}
\end{mySubfigure}
\end{myFigure}

By comparing the numbers in \autoref{fig:04-single-stream-throughput-b} to that from
\autoref{fig:04-multi-stream-throughput-d}, we see that the performance of TCP is also greater than
that of \libmsquic{}-based implementation when using 32 parallel streams per connection streams.

\autoref{fig:04-lab-single-stream-latency} shows the latency measurement results in the LAN
environment. These results show that TCP exhibits substantially lower latencies than either QUIC
implementation. Very similar results were also obtained on the Windows workstation.

\begin{myFigure}{fig:04-lab-single-stream-latency}{Single stream latency measurements (LAN)}
\begin{mySubfigure}{0.49\linewidth}{fig:04-lab-single-stream-latency-a}{\SI{256}{\byte} messages}
\footnotesize
\input{plots/04-lab-single-stream-latency-a.tex}
\end{mySubfigure}
\begin{mySubfigure}{0.49\linewidth}{fig:04-lab-single-stream-latency-b}{\SI{4096}{\byte} messages}
\footnotesize
\input{plots/04-lab-single-stream-latency-b.tex}
\end{mySubfigure}
\end{myFigure}

It is worth noticing that in \autoref{fig:04-lab-single-stream-latency-a}, the latency of managed
QUIC implementation is very low until it reaches a certain threshold and then increases rapidly as a
consequence of increased GC pauses.

The results presented in this section suggest that neither QUIC implementation can compete with the
\SslStream{} in neither local LAN networks nor over the loopback interface (even with the MTU
reduction we explained in \autoref{sec:04-localhost-mtu}). It is possible, however unlikely, that
the performance over real \SI[per-mode=symbol]{10}{\giga\bit\per\second} network will yield
different results. A possible explanation for these results is that the TCP networking stack is far
better optimized than UDP\@. An example of such optimization is \textit{TCP offload
engine}~\cite{wiki:tcp-offload-engine} which offloads the TCP/IP stack implementation onto the
controller on the network interface card.

\subsection{Performance in Simulated Cellular Network}

In the last test, we compared the resilience of QUIC and TCP+TLS implementations in a simulated
network with delay and packet loss. When choosing the values of delay and packet loss parameters, we
tried to approximate a 4G cellular broadband network's behavior. We performed a trivial connection
speed test of local 4G network using a smartphone and the
\href{https://speedtest.net}{\texttt{speedtest.net}} testing web application. This test reported
latency of \SI{25}{\milli\second}, which we then chose as the latency in the simulated network in
our tests.

As for the packet loss in 4G network, in 2012, \citeauthor{measuring4G} have measured
characteristics of 4G networks of major US cellular network carriers and measured packet loss
percentages from 0.004\% to 0.1\%~\cite{measuring4G}. Even though the cellular network technology
may have evolved since then, we could not find more recent measurements. Therefore, we will use
those two packet loss values in our test parameters.

It should be noted that our approximation of the cellular network is not perfect. The latency in a
real-world network is not constant but has a random distribution, and the \SI{25}{\milli\second} we
measured earlier was only the mean value. Similarly, packet loss does not affect each packet
independently. Instead, packets are often lost in bursts. There are multiple models for modeling
more realistic packet loss, such as the Gilbert-Elliot model~\cite{wiki:burst-error}. However, we
were unable to find enough data to compute appropriate parameters for such models.

To simulate the 4G network, we used the \textit{Traffic Control} capabilities of the Linux kernel.
The necessary configuration can be set using the \texttt{tc} utility. For example, the following
command was issued on both machines to configure \SI{25}{\milli\second} lag and 0.004\% packet loss:

\begin{myVerbatim}
tc qdisc change dev eth0 root netem delay 12ms loss 0.004%
\end{myVerbatim}

Note that the set delay is half of the desired latency because the traffic shaping done by
\texttt{tc} applies only to outbound traffic. Each endpoint, therefore, adds half of the total
latency.

\autoref{fig:04-lab-loss-throughput} shows the measured throughput for 0.004\% packet loss
(\autoref{fig:04-lab-loss-throughput-a}) and 0.1\% packet loss
(\autoref{fig:04-lab-loss-throughput-b}). As expected, The increase in packet loss leads to a
decrease in throughput. In the presence of only low packet loss, the TCP implementation exhibits an
order of magnitude greater throughput. However, the throughput of the TCP protocol decreased at a
greater rate than that of the two QUIC implementation. The throughput of the two QUIC implementation
is comparable in all tests.

\begin{myFigure}{fig:04-lab-loss-throughput}{Single stream throughput measurements in simulated network}
\begin{mySubfigure}{0.49\linewidth}{fig:04-lab-loss-throughput-a}{\SI{25}{\milli\second} latency, 0.004\% loss}
\footnotesize
\input{plots/04-lab-loss-throughput-a.tex}
\end{mySubfigure}
\begin{mySubfigure}{0.49\linewidth}{fig:04-lab-loss-throughput-b}{\SI{25}{\milli\second} latency, 0.1\% loss}
\footnotesize
\input{plots/04-lab-loss-throughput-b.tex}
\end{mySubfigure}
\end{myFigure}

Upon closer inspection, it turns out that the managed QUIC implementation is bottlenecked by the
stream buffering implementation. In \autoref{sec:03-send-stream}, we explained the design behind the
sending part of the stream, namely the \SendStream{} class. The implementation uses up to eight
\SI{16}{\kibi\byte} buffers to store the application. Once these buffers are filled, subsequent
calls to the \method{Write} or \method{WriteAsync} method will block until some data are
acknowledged. The application data must be buffered until the other endpoint acknowledges their
reception. This takes at least one roundtrip, i.e., \SI{25}{\milli\second} in our test. Our
implementation can, therefore, send at most $8 \cdot 16 = 128 \si{\kibi\byte}$ on a single stream during
each round trip. This sending rate implies a theoretical maximum throughput as
$128\si{\kibi\byte} \div 0.025 = 5 \si[per-mode=symbol]{\mebi\byte\per\second}$.

This hypothesis can be tested by doubling the delay on the simulated network.
\autoref{fig:04-lab-loss-throughput-long} shows measurements with \SI{50}{\milli\second} latencies
(\autoref{fig:04-lab-loss-throughput-a-50}) and \SI{100}{\milli\second} latencies
(\autoref{fig:04-lab-loss-throughput-a-100}). In these variations of the test we see that the
throughput of the managed implementation halves everytime the latency is doubled.

\begin{myFigure}{fig:04-lab-loss-throughput-long}{Single stream throughput measurements in simulated network with long delays}
\begin{mySubfigure}{0.49\linewidth}{fig:04-lab-loss-throughput-a-50}{\SI{50}{\milli\second} latency, 0.004\% loss}
\footnotesize
\input{plots/04-lab-loss-throughput-a-50.tex}
\end{mySubfigure}
\begin{mySubfigure}{0.49\linewidth}{fig:04-lab-loss-throughput-a-100}{\SI{100}{\milli\second} latency, 0.004\% loss}
\footnotesize
\input{plots/04-lab-loss-throughput-a-100.tex}
\end{mySubfigure}
\end{myFigure}

This performance limitation could be fixed simply by increasing the number of buffers used to buffer
application data. However, a large number of such buffers would lead to inefficient memory
utilization in low-latency networks. Instead, the amount of data buffered should be changed
dynamically based on information such as the congestion window's size or the flow control limit
advertised by the other endpoint. We leave the implementation of such dynamic buffering for future
work.

\autoref{fig:04-lab-loss-throughput} shows the latency measurements results. Each implementation
exhibits the same latency for both 0.004\% (\autoref{fig:04-lab-loss-latency-a}) and 0.1\%
(\autoref{fig:04-lab-loss-latency-b}) packet loss. This means that 0.1\% is still low enough to not
affect the 99th percentile of latency.

\begin{myFigure}{fig:04-lab-loss-latency}{Single stream latency measurements in simulated network}
\begin{mySubfigure}{0.49\linewidth}{fig:04-lab-loss-latency-a}{\SI{25}{\milli\second} latency, 0.004\% loss}
\footnotesize
\input{plots/04-lab-loss-latency-a.tex}
\end{mySubfigure}
\begin{mySubfigure}{0.49\linewidth}{fig:04-lab-loss-latency-b}{\SI{25}{\milli\second} latency, 0.1\% loss}
\footnotesize
\input{plots/04-lab-loss-latency-b.tex}
\end{mySubfigure}
\end{myFigure}

All implementations show almost the same latency in almost all tests. The single exception is our
implementation, which exhibits greater latencies for \SI{4096}{\byte} messages. This is because this
message size does not fit a single packet, our pacer implementation spaces out individual packets
one by one. The \libmsquic{}-based implementation, on the other hand, sends small batches of packets
for each tick of the pacer and the small batches are enough to transmit the entire \SI{4096}{\byte}
message.

By using only a single stream in each QUIC connection in these tests, it could be argued that we did
not utilize one of the big advantages of QUIC over TCP\@ --- the reduced \gls{head-of-line-blocking}.
\autoref{fig:04-lab-multi-loss-throughput} shows the measurements of throughput similar to the one
in \autoref{fig:04-lab-loss-throughput}, but with QUIC connections transmitting data in 32 parallel
streams. Contrary to the expectation, the throughput of the QUIC implementation increases only
slightly.

\begin{myFigure}{fig:04-lab-multi-loss-throughput}{Multiple stream throughput measurements in simulated network}
\begin{mySubfigure}{0.49\linewidth}{fig:04-lab-multi-loss-throughput-a}{32 streams, \SI{25}{\milli\second} latency, 0.004\% loss}
\footnotesize
\input{plots/04-lab-multi-loss-throughput-a.tex}
\end{mySubfigure}
\begin{mySubfigure}{0.49\linewidth}{fig:04-lab-multi-loss-throughput-b}{32 streams, \SI{25}{\milli\second} latency, 0.1\% loss}
\footnotesize
\input{plots/04-lab-multi-loss-throughput-b.tex}
\end{mySubfigure}
\end{myFigure}

Because we were sending data using multiple streams, the managed QUIC implementation was no longer
blocked by the limited buffering discussed above. Instead, it was limited by the size of the
congestion window. When investigated more closely, it turns out that the throughput of both QUIC
implementation was very high at the beginning, possibly even higher than that of TCP\@. However,
during the \SI{5}{\second} warmup period, The packet loss led to a significant reduction of the
congestion window, and the throughput was reduced to the values we can see on the plot.

We believe that this behavior was caused by the too uniform distribution of the packet loss we used
in our testing environment. In an environment with uniform packet loss, almost all lost packets
induce a separate congestion event and reducing the congestion window. When packets are lost in
bursts --- as is more common in real-world networks --- the loss of multiple consecutive packets is
handled as a single congestion event, leading to less frequent reductions of the congestion window.

In this case, the protocol's behavior depends on the implementation of the congestion control
algorithm. The managed QUIC implementation uses an adaptation of the NewReno algorithm described by
the QUIC specification~\autocite[Section~7]{draft-ietf-quic-recovery}. The \libmsquic{}-based QUIC
implementation and TCP used the CUBIC algorithm~\cite{rfc8312}. Because of the substantial
difference in performance, we consider it worthwhile to experiment with different congestion control
algorithms for our implementation. Studying the behavior and implementation of TCP congestion
control in the Linux kernel should also provide useful insight for improving our implementation's
behavior. However, due to the complexity, we leave such experiments and investigations for future
work.

\section{Interop between QUIC Implementations}

The benchmarking application can also test the managed QUIC implementation's interoperability with
the \libmsquic{}-based implementation. This can be done by starting one application in
\texttt{server} mode using one implementation provider and another instance of the application in
\texttt{client} mode using the other provider.

More exhaustive interop tests could be achieved using the open-source QUIC Interop Test
Runner~\cite{QuicInteropRunner}. However, this level of testing is outside of the scope of this
thesis and is subject to future work.

\section{Summary}

Compared to \libmsquic{}-based QUIC implementation provider, the managed QUIC implementation
developed as part of this thesis can provide higher throughput when using a small number of parallel
connections while maintaining comparable latencies. However, in the current state, the managed QUIC
implementation does not scale very well, and both throughput and observed latencies degrade when a
large number of connections are created. When under substantial load, the \libmsquic{}-based
implementation delivered throughput up to four times greater than our implementation.

However, in our tests, the TCP+TLS stack has significantly larger throughput and lower latencies
than either QUIC implementation in our LAN environment. A possible explanation of this is the
presence of hardware optimizations for the TCP stack on the network interface card.

We also compared the QUIC and TCP implementations in a simulated 4G network. In our tests, we again
measured significantly higher throughput using TCP+TLS stack than either QUIC implementation, even
if we used multiple streams in QUIC connections. As for latencies, all implementations exhibit
almost identical latency, except for our implementation, which had larger latencies once the
messages exceeded the maximum size of a QUIC packet.

From our measurements, we conclude that our prototype QUIC implementation will need much work in
order to provide performance similar to that of \libmsquic{}. We have identified some of the
underlying problems in the current managed implementation and proposed future improvements. However,
both QUIC implementations will require a significant amount of optimizations before reaching
performance comparable to present TCP implementations.

\chapter{Developer Documentation}

The Managed QUIC implementation developed in this thesis is contained in the
\filename{System.Net.Quic} project in a fork of the official \dotnet{} runtime repository. The
source code is attached in \todo{path to the src in attachments.}

The source code for the \filename{System.Net.Quic} project is located inside the
\filename{src/dotnet-runtime/src/libraries/System.Net.Quic/} directory. The directory listing in
\autoref{lst:05-system-net-quic-structure} shows the structure of the project directory. The figure
also emphasises the \filename{Managed} and \filename{UnitTests} directories which contain the code
developed as part of this thesis. These directories are the main focus of this chapter.

\newcommand{\mydtcomment}[1]{\DTcomment{#1}}

% use smaller width to make the tree more compact
\DTsetlength{0.2em}{0.6em}{0.2em}{0.4pt}{1.6pt}
\renewcommand{\DTstylecomment}[1]{{\footnotesize\rmfamily #1}}
\renewcommand{\DTstyle}[1]{{\footnotesize\texttt{#1}}}
\definecolor{dtemphcolor}{rgb}{0,0,0.75}
\newcommand{\dtemph}[1]{\textcolor{dtemphcolor}{\emph{#1}}}

\begin{myListing}[Directory structure of the System.Net.Quic project.]{lst:05-system-net-quic-structure}{Directory structure of the System.Net.Quic project. Emphasised items contain the implementation developed in this thesis.}
\dirtree{%
  .1 {src/dotnet-runtime/src/libraries/System.Net.Quic/}.
  .2 {ref}\mydtcomment{Refererence assembly code}.
  .2 {src}\mydtcomment{Main library source code}.
  .3 {Resources}.
  .4 {Strings.resx}\mydtcomment{Definition of localizable strings like exception messages.}.
  .3 {System}.
  .4 {Net}.
  .5 {Quic}.
  .6 {Implementations}\mydtcomment{Root directory for all QUIC implementations}.
  .7 {\dtemph{Managed}}\mydtcomment{\dtemph{This thesis' implementation sources}}.
  .7 {MsQuic}\mydtcomment{\libmsquic{}-based implementation sources}.
  .7 {Mock}\mydtcomment{Mock implementation used only in tests}.
  .6 {Interop}\mydtcomment{Imports from native libraries}.
  .3 {System.Net.Quic.csproj}.
  .2 {tests}\mydtcomment{Library tests source code}.
  .3 {certs}\mydtcomment{X.509 certificates used in tests}.
  .4 {cert.crt}\mydtcomment{Public certificate file}.
  .4 {cert.key}\mydtcomment{Private key file}.
  .3 {FunctionalTests}\mydtcomment{Tests against the public API}.
  .4 {System.Net.Quic.Functional.Tests.csproj}.
  .3 {\dtemph{UnitTests}}\mydtcomment{\dtemph{Managed implementation unit tests}}.
  .4 {\dtemph{System.Net.Quic.Unit.Tests.csproj}}.
  .2 {Directory.Build.props}.
  .2 {System.Net.Quic.sln}.
}
\end{myListing}

\todo{sections for: development (environment setup, setup for debugging)? running tests?}

\section{QUIC Implementation Providers}

The \filename{System.Net.Quic} project contains multiple implementations of the QUIC protocol. This
is not achieved via polymorphism of the QUIC API classes, but rather using indirection to
\textit{implementation providers}. Each \QuicListener{}, \QuicConnection{}, and \QuicStream{}
instance contains a reference to a \QuicListenerProvider{}, \QuicConnectionProvider{}, or
\QuicStreamProvider{} instance, respectively. Implementations of each provider are provided by each
QUIC implementation in the \filename{System.Net.Quic} library. Figure
\autoref{fig:05-quic-impl-providers} illustrates this indirection layer using a class diagram.

\begin{myFigure}{fig:05-quic-impl-providers}{Implementation providers for the QUIC API classes}

\resizebox{0.95\linewidth}{!}{\input{img/05-quic-implementation-providers.pdf_tex}}

\end{myFigure}

Creation of new \QuicListener{} and \QuicConnection{} instances is implemented using the
\gls{abstract-factory}~\cite{wiki:abstract-factory-pattern}. There are multiple implementations of
\QuicImplementationProvider{} class, one for each QUIC implementation. Instances of
\QuicImplementationProvider{} class are available as static properties of the
\class{QuicImplementationProviders} static class. There are five such providers:

\begin{itemize}

  \litem[]{\texttt{Managed}} Managed implementation with TLS backed by \libopenssl{} fork with QUIC enabling API.

  \litem[]{\texttt{ManagedMockTls}} Managed implementation with mock TLS, which does not depend on external TLS implementation, but cannot interoperate with other QUIC implementations.

  \litem[]{\texttt{MsQuic}} QUIC implementation backed by \libmsquic{} native library.

  \litem[]{\texttt{Mock}} A mock QUIC implementation for use in tests.

  \litem[]{\texttt{Default}} The default implementation. Same as \texttt{Managed}, but can be redirected to other implementations by setting \texttt{DOTNETQUIC_PROVIDER} environment variable to the desired provider name.

\end{itemize}

The \QuicListener{} and \QuicConnection{} classes have a constructor overload which accepts an
instance of the \QuicImplementationProvider{} to be used. This way, the QUIC implementation can be
selected during runtime. This also allows reusing a suite of functional tests for all
implementations by simply changing the implementation provider.

\section{Managed QUIC Implementation}

The source code for the managed implementation developed in this thesis is located under the
\filename{System.Net.Quic/src/System/Net/Quic/Implementations/Managed/} subdirectory.
\autoref{lst:05-managed-quic-structure} outlines the directory structure of the implementation.

\begin{myListing}{lst:05-managed-quic-structure}{Directory structure of the managed QUIC implementation}
\dirtree{%
  .1 {System.Net.Quic/src/System/Net/Quic/Implementations/Managed}.
  .2 {Internal}\mydtcomment{Internal code of the implementation}.
  .3 {Crypto}\mydtcomment{Cryptographic facilities}.
  .3 {Frames}\mydtcomment{Definition of QUIC frames}.
  .3 {Headers}\mydtcomment{Definition of QUIC packet headers}.
  .3 {Packets}\mydtcomment{QUIC packet number spaces handling}.
  .3 {Parsing}\mydtcomment{Parsing of QUIC primitives}.
  .3 {Recovery}\mydtcomment{Loss detection and recovery}.
  .3 {Sockets}\mydtcomment{Servicing socket IO}.
  .3 {Streams}\mydtcomment{Stream buffering}.
  .3 {Tls}\mydtcomment{TLS integration}.
  .4 {Mock}\mydtcomment{Mock TLS implementation}.
  .4 {OpenSsl}\mydtcomment{OpenSSL TLS integration}.
  .3 {Tracing}\mydtcomment{Tracing and logging facilities}.
  .2 {ManagedQuicConnection.cs}\mydtcomment{Implementation of public API}.
  .2 {ManagedQuicConnection.Frames.cs}\mydtcomment{Processing of QUIC frames}.
  .2 {ManagedQuicConnection.Packets.cs}\mydtcomment{Processing of QUIC packets}.
  .2 {ManagedQuicConnection.Recovery.cs}\mydtcomment{Handling packet loss}.
  .2 {ManagedQuicConnection.Stream.cs}\mydtcomment{Stream management}.
  .2 {ManagedQuicImplementationProvider.cs}\mydtcomment{Abstract factory}.
  .2 {ManagedQuicListener.cs}\mydtcomment{Implementation of public API}.
  .2 {ManagedQuicStream.cs}\mydtcomment{Implementation of public API}.
}
\end{myListing}

The implementation is exposed using the \ManagedQuicListener{}, \ManagedQuicConnection{},
\ManagedQuicStream{} implementation provider classes and the \ManagedQuicImplementationProvider{}
factory. The source code for these classes can be found in the root directory of the implementation.

\subsection{Supporting Data Structures}

The QUIC implementation requires a few specialized data structures.

\todo{should this be in analysis or documentation?}

\subsubsection{RangeSet}

Efficient representation of acknowledged ranges of packet numbers and maintaining information about
which parts of the QUIC stream have been sent or acknowledged require a data structure capable of
efficiently performing set operations on ranges of integers. The \RangeSet{} class represents such a
set.

It is expected that the number of ranges in one \RangeSet{} instance will be rather small. Therefore
our implementation uses a sorted \class{List<T>} instance to store individual ranges. The ranges
themselves are represented by \class{RangeSet.\allowbreak{}Range} struct which contains the first and last
element represented by the range.

\subsubsection{PacketNumberWindow}

The \PacketNumberWindow{} class is used to test if QUIC packet the given packet number has already
been received. Internally, it contains two 64-bit fields: a bitmask marking the received bits and an
integer representing the packet number represented by the least significant bit of the mask. Marking
a packet number as received consists of setting a particular bit in the mask, shifting the other
bits if necessary. \todo{is this clear or does it need an illustration?}

The data structure can maintain state only for 64 consecutive packet numbers. This leads to packets
with packet number under the window being defensively marked as already received and, therefore,
discarded. The tollerance of the loss detection algorithm outlined in
\autoref{sec:02-loss-detection} allow way smaller packet reordering. In fact, the specification
suggests maximum reordering tollerance of 3, meaning that receiving a packet number should mark all
packets as lost except the previous two. Therefore, it is very likely that any newly received packet
with packet number under the window maintained by the \PacketNumberWindow{} instance is already
deemed lost and its data retransmitted in a different packet. Thus the width of 64 of the packet
number window is sufficient.

\todo{possibly include segmented ArrayBuffer, if I decide to implement it myself}

\subsection{Implementation Overview}



\section{Tests Implementation}

The QUIC implementation is covered by an extensive suite of unit tests and functional tests. The
unit tests focus on the correctness of the individual parts of the managed QUIC implementation, and
that the QUIC packets sent by the implementation conform to the QUIC protocol specification.
Functional tests, on the other hand, tests high-level functionality like being able to send and
receive data, and that the public API behavior conforms to the public API specification.

\subsection{Unit Tests}

\subsection{Functional Tests}

\chapter{User Documentation}\label{chap:06-user-docs}

This section provides guidance on how to obtain the build of the \dotnet{} runtime with managed QUIC
implementation, how to install it, and how to use the code to develop other applications.

\section{Getting Started}

This thesis provides a branch of the \dotnet{} runtime codebase with managed QUIC implementation.
Since our branch contains changes only in the \SystemNetQuicDll{}, the easiest way of composing a
fully working \dotnet{} distribution is obtaining a full SDK installation of the latest master
development version of \dotnet{}~6 and replacing the \SystemNetQuicDll{}. This section explains how
to do this without affecting other \dotnet{} SDK installations present on the machine.

The managed implementation depends on a particular \libopenssl{} version to interoperate with other
QUIC implementations. This section also explains how to deploy a locally built \libopenssl{} to be
automatically used by user code.

The complete setup process consists of the following steps, which we explained in greater detail in
the subsequent subsections:

\begin{enumerate}

  \item Build the \SystemNetQuicDll{} library from our branch of the \dotnet{} runtime sources.

  \item Compose a new local \dotnet{}~6 SDK installation with locally built \SystemNetQuicDll{}.

  \item \textit{(optional)} Compile a QUIC-supporting \libopenssl{} from source and deploy it.

  \item Configure the development environment to use the new \dotnet{} installation when compiling
and running user applications.

\end{enumerate}

\subsection{Building the System.Net.Quic Library from Source}\label{sec:06-build-runtime}

The source code for the \dotnet{} runtime with managed QUIC implementation is part of this thesis'
attachments in the \filename{src/dotnet-runtime/} directory. The latest version of the source code
can also be found on the thesis author's GitHub~\cite{githubRzikmRuntimelab}. For the remainder of
this section, all paths will be relative to the \dotnet{} runtime repository directory.

The \dotnet{} runtime repository contains a descriptive guide on how to build the sources. The
necessary prerequisites are listed in files inside the \filename{docs/workflow/requirements/}
directory, separately for each operating system. Once all necessary prerequisities are installed,
the entire \dotnet{} runtime can be built using the \filename{build.cmd} batch file (on Windows) or
\filename{build.sh} script (on Linux). The arguments are the same for both operating systems.

\begin{myVerbatim}
> ./build.cmd -subset clr+libs -configuration release
\end{myVerbatim}

The above command will build the Common Language Runtime (CLR) and all libraries in Release
configuration. The artifacts are available in the \filename{artifacts/bin/System.Net.Quic}
directory. The important artifacts from this directory are:

\begin{itemize}

  \item \filename{ref/net6.0-Release/System.Net.Quic.dll}: The so-called reference
assembly~\cite{ReferenceAssemblyDocs} that specifies the public API of the library.

  \item \filename{net6.0-{OS}-Release/System.Net.Quic.dll}: Where \verb|{OS}| is the identifier for
the operating system running on the machine. This is the \dotnet{} assembly with the actual QUIC
implementation.

\end{itemize}

\subsection{Creating a Local Installation of \dotnet{}}\label{sec:06localdotnet}

Now we need to download the latest \dotnet{}~6 SDK\@. A zip archive containing the SDK can be
downloaded from a link listed in the official SDK installer GitHub
repository~\cite{dotnetSdkGithub}. Download the ``Master (6.0.x Runtime)'' build for your platform
and extract it to a convenient location. In the remainder of this guide, the directory containing
the extracted contents will be referred to as \filename{DOTNET_ROOT}.

The \SystemNetQuicDll{} produced in the previous subsection must be copied to appropriate locations
in the \filename{DOTNET_ROOT}. The reference \SystemNetQuicDll{} assembly should be copied over the
existing one in the
\filename{DOTNET_ROOT/packs/Microsoft.NETCore.App.Ref/6.0.0-{version}/ref/net6.0/} directory, and
the implementation assembly to should be copied to the
\filename{DOTNET_ROOT/shared/Microsoft.NETCore.App/6.0.0-{version}/} directory, overwriting the
existing files.

\subsection{Adding the OpenSSL Library}\label{sec:06openssl}

The implementation requires a QUIC-supporting \libopenssl{} library build from a development branch
maintained by Akamai. The appropriate source codes can be found in the \filename{extern/openssl}
directory in the thesis attachments. The source code is also available online on Akamai's
GitHub~\cite{AkamaiOpensslGithub}. The implementation has been developed and tested with the
\texttt{OpenSSL_1_1_1g-quic} branch of the code, but other QUIC-enabled branches of \libopenssl{}
version 1.1.1 should work as well.

Before building the \libopenssl{} library form source, check the \filename{NOTES.{OS}} file in the
repository and make sure all prerequisites are installed on the machine. After that, the
\libopenssl{} library can be built by running the following command inside the repository. Note that
for Windows OS, you must run these commands using the \textit{Developer command prompt for Visual
  Studio} in order to have the necessary tools in \texttt{PATH}.

\begin{myVerbatim}
# Windows
> perl Configure VC-WIN64A
> nmake

# Linux
> ./config
> make
\end{myVerbatim}

This will produce the \libname{libcrypto} and \libname{libssl} libraries in the \libopenssl{}
repository root. On windows, these libraries are named \texttt{libcrypto-1_1-x64.dll} and
\texttt{libssl-1_1-x64.dll}. These libraries are loaded by the managed QUIC implementation during
runtime and, therefore, must be present in a location where the OS loader can find them. This can be
achieved by putting the libraries in any of the following locations:

\begin{itemize}

  \item Next to the compiled program executable.

  \item A directory listed in the \texttt{PATH} environment variable

  \item \textit{(preferred)} next to the \SystemNetQuicDll{} library in the \dotnet{} installation
directory, i.e., \filename{DOTNET_ROOT/shared/Microsoft.NETCore.App/6.0.0-{version}/}.

\end{itemize}

Note that if there is already a different version of \libopenssl{} installed on the system, it is
necessary to ensure that the system loads the correct \libopenssl{} version. This is different for
each operating system:

\begin{itemize}

  \item On Windows, this can be ensured by placing the DLL files in the same directory as the
program executable. The entire library search process, including the order of directories searched,
is described in Windows documentation~\cite{windowsDllSearch} in detail.

  \item On Linux, this can be achieved by defining the \texttt{LD_LIBRARY_PATH} environment variable
to the directory containing the \libopenssl{} libraries. Additional information about loading of
dynamic libraries on Linux can be found in the manual pages for \texttt{ld.so}. These manual pages
are also available online~\cite{linuxDllSearch}.

\end{itemize}

\subsection{Configuring the Development Environment}\label{sec:06-env-vars}

Lastly, we need to configure the environment variables so that the \dotnet{} SDK installation
created in the previous step is used when building the user code. For this, the following
environment variables need to be defined correctly.

\begin{description}

        \ditem{\texttt{DOTNET_ROOT}} Path to the \dotnet{} installation directory. This instructs
the build process to use the SDK installed in this directory. Use the path to the local \dotnet{}~6
SDK installation created in \autoref{sec:06localdotnet}.

        \ditem{\texttt{DOTNET_MULTILEVEL_LOOKUP}} Set this to ``0''. This instructs the build
process not to look for SDK installation in other places than \filename{DOTNET_ROOT}.

        \ditem{\texttt{PATH}} Prepend the \filename{DOTNET_ROOT} directory to the beginning of the
\texttt{PATH} variable to make sure the \filename{dotnet} executable from the \filename{DOTNET_ROOT}
is used over the system-wide installed one.
\end{description}

After configuring the variables, check the output of the \verb|dotnet --info| command. Assuming
\texttt{DOTNET_ROOT} is \filename{C:\dotnet\\}, then the output should be similar to the
\autoref{lst:06-dotnet-info-output}. The list of installed \dotnet{} runtimes should contain
\texttt{Microsoft.NETCore.App} from the local \dotnet{}~6 SDK installation prepared in
\autoref{sec:06-local-dotnet} (check the path inside the brackets). Note that the listing contains
version numbers of the latest \dotnet{}~6 SDK at the time of writing this text, and the SDK
installer would be updated since then to a newer version.

\begin{myListingVerbatim}[Output of the \texttt{dotnet --info} command in configured environment]{lst:06-dotnet-info-output}{Output of the \texttt{dotnet --info} command in correctly configured environment. The unimportant portions of the output in grey has been left out brevity}
> dotnet --info
&color{colorunimportant}.NET SDK (reflecting any global.json):
&color{colorunimportant} ...

&color{colorunimportant}Runtime Environment:
&color{colorunimportant} ...

&color{colorunimportant}Host (useful for support):
&color{colorunimportant} ...

.NET SDKs installed:
  6.0.100-alpha.1.20563.2 [&textbf{C:\dotnet\sdk}]

.NET runtimes installed:
  Microsoft.AspNetCore.App 6.0.0-alpha.1.20526.6 [&textbf{C:\dotnet\shared\}&textcolor{colorunimportant}{...}]
  Microsoft.NETCore.App 6.0.0-alpha.1.20560.10 [&textbf{C:\dotnet\shared\}&textcolor{colorunimportant}{...}]
  Microsoft.WindowsDesktop.App 6.0.0-alpha.1.20560.7 [&textbf{C:\dotnet\shared\}&textcolor{colorunimportant}{...}]

&color{colorunimportant}To install additional .NET runtimes or SDKs:
&color{colorunimportant}  https://aka.ms/dotnet-download
\end{myListingVerbatim}

With the environment configured as described, \dotnet{} applications can be compiled against the
\dotnet{}~6 SDK either using the \verb|dotnet build| command-line command, or using Visual Studio or
any other IDE\@.

\subsection{Creating a Sample \dotnet{}~6 Project}

The last step is creating a new project and configuring it to use \dotnet{}~6. This section
demonstrates how this can be done using the \verb|dotnet| command-line tool. Assuming the
environment variables have been configured as specified in \autoref{sec:06-env-vars}, you can create
a new project using the following commands:

\begin{myVerbatim}
> mkdir hello-net6
> cd hello-net6
> dotnet new console
\end{myVerbatim}

These commands will create a \filename{hello-net6/hello-net6.csproj} file. Use of \dotnet{}~6
preview requires minor changes to the project file. Namely changing the \xmltag{TargetFramework}
property to \verb|net6.0|. The modified project file contents are listed in
\autoref{lst:06-csproj-net6.0}.

\begin{myListingXml}{lst:06-csproj-net6.0}{Project file for \dotnet{}~6 console application project.}{Sdk}
<Project Sdk="Microsoft.NET.Sdk">

  <PropertyGroup>
    <OutputType>Exe</OutputType>
    <TargetFramework>net6.0</TargetFramework>
    <RootNamespace>hello_net6</RootNamespace>
  </PropertyGroup>

</Project>
\end{myListingXml}

Lastly, the NuGet package feed must be configured to use the correct package source for the preview
packages for \dotnet{}~6. This can be done by creating a \filename{NuGet.Config} file in the project
directory with contents as listed in \autoref{lst:06-nuget-config}. The contents can also be copied
from the webpage from which the \dotnet{}~6 SDK was downloaded~\cite{dotnetSdkGithub}.

\begin{myListingXml}[basicstyle=\ttfamily\scriptsize]{lst:06-nuget-config}{NuGet configuration file for \dotnet{}~6 projects.}{key,value}
<configuration>
  <packageSources>
    <add key="dotnet6"
      value="https://pkgs.dev.azure.com/dnceng/public/_packaging/dotnet6/nuget/v3/index.json" />
  </packageSources>
</configuration>
\end{myListingXml}

To check that the development environment is configured correctly, the program in
\autoref{lst:06-hello-net6} can be used.

\begin{myListingCsharp}{lst:06-hello-net6}{\csharp{} program for testing the SDK installation.}{Program,FileVersionInfo,Console,Path,QuicImplementationProviders}{}
using System;
using System.Diagnostics;
using System.IO;
using System.Net.Quic;

namespace hello_net6._0
{
    class Program
    {
        static void |Main|(string[] args)
        {
            string assemblyPath = typeof(object).Assembly.Location;
            string assemblyDir = Path.|GetDirectoryName|(assemblyPath);
            var info = FileVersionInfo.|GetVersionInfo|(assemblyPath);
            Console.|WriteLine|($"Hello from .NET {info.ProductVersion}");
            Console.|WriteLine|($"Runtime location: {assemblyDir}");
            Console.|WriteLine|($"QUIC: {QuicImplementationProviders.Default}");
        }
    }
}
\end{myListingCsharp}

Running the program should produce output similar to the following:

\noindent\begin{minipage}{\textwidth}
\begin{myVerbatim}
> dotnet run
Hello from .NET 6.0.0-alpha.1.20560.10+72b7d236ad634c2280c73499ebfc2b594995ec06
Runtime location: C:\dotnet\shared\Microsoft.NETCore.App\6.0.0-alpha.1.20560.10
QUIC: System.Net.Quic.Implementations.Managed.ManagedQuicImplementationProvider
\end{myVerbatim}
\end{minipage}\medskip

If the output lists a different runtime location, verify that the environment variables have been
set correctly. If the QUIC provider is different, it means that the \SystemNetQuicDll{} with our
managed QUIC implementation was not copied to the correct directory.

\subsection{Deploying \dotnet{} Applications}

The applications built against the preview \dotnet{}~6 SDK will run only if the environment is
configured according to \autoref{sec:06-env-vars}. In order to run the applications outside the
configured environment, it is necessary to build them as
\textit{self-contained}~\cite{SelfContainedPublishDocs}. Self-contained builds of \dotnet{}
applications contain a copy of the \dotnet{} runtime and other necessary binaries. Self-contained
applications can be built using following command:

\begin{myVerbatim}
> dotnet publish --self-contained --runtime <RID>
\end{myVerbatim}

Where \texttt{<RID>} is the runtime identifier for which to publish. Commonly used values are
\texttt{win-x64} and \texttt{linux-x64}. Full list of supported runtime identifiers can be found in
the official \dotnet{} documentation~\cite{dotnetRIDs}.

Unfortunately, the packaged application contains the unmodified \SystemNetQuicDll{} without managed
QUIC support. Therefore, the last step is manually overwriting the \SystemNetQuicDll{} with the one
locally built from the \dotnet{} runtime sources in \autoref{sec:06-build-runtime}. The modified
\libopenssl{} libraries also need to be copied over to the application directory.

\section{Simple Echo Server using QUIC}\label{sec:06tutorial}

This section is a walkthrough on how to use QUIC in \dotnet{}. In this section, we will create a
trivial echo server. When a new connection is established, clients will open a single bidirectional
stream and send arbitrary data over it. The server will then echo the data back to the client using
the same stream. For simplicity, we will use a single \dotnet{} program to represent both client and
server and conduct the connection over the loopback network interface. We will also omit error
checking from this example for brevity. The complete source code for this example can be found in
the \filename{src/supplementary/samples/Echo/} directory in the thesis attachments.

\subsection{Echo Server Implementation}

The use of TLS 1.3 for encryption is mandatory for QUIC\@. This requires providing an X.509
certificate on the server's side. The current QUIC API requires that the certificate and the private
key be saved in separate files in the PEM format. Certificate files that can be used in this example
are provided in the attachments. The public certificate file is at \filename{certs/cert.crt} and the
private key is at \filename{certs/cert.key}. Alternatively, a new certificate can be created using
the \texttt{openssl} command-line utility using the following commands.

\begin{myVerbatim}
> openssl req -x509 -newkey rsa -keyout key.pem -out cert.pem -days 365 -nodes
\end{myVerbatim}

In order to accept incoming QUIC connections, we need to create an instance of \QuicListener{}.
\autoref{lst:06-echo-quic-listener} shows how the \QuicListener{} can be created and provided with:

\begin{itemize}

  \item listening endpoint;

  \item identifier of the application-layer protocol to be used\footnote{QUIC is not intended to be
        used standalone, but as a transport protocol for other application-layer protocols. Since
        servers can support multiple versions of the application protocol, QUIC uses the \gls{alpn}
        extension to TLS to negotiate the application-layer protocol as part of the QUIC connection
        handshake.}, even though we are not implementing any standard protocol, we still need to
        provide one; for our example, we chose to use \texttt{"echo"} as the \gls{alpn} identifier;
        and

  \item paths to the X.509 certificate and private key file.

\end{itemize}

The \QuicListener{} class implements \interface{IDisposable} and, therefore, we can also use the
\keyword{using} statement to make sure the \QuicListener{} is closed when the method returns.

\begin{myListingCsharp}{lst:06-echo-quic-listener}{Creating and starting a new \QuicListener{} for the echo server}{Task,QuicListener,QuicListenerOptions,SslServerAuthenticationOptions,List,SslApplicationProtocol,IPEndPoint,CancellationToken}{}
public static async Task<int> |RunServer|(IPEndPoint listenEp,
    string certificateFile, string privateKeyFile, CancellationToken token)
{
    using QuicListener listener = new QuicListener(new QuicListenerOptions
    {
        ListenEndPoint = listenEp,
        CertificateFilePath = certificateFile,
        PrivateKeyFilePath = privateKeyFile,
        ServerAuthenticationOptions = new SslServerAuthenticationOptions
        {
            ApplicationProtocols = new List<SslApplicationProtocol>
            {
                new SslApplicationProtocol("echo")
            }
        }
    });

    // ...
}
\end{myListingCsharp}

We can then use the \method{AcceptConnectionAsync} method to wait for new incoming connections
asynchronously. \autoref{lst:06-echo-server-accept} shows how to accept new connections and process
them asynchronously in a separate \class{Task} so that multiple connections can be served
concurrently.

\begin{myListingCsharp}{lst:06-echo-server-accept}{Accepting new connections on \QuicListener{}}{List,Task,QuicConnection}{}
    // QuicListener must be started before accepting connections.
    listener.|Start|();

    // tasks that need to be awaited when trying to exit gracefully
    List<Task> tasks = new List<Task>();

    try
    {
        QuicConnection conn;
        while ((conn = await listener.|AcceptConnectionAsync|(token)) != null)
        {
            // copy the connection into a variable with narrower scope which
            // can be safely captured inside the lambda function
            QuicConnection captured = conn;
            var task = Task.|Run|(
                () => |HandleServerConnection|(captured, token));
            tasks.|Add|(task);
        }
    }
    finally
    {
        // wait until all connections are closed
        await Task.|WhenAll|(tasks);
    }
\end{myListingCsharp}

\autoref{lst:06-echo-server-loop} Shows the implementation of \method{HandleServerConnection} which
does the actual echoing of the incoming data. The \QuicStream{} is accepted using the
\method{AcceptStreamAsync} method on the \QuicConnection{} class and, because it is a bidirectional
stream, we can use it to send the data back to the client. Lastly, once all data is sent, we
gracefully close the connection using the \method{CloseAsync} method.

\begin{myListingCsharp}{lst:06-echo-server-loop}{Echo server reading and writing data to \QuicStream{}}{Task,QuicConnection,CancellationToken,QuicStream}{}
public static async Task |HandleServerConnection|(QuicConnection connection,
    CancellationToken token)
{
    try
    {
        QuicStream stream = await connection.|AcceptStreamAsync|(token);

        byte[] buffer = new byte[4 * 1024];

        int read;
        while ((read = await stream.|ReadAsync|(buffer, token)) > 0)
        {
            await stream.|WriteAsync|(buffer, 0, read, token);
            await stream.|FlushAsync|(token);
        }
    }
    finally
    {
        // gracefully close the connection with 0 error code
        await connection.|CloseAsync|(0);
    }
}
\end{myListingCsharp}

\subsection{Echo Client Implementation}

The client implementation is more straightforward than that of the server.
\autoref{lst:06-echo-client-connection} shows how to create a client \QuicConnection{} using the
using the server endpoint address and the \gls{alpn} identifier. The connection can be then
established by calling the \method{ConnectAsync} method.

\begin{myListingCsharpNoPageBreak}{lst:06-echo-client-connection}{Creating a client \QuicConnection{} instance}{Task,IPEndPoint,CancellationToken,QuicConnection,QuicClientConnectionOptions,SslClientAuthenticationOptions,List,SslApplicationProtocol}{}
public static async Task<int> |RunClient|(IPEndPoint serverEp,
    CancellationToken token)
{
    using var client = new QuicConnection(new QuicClientConnectionOptions
    {
        RemoteEndPoint = serverEp,
        ClientAuthenticationOptions = new SslClientAuthenticationOptions
        {
            ApplicationProtocols = new List<SslApplicationProtocol>
            {
                new SslApplicationProtocol("echo")
            }
        }
    });

    await client.|ConnectAsync|(token);

    // ...
}
\end{myListingCsharpNoPageBreak}

Once the connection is established, the client opens a bidirectional \QuicStream{} using the
\method{OpenBidirectionalStream} method. The \QuicStream{} can be used like any other \Stream{}
instance. \autoref{lst:06-echo-client-loop} shows the rest of the echo client implementation.

\begin{myListingCsharp}{lst:06-echo-client-loop}{Sending standard input via \QuicStream{}}{QuicStream,Task,Encoding,Console}{}
    try
    {
        await using QuicStream stream = client.|OpenBidirectionalStream|();

        // spawn a reader task to not let server be flow-control blocked
        _ = Task.|Run|(async () =>
        {
            byte[] arr = new byte[4 * 1024];
            int read;
            while ((read = await stream.|ReadAsync|(arr, token)) > 0)
            {
                string s = Encoding.ASCII.|GetString|(arr, 0, read);
                Console.|WriteLine|($"Received: {s}");
            }
        });

        string line;
        while ((line = Console.|ReadLine|()) != null)
        {
            // convert into ASCII byte array before sending
            byte[] bytes = Encoding.ASCII.|GetBytes|(line);
            await stream.|WriteAsync|(bytes, token);
            // flush the stream to send the data immediately
            await stream.|FlushAsync|();
        }

        // once all stdin is written, close the stream
        stream.|Shutdown|();

        // wait until the server receives all data
        await stream.|ShutdownWriteCompleted|(token);
    }
    finally
    {
        // gracefully close the connection with 0 error code
        await client.|CloseAsync|(0, token);
    }
\end{myListingCsharp}

\subsection{A More Complex Example Application}

In this thesis, we developed a more complex version of the echo server example from the previous
section for benchmarking purposes. The source code for this benchmarking application can be found in
the \filename{src/supplementary/benchmark/ThroughputTests/} directory of the thesis attachments. We
will provide more information about the application in \autoref{sec:04-benchmark-app}.

\section{QUIC API Reference}\label{sec:06-api}

This section describes the API designed by the \dotnet{} development team to expose QUIC to other
developers. As mentioned in the introduction chapter, the current design is a work-in-progress and
is subject to change in the future. All of the mentioned classes are located in the
\namespace{System.Net.Quic} namespace.

\subsection{QuicListener Class}

The \class{QuicListener} class is the equivalent of the \class{TcpListener} for TCP connections.
Servers use this class to accept incoming QUIC connections.

\begin{description}

    \ditemctor{QuicListener}{\QuicListenerOptions{}} Constructor.

    \ditemproperty{IPEndPoint}{ListenEndPoint}{\propget} The IP endpoint being listened to for new
connection. Read-only.

    \ditemmethod[]{\ValueTaskOf{\QuicConnection{}}}{AcceptConnectionAsync}{\CancellationToken{}}
Accepts a new incoming QUIC Connection.

    \ditemmethod[\keyword]{void}{Start}{} Starts listening.

    \ditemmethod[\keyword]{void}{Close}{} Stops listening and closes the listener. Does not close
already accepted connections.

\end{description}

\subsection{QuicListenerOptions Class}

The \class{QuicListenerOptions} class holds all configuration used to construct new
\class{QuicListener} instances.

\begin{description}

    \ditemproperty{SslServerAuthenticationOptions}{\ServerAuthenticationOptions{}}{\propgetset} SSL
related options like certificate selection/validation callbacks, and supported protocols for ALPN\@.

    \ditemproperty[\keyword]{string}{CertificateFilePath}{\propgetset} Path to the X.509 certificate
used by the server.

    \ditemproperty[\keyword]{string}{CertificateKeyPath}{\propgetset} Path to the private key for
the used X.509 certificate.

    \ditemproperty{IPEndPoint}{ListenEndPoint}{\propgetset} The IP endpoint to listen on.

    \ditemproperty[\keyword]{int}{ListenBacklog}{\propgetset} Number of connection to be held
waiting for acceptance by the application. Upon reaching this limit, further connections will be
refused.

    \ditemproperty[\keyword]{long}{MaxBidirectionalStreams}{\propgetset} Limit on the number of
bidirectional streams the client can open in an accepted connection.

    \ditemproperty[\keyword]{long}{MaxUnidirectionalStreams}{\propgetset} Limit on the number of
unidirectional streams the client can open in an accepted connection.

    \ditemproperty{TimeSpan}{IdleTimeout}{\propgetset} The period of inactivity after which the
connection will be closed via idle timeout.

\end{description}

\subsection{QuicConnection Class}

The \QuicConnection{} class represents the QUIC connection itself. Clients open new connections by
creating a new instance of this class and calling the \method{Connect\allowbreak{}Async} method. Servers receive
new connections using the \class{QuicListener} class.

\begin{description}

    \ditemctor{QuicConnection}{\QuicClientConnectionOptions{}} Constructor. The newly created
instance must be explicitly connected using the \method{ConnectAsync} method.

    \ditemproperty[\keyword]{bool}{Connected}{\propget} Indicates whether the \QuicConnection{} is
connected (the handshake has completed).

    \ditemproperty{IPEndPoint}{LocalEndPoint}{\propget} Local IP endpoint of the connection.

    \ditemproperty{IPEndPoint}{RemoteEndPoint}{\propget} Remote IP endpoint of the connection.

    \ditemmethod{ValueTask}{ConnectAsync}{\CancellationToken{}} Connects to the remote endpoint.

    \ditemmethod{QuicStream}{OpenUnidirectionalStream}{} Opens a new unidirectional stream. Throws a
\class{QuicException} if the stream cannot be opened.

    \ditemmethod{QuicStream}{OpenBidirectionalStream}{} Opens a new bidirectional stream. Throws a
\class{QuicException} if the stream cannot be opened.

    \ditemmethod[]{\ValueTaskOf{\QuicStream{}}}{AcceptStreamAsync}{\CancellationToken{}} Accepts an
incoming stream.

    \ditemmethod{ValueTask}{CloseAsync}{\Long{}, \CancellationToken{}} Closes the connection with
the specified given error code and terminates all active streams.

    \ditemmethod[\keyword]{long}{GetRemoteAvailableUnidirectionalStreamCount}{} Gets the maximum
number of unidirectional streams that this endpoint can open.

    \ditemmethod[\keyword]{long}{GetRemoteAvailableBidirectionalStreamCount}{} Gets the maximum
number of bidirectional streams that this endpoint can open.

\end{description}

\subsection{QuicClientConnectionOptions}

The \class{QuicClientConnectionOptions} is used by clients to configure new QUIC conections.

\begin{description}

    \ditemproperty{SslClientAuthenticationOptions}{ClientAuthenticationOptions}{\propgetset} Client
authentication options to use when establishing the connection.

    \ditemproperty{IPEndPoint}{LocalEndPoint}{\propgetset} The local IP endpoint that will be bound
to.

    \ditemproperty{IPEndPoint}{RemoteEndPoint}{\propgetset} The IP endpoint to connect to.

    \ditemproperty[\keyword]{long}{MaxBidirectionalStreams}{\propgetset} Limit on the number of
bidirectional streams the server can open.

    \ditemproperty[\keyword]{long}{MaxUnidirectionalStreams}{\propgetset} Limit on the number of
unidirectional streams the server can open.

    \ditemproperty{TimeSpan}{IdleTimeout}{\propgetset} The period of inactivity after which the
connection will be closed via idle timeout.

\end{description}

\subsection{QuicStream Class}

The \QuicStream{} class represents a single stream in a QUIC connection and derives from the
abstract \class{Stream} class. The \class{Stream} class is a bidirectional stream abstraction and
since not all QUIC streams are bidirectional, user should check if the specific \QuicStream{}
instance supports supports the operation by inspecting the \method{CanRead} and \method{CanWrite}
properties. Invoking write methods on read-only --- more specifically, incoming unidirectional ---
stream will cause an \class{InvalidOperationException} to be thrown and vice versa.

The list below mentions the members specific for the \class{QuicStream} class and some important
members inherited from the \class{Stream} class.

\begin{description}

    \ditemproperty[\keyword]{long}{StreamId}{\propget} The Stream ID\@.

    \ditemproperty[\keyword]{bool}{CanRead}{\propget} Returns \keyword{true} if the stream supports
reading.

    \ditemproperty[\keyword]{bool}{CanWrite}{\propget} Returns \keyword{true} if the stream supports
reading.

    \ditemmethod[\keyword]{void}{AbortRead}{\Long} Aborts the receiving part of the stream with the
provided error code.

    \ditemmethod[\keyword]{void}{AbortWrite}{\Long} Aborts the sending part of the stream with the
provided error code.

    \ditemmethod[\keyword]{int}{Read}{\SpanOf{\Byte{}}} Reads the content of the stream into the
provided buffer, blocks if no data is available. Returns 0 only when there will be no more data in
the stream.

    \ditemmethod[]{\ValueTaskOf{\keyword{int}}}{ReadAsync}{\MemoryOf{\Byte{}}, \CancellationToken{}}
Reads the content of the stream into provided buffer, blocks until some data is available. Returns 0
only when there will be no more data in the stream.

    \ditemmethod[\keyword]{void}{Write}{\SpanOf{\Byte{}}} Writes the content of the provided buffer
into the stream, returns when the data have been buffered internally.

    \ditemmethodWithComment{ValueTask}{WriteAsync}{*, \CancellationToken{}}{multiple overloads}
Multiple overloads of this method offer writing from various types of buffers:
\ReadOnlyMemoryOf{\Byte{}}, \ReadOnlySequenceOf{\Byte{}}, and
\genericClass{ReadOnlyMemory\allowbreak{}}{\ReadOnlyMemoryOf{\Byte{}}}. The last one can be used to
perform \textit{Vectored I/O}~\cite{wiki:vectored-io}\@. The returned task completes when the
provided data have been buffered internally and the buffers can be reused for other purposes.

    \ditemmethodWithComment{ValueTask}{WriteAsync}{*, \bool{}, \CancellationToken{}}{multiple
overloads} Like the methods above, but also allow specifying that the provided data are the last on
the stream and that the stream should be gracefully closed.

    \ditemmethod{ValueTask}{ShutdownWriteCompleted}{\CancellationToken{}} The returned task
completes when the stream shutdown completes. Meaning that acknowledgment from the peer is received.

    \ditemmethod{ValueTask}{Shutdown}{} Gracefully closes the writing direction of the stream,
indicating that no more data will be sent.

\end{description}

\subsection{Exceptions}

The QUIC API can throw the following exceptions:

\begin{description}

    \ditem{\ditemsrcsize\class{QuicException}} Base class for all thrown exceptions, used when a
more specific exception is not available

    \ditem{\ditemsrcsize\class{QuicConnectionAbortedException}} Thrown when the connection is
forcibly closed either by the transport or by the remote endpoint.

    \ditem{\ditemsrcsize\class{QuicStreamAbortedException}} Thrown when the stream was aborted by
the remote endpoint.

    \ditem{\ditemsrcsize\class{QuicOperationAbortedException}} Thrown when the pending operation was
aborted by the local endpoint.

\end{description}


\chapter*{Conclusion}
\addcontentsline{toc}{chapter}{Conclusion}

\begin{itemize}

    \item Evaluation
    \begin{itemize}

        \item Types of benchmarks
        \item Performance charts
        \item Robustness tests and simulated network issues
        \item Conclusion

    \end{itemize}

    \item Integration to .NET and out-of-runtime build??

    \item Documentation
        \begin{itemize}

            \item Obtaining the source code

        \end{itemize}

        \item Building the library
        \begin{itemize}

            \item Standalone
            \item Part of dotnet runtime

        \end{itemize}

        \item API documentation
        \begin{itemize}

            \item Documented dispose and abort model

        \end{itemize}

        \item Samples
        \begin{itemize}

            \item Actual executable code samples

        \end{itemize}

    \item Conclusion
    \begin{itemize}

        \item Goal evaluation
        \item Future work

    \end{itemize}

\end{itemize}


%%% Bibliography
%%% Bibliography (literature used as a source)
%%%
%%% We employ bibTeX to construct the bibliography. It processes
%%% citations in the text (e.g., the \cite{...} macro) and looks up
%%% relevant entries in the bibliography.bib file.
%%%
%%% The \bibliographystyle command selects, which style will be used
%%% for references from the text. The argument in curly brackets is
%%% the name of the corresponding style file (*.bst). Both styles
%%% mentioned in this template are included in LaTeX distributions.

\bibliographystyle{unsrtnat}     %% [number]

\renewcommand{\bibname}{Bibliography}

%%% Generate the bibliography. Beware that if you cited no works,
%%% the empty list will be omitted completely.

\bibliography{bibliography}

%%% If case you prefer to write the bibliography manually (without bibTeX),
%%% you can use the following. Please follow the ISO 690 standard and
%%% citation conventions of your field of research.

% \begin{thebibliography}{99}
%
% \bibitem{lamport94}
%   {\sc Lamport,} Leslie.
%   \emph{\LaTeX: A Document Preparation System}.
%   2nd edition.
%   Massachusetts: Addison Wesley, 1994.
%   ISBN 0-201-52983-1.
%
% \end{thebibliography}


%%% Figures used in the thesis (consider if this is needed)
\listoffigures

%%% Tables used in the thesis (consider if this is needed)
%%% In mathematical theses, it could be better to move the list of tables to the beginning of the thesis.
\listoftables

%%% Code listings used in the thesis (consider if this is needed)
\lstlistoflistings

%%% Attachments to the master thesis, if any. Each attachment must be
%%% referred to at least once from the text of the thesis. Attachments
%%% are numbered.
%%%
%%% The printed version should preferably contain attachments, which can be
%%% read (additional tables and charts, supplementary text, examples of
%%% program output, etc.). The electronic version is more suited for attachments
%%% which will likely be used in an electronic form rather than read (program
%%% source code, data files, interactive charts, etc.). Electronic attachments
%%% should be uploaded to SIS and optionally also included in the thesis on a~CD/DVD.
%%% Allowed file formats are specified in provision of the rector no. 72/2017.
\appendix
\chapter{Attachments}

\todo{builds of dotnet runtime for linux and windows, if they fit the 850 limit}

Contents of the attachment archive:

\dirtree{%
  .1 {/}.
  .2 {certs}\mydtcomment{Certificates for use in tests and attached samples}.
  .3 {cert.crt}\mydtcomment{Public certificate file}.
  .3 {cert.key}\mydtcomment{Private key file}.
  .3 {cert.pfx}\mydtcomment{Certificate and private key combined in PFX format}.
  .2 {extern}\mydtcomment{Third-party source code}.
  .3 {openssl}\mydtcomment{Copy of the QUIC-enabled fork of \libopenssl{} by Akamai}.
  .2 {src}\mydtcomment{Root for all source code}.
  .3 {dotnet-runtime}\mydtcomment{\dotnet{} runtime fork with managed QUIC implementation}.
  .3 {supplementary}\mydtcomment{Root of additional \dotnet{} projects}.
  .4 {benchmark}\mydtcomment{Root of all benchmark projects}.
  .5 {ThroughputTests}\mydtcomment{Benchmarking application used in \autoref{chap:04-evaluation}}.
  .6 {run.ps1}\mydtcomment{Script used to automate running the benchmarking application}.
  .4 {samples}\mydtcomment{Root of all sample projects}.
  .5 {Echo}\mydtcomment{Simple echo server and client implementation from \autoref{chap:06-user-docs}}.
  .5 {hello-net6.0}\mydtcomment{Hello world application for testing \dotnet{}~6 installation}.
  .4 {ManagedQuic.sln}\mydtcomment{\dotnet{} Solution file for all projects in subdirectory}.
  .2 {tex}\mydtcomment{\LaTeX{} source code for this thesis}.
  .3 {measurements}\mydtcomment{Raw CSV data with measurements presented in \autoref{chap:04-evaluation}}.
  .2 {README.md}\mydtcomment{File describing the contents of the archive}.
}


%%% Abbreviations used in the thesis, if any, including their explanation
%%% In mathematical theses, it could be better to move the list of abbreviations to the beginning of the thesis.
\printglossary

\openright
\end{document}
